\documentclass[11pt]{amsart}
\usepackage[T1]{fontenc}
\usepackage[utf8]{inputenc}
\usepackage{lmodern,microtype,amsmath,amssymb,amsthm,mathtools}
\usepackage[margin=26mm]{geometry}
\usepackage{xcolor}
\usepackage[colorlinks=true,linkcolor=blue!45!black,citecolor=blue!45!black,urlcolor=blue!45!black]{hyperref}
\usepackage{xurl}
\setlength{\emergencystretch}{3em}
\newcommand{\workbenchurl}{https://github.com/KokunoYumeto/schemes-and-moduli-workbench}
\newcommand{\recordurl}{https://github.com/KokunoYumeto/mathematics-workbenches}
\newcommand{\B}{\mathbb B}
\newcommand{\Z}{\mathbb Z}
\newcommand{\N}{\mathbb N}
\newcommand{\Ab}{\mathbf{Ab}}
\newcommand{\Mod}{\mathrm{Mod}}
\newcommand{\Qcoh}{\operatorname{QCoh}}
\DeclareMathOperator{\Spec}{Spec}
\DeclareMathOperator{\Id}{Id}
\DeclareMathOperator{\im}{im}
\DeclareMathOperator{\Ann}{Ann}
\DeclareMathOperator{\supp}{supp}
\DeclareMathOperator{\End}{End}
\DeclareMathOperator{\Hom}{Hom}
\newtheorem{theorem}{Theorem}[section]
\newtheorem{proposition}[theorem]{Proposition}
\newtheorem{lemma}[theorem]{Lemma}
\newtheorem{corollary}[theorem]{Corollary}
\theoremstyle{definition}
\newtheorem{definition}[theorem]{Definition}
\newtheorem{example}[theorem]{Example}
\theoremstyle{remark}
\newtheorem{remark}[theorem]{Remark}
\title[Projectives over inverse semirings]{Projective semimodules over inverse semirings, after Faul, Goswami, Joubert and Manuell}
\author{}
\date{25 November 2026}
\begin{document}
\begin{abstract}
Faul, Goswami, Joubert and Manuell introduced zero adjunction and described inverse semirings with two additive idempotents. We classify finitely generated projectives over every such semiring by a projective ring module and a finite poset decomposition of its reduction. For arbitrary inverse semirings, we give the semimodule counterpart of Caviglia, Faul, Manuell and Mesiti's structure theorem and identify an action category by explicit directed coefficient colimits. Projectivity is equivalent to projectivity of the idempotent support and ordinary projectivity of its amplitude, generated by a canonical ideal. A finite support frame computes that ideal by scalar constraints. A single category of coefficient groups models all projective amplitudes on a fixed support, including every morphism, by projectives generated by an explicit ideal.

Compatible coefficient quotients preserve projective isomorphism classes when their kernels in the upper multiplicative-idempotent rings are locally nilpotent. A larger Jacobson quotient always detects isomorphism and lifts isomorphisms between existing projectives. When those upper rings are exchange rings, every compatible coefficient quotient lifts all finite projectives at their exact supports. Under semiperfectness, with arbitrary support semiring, we classify every projective amplitude by finite multiplicities of local idempotents in the canonical action ideal. These results retain support and differ from global semiring radical quotients.

We also obtain labelled-poset classifications for zero-adjoined rings, their iterates, finite distributive support extensions and finite group semirings, together with ordinary and strong Grothendieck groups and an affine equivalence for ordinary quasi-coherent semimodule sheaves. Counterexamples specify the limits of completion, radical lifting and uniform decision procedures. The unrestricted support classification remains open. All principal statements were re-derived by separate model workers before proof comparison; no human has checked the mathematics.
\end{abstract}
\maketitle

\section{The construction, statements and context}
A semiring has an additive commutative monoid with identity, an associative unital multiplication, distributivity, and an absorbing additive zero. An \emph{inverse semiring} has an inverse additive monoid. Faul, Goswami, Joubert and Manuell's Definition~3.6 \cite{FGJM} associates to an inverse semiring $R$ the semiring $R_0$ obtained by adjoining a new additive identity and multiplicative absorber. Unless a classification explicitly permits noncommutative rings, $R$ is a commutative unital ring, including the zero ring. Write
\[
 S=R_0=R\sqcup\{\tau\},\qquad e=0_R.
\]
Thus $e$ is an element of the old ring, $e+r=r$ and $er=e$ for $r\in R$, while $\tau+s=s$ and $\tau s=\tau$ for every $s\in S$. The multiplicative unit remains $1_R$, even when $R$ is the zero ring. The Boolean semifield is $\B=\{0,1\}$ with addition maximum and multiplication minimum.

Semimodules and their homomorphisms preserve the additive zero and have a unital scalar action. Projective means projective with respect to surjective homomorphisms. A finitely generated projective is therefore a retract of a finite free semimodule. This convention is different from \emph{strong projectivity}, which requires an isomorphism $P\oplus P'\simeq S^n$ for some $P'$; see Katsov, Nam and Zumbr\"agel \cite[Definition~4.1]{KNZ}. Finite direct sums of semimodules are their Cartesian products with componentwise operations.

For a join-semilattice $L$ with least element $0_L$, a \emph{diagram of $R$-modules over $L$} means a functor $A:L\to R\text{-}\Mod$, where the order is regarded as a category. Write $\rho_{l,k}:A_l\to A_k$ for its maps. A morphism of such diagrams consists of a zero-and-join preserving map $u:L\to L'$ and a natural family of linear maps $A_l\to A'_{u(l)}$. In particular, the module $A_{0_L}$ is allowed to be nonzero.

For a finite poset $P$, let $\Id(P)$ be its lattice of order ideals. Given finitely generated projective $R$-modules $P_p$ for $p\in P$, form the diagram
\begin{equation}\label{eq:labelled}
 A_I=\bigoplus_{p\in I}P_p\quad (I\in\Id(P)),
\end{equation}
whose maps are the inclusions of summands. Zero labels are allowed. Let $\mathcal P(P,(P_p))$ denote the corresponding semimodule in Theorem~\ref{thm:main}.

Section~\ref{sec:general} gives the module counterpart of Caviglia, Faul, Manuell and Mesiti's Theorem~8.5. If an inverse semiring corresponds to $(Q,F)$, then its modules are diagrams of abelian groups over $Q$-semimodules with compatible bilinear action maps. This is a description for all inverse semirings, not only those arising from rings. We give exact finite matrix descriptions of their projectives. Theorem~\ref{thm:canonical-ideal} separates ordinary projectivity over an explicit additive action category from generation by a canonical support ideal. It removes the presentation-dependent support-section search; classifying arbitrary supports and ordinary action-category projectives remains open. For all the two-idempotent inverse semirings of Construction~3.5 of Faul, Goswami, Joubert and Manuell, the following theorem gives an intrinsic classification. Theorem~\ref{thm:finite-lattice-projectives} also resolves finite distributive support extensions with a single top coefficient ring. The corresponding intrinsic problem for an unrestricted pair $(Q,F)$ remains open.

Let now $R$ be any unital ring, not necessarily commutative, let $A$ be an $R$--$R$-bimodule, and let $f:A\to R$ be a bimodule map satisfying $f(a)b=af(b)$. Construction~3.5 is $T=I(f)=R\sqcup A$: old sums and products are those of $R$, lower sums are those of $A$, $r+a=r+f(a)$ in $R$, lower multiplication is $ab=f(a)b$, and mixed products are the bimodule actions. Its zero is $0_A$ and its unit is $1_R$. Put $I=f(A)$ and $B=R/I$. The module assertions of the labelled-poset theorem also hold for noncommutative rings, with left projective labels; Lemma~\ref{lem:noncomm} records this extension.

\begin{theorem}[projectives over two-idempotent inverse semirings]\label{thm:twolevelprojectives}
Isomorphism classes of finitely generated projective left $I(f)$-semimodules are classified by triples
\[
 (P,N,\phi),\qquad \phi:P/IP\xrightarrow{\sim}N_{\mathrm{top}},
\]
where $P$ is a finitely generated projective left $R$-module and $N$ is a finitely generated projective left $B_0$-semimodule. Thus $N$ is a finite poset $J$ labelled by finite projective left $B$-modules $Q_j$, with $N_H=\bigoplus_{j\in H}Q_j$ for $H\in\Id(J)$ and $N_{\mathrm{top}}=\bigoplus_{j\in J}Q_j$.

Two triples are equivalent precisely when there are an $R$-module isomorphism $u:P\to P'$ and a $B_0$-semimodule isomorphism $v:N\to N'$ such that
\[
 v_{\mathrm{top}}\phi=\phi'\overline u.
\]
Here $v$ is a full diagram isomorphism, so it may mix summands while preserving every order-ideal submodule. Direct sums of triples correspond to direct sums of projectives.

If $f$ is injective, identify $A$ with $I$. The corresponding semimodule has group
\begin{equation}\label{eq:twolevelfiber}
 P_H=\{p\in P:\phi(p\bmod IP)\in N_H\}
\end{equation}
over $H$, with inclusion maps and least group $IP$. Old ring scalars act in each fiber; lower scalars in $I$ act by multiplication into $IP$.

For arbitrary $f$, reduction $I(f)\to I(I\hookrightarrow R)$ induces a bijection on isomorphism classes of finite projectives, compatible with direct sum. Every class is realized by lifting an idempotent matrix; the gluing triples classify these lifts. This last assertion is about isomorphism classes, and does not assert an equivalence of the full projective categories.
\end{theorem}

For a commutative $R$, the data have a geometric reading: a vector bundle on $\Spec R$ carries a finite distributive family of subbundles after restriction to the closed subscheme $\Spec(R/I)$. The incidence poset and the projective successive quotients specify that family. The kernel of $f$ can alter morphisms but creates no additional isomorphism classes of finite projectives. For $A=0$, the theorem reduces to the zero-adjoined labelled-poset classification; for $f=1_R$, it gives a ring-projective module and an independent finite distributive lattice.

\begin{theorem}[spectrum, diagrams and projectives]\label{thm:main}
For every commutative unital ring $R$ the following hold.
\begin{enumerate}
\item The ideals of $S$ are $\{\tau\}$ and $J\sqcup\{\tau\}$ for the ideals $J$ of $R$. The prime spectrum is $\Spec R$ together with one new point $\eta=\{\tau\}$, contained in every nonempty open. Its open sets are the empty set and $\{\eta\}\cup U$ for $U\subseteq\Spec R$ open. The stalk at $\eta$ is $\B$, and the stalk at $\mathfrak p\cup\{\tau\}$ is $(R_{\mathfrak p})_0$.
\item The congruences of $S$ are indexed by an ideal $J$ of $R$ and the choice of whether $\tau$ is a singleton class. Their quotients are respectively $(R/J)_0$ and $R/J$.
\item The category of $S$-semimodules is equivalent to the category of diagrams of $R$-modules over join-semilattices with zero.
\item Every finitely generated projective $S$-semimodule is $\mathcal P(P,(P_p))$ for a finite poset labelled by finitely generated projective $R$-modules. Two such semimodules are isomorphic exactly when a poset isomorphism identifies the isomorphism classes of all their labels.
\end{enumerate}
\end{theorem}

The theorem describes what the two additive zeros remember. The usual ring spectrum is still present, and the new point sees the semilattice of additive idempotents of a semimodule. Projectivity couples a finite distributive lattice to projective ring modules. A nonzero ordinary $R$-module, regarded as an $S$-semimodule through $S\to R$, has only one additive idempotent but is \emph{not} projective over $S$. Indeed, the least group of a free $S$-semimodule is zero, and this property passes to retracts. The projectives with one idempotent consist only of the zero object.

Let $V(S)$ be the commutative monoid of isomorphism classes of finitely generated projective $S$-semimodules, and let $K_0(S)$ be its group completion. Let $SV(S)$ and $SK_0(S)$ use strongly projective semimodules instead. Let $\mathcal C_R$ be the set of isomorphism classes of nonempty connected finite posets labelled by isomorphism classes of finitely generated projective $R$-modules. Connected means that the undirected graph joining comparable distinct elements is connected.

\begin{theorem}[Grothendieck monoids and groups]\label{thm:Kzero}
There are canonical isomorphisms
\[
 V(R_0)\simeq\N^{(\mathcal C_R)},\qquad
 K_0(R_0)\simeq\Z^{(\mathcal C_R)},\qquad
 SV(R_0)\simeq\N,\qquad SK_0(R_0)\simeq\Z.
\]
Every finitely generated strongly projective $R_0$-semimodule is free. Under extension of scalars along $R_0\to R$, a labelled poset maps to $\bigoplus_{p\in P}P_p$.
\end{theorem}

Here the labels belong to $V(R)$, not merely to $K_0(R)$. A one-point poset labelled $A\oplus B$ and two isolated points labelled $A,B$ have different idempotent lattices. Passing to the group completion of $V(R_0)$ does not identify these objects.

The projective classification extends to every iterate of the construction. Write $R^{[k]}$ for the ring $R$ with $k\ge1$ zeros adjoined successively, and identify its additive idempotents with $C_{k+1}=\{c_0<\cdots<c_k\}$; addition is maximum, multiplication is minimum, and $c_k=0_R$. The scalar unit is $1_R$, whose additive idempotent is $c_k$.

\begin{theorem}[projectives after repeated zero adjunction]\label{thm:iteratedprojective}
For each $k\ge1$, finitely generated projective $R^{[k]}$-semimodules are classified by the following data, up to isomorphism:
\begin{enumerate}
\item a finite poset $P$ and an order-preserving colour $h:P\to\{1,\ldots,k\}$;
\item for every $p\in P$ and $1\le i\le h(p)$, a greatest element $p_i$ of $\{q\le p:h(q)\le i\}$, with $h(p_i)=i$;
\item a finitely generated projective $R$-module $P_p$ at each point with $h(p)=k$, allowing zero modules.
\end{enumerate}
The associated semimodule has idempotent lattice $\Id(P)$, with the scalar $c_i$ acting on an ideal by intersection with $U_i=\{p:h(p)\le i\}$, and has group
\[
 A_I=\bigoplus_{\substack{p\in I\\h(p)=k}}P_p
\]
over $I$, with inclusion maps. A scalar $c_i$ for $i<k$ sends any vector to the group zero over $I\cap U_i$. An old ring scalar acts in the group over $I$. Isomorphisms are precisely colour-preserving poset isomorphisms identifying the module labels.

Let $\mathcal C_{R,k}$ be the isomorphism classes of nonempty connected data of this kind. Then
\[
 V(R^{[k]})\simeq\N^{(\mathcal C_{R,k})},\qquad
 K_0(R^{[k]})\simeq\Z^{(\mathcal C_{R,k})},\qquad
 SK_0(R^{[k]})\simeq\Z.
\]
Every finitely generated strongly projective $R^{[k]}$-semimodule is free.
\end{theorem}

The condition says that truncating any join-irreducible support to a smaller positive level again gives a join-irreducible support, and that every intervening level occurs. It is an intrinsic condition on a finite coloured poset, so the statement removes the need to search over idempotent matrices. When $k=1$ the condition is automatic and gives Theorems~\ref{thm:main}(4) and~\ref{thm:Kzero}. We prove that case first, because its group-splitting argument carries over unchanged; Section~\ref{sec:iteration} then proves the additional support condition.

We use the prime-spectrum semiring schemes described by Lorscheid \cite[\S1.1 and \S1.4.1]{Lorscheid2026}. For $S=R_0$ this is also the spectrum of its associated semiring blueprint: all its ideals are subtractive. Our sheaves are ordinary semimodule sheaves. We call one quasi-coherent if on a cover by distinguished affine opens it is the sheaf associated, by localization, to a semimodule. This fixes the convention and makes no assertion about different categories of blue modules.

\begin{theorem}[quasi-coherent sheaves]\label{thm:QC}
Put $Y=\Spec R$ and $X=\Spec R_0$. In the convention just specified, $\Qcoh(X)$ is equivalent to diagrams $L\to\Qcoh(Y)$ over join-semilattices with zero. For such a diagram $(\mathcal A_l)$, its value on a nonempty open $\{\eta\}\cup U$ is
\begin{equation}\label{eq:sheaf}
 \coprod_{l\in L}\Gamma(U,\mathcal A_l).
\end{equation}
The group over each label is added using the diagram maps. When $U=\varnothing$, each section group is zero and the value is $L$ as a $\B$-semimodule. Consequently association and global sections give inverse equivalences between $R_0$-semimodules and $\Qcoh(X)$.
\end{theorem}

The chain result has an extension in which ring modules can first appear
at nonprincipal order ideals. Let $C$ be a finite poset,
$L=\Id(C)$, and let $R$ be any unital ring. Define
\[
 T=G_L(R)=\bigl(\{0_R\}\times L\bigr)\cup\bigl(R\times\{C\}\bigr)
 \subseteq R\times L,
\]
where union and intersection are the operations on $L$.
Its unit is $(1_R,C)$ and its additive zero is $(0_R,\varnothing)$.
This is the finite distributive case of the support extension studied
in Section~\ref{sec:latticeextension} below.

\begin{theorem}[projectives over finite distributive support extensions]
\label{thm:finite-lattice-projectives}
The finitely generated projective left $G_{\Id(C)}(R)$-semimodules are
classified by the following data.
\begin{enumerate}
\item A finite poset $P$ and an order-preserving map $c:P\to C$ such that,
for every $p\in P$ and $a\le c(p)$, the set
\[
 \{q\le p:c(q)\le a\}
\]
has a greatest element $p_a$, with $c(p_a)=a$.
\item A finitely generated projective left $R$-module $Q_A$ at every
\emph{admissible} order ideal $A\subseteq P$. An ideal is admissible if,
for every $a\in C$, there is an element $p_a\in A$ of color $a$ such that
\begin{equation}\label{eq:admissible-ideal}
 \{p\in A:c(p)\le a\}=\mathord\downarrow p_a.
\end{equation}
Zero labels are permitted.
\end{enumerate}
The support lattice and group fibers are
\[
 N=\Id(P),\qquad lI=\{p\in I:c(p)\in l\},\qquad
 G_I=\bigoplus_{\substack{A\subseteq I\\A\ \mathrm{admissible}}}Q_A.
\]
The transition maps are inclusions. Old ring scalars act linearly in
each group; a scalar $(0_R,l)$ with $l\ne C$ gives the zero element of
the group over $lI$.

Two data sets give isomorphic semimodules exactly when a colored-poset
isomorphism carries each label to an isomorphic $R$-module. If $C$ is
empty, $P$ is empty and its unique ideal is admissible; the theorem is
ordinary projective $R$-module theory. If $C$ has a greatest element,
the admissible ideals are exactly the principal ideals at points of
that color.
\end{theorem}

For example, take $C$ to be a two-point antichain and $P=C$ with its
identity coloring. The free semimodule $G_L(R)$ has its $R$ label on
$A=P$, while every proper support fiber has zero group. That ideal is
not principal. Labels only on join-irreducibles would therefore omit
the free rank-one object. The proof in
Section~\ref{sec:finite-lattice-projectives} gives a finite construction
of a free retraction from the displayed intrinsic data.


There is also an intrinsic classification with several nonzero coefficient
fibers and a support multiplication different from meet. For a finite group
$\Gamma$, the ordinary monoid semiring $R_0[\Gamma]$ has support semiring
$\mathcal P(\Gamma)$ with union and setwise multiplication, and group $R^A$
over $A\subseteq\Gamma$. Group symbols commute with coefficients, including
when $R$ is noncommutative.

\begin{theorem}[finite group monoid semirings]
\label{thm:finite-group-projectives}
Let $R$ be any unital ring and $\Gamma$ a finite group. Finitely generated
projective left $R_0[\Gamma]$-semimodules are classified by finite
$\Gamma$-posets $P$ and, for one representative $p$ of each orbit, a
finitely generated projective left $R[\Gamma_p]$-module $V_p$, where
$\Gamma_p$ is the stabilizer of $p$. Transport these modules along their
orbits. The support and its groups are
\[
 L=\Id(P),\qquad M_I=\bigoplus_{p\in I}V_p,
\]
with inclusion transitions and the transported $\Gamma$-action. Zero
labels are allowed. Isomorphism means a $\Gamma$-equivariant poset
isomorphism with compatible label isomorphisms; changing an orbit
representative conjugates its stabilizer and transports its module.
No invertibility assumption on $|\Gamma|$ is required.
\end{theorem}

For $R=0$ this says that every finite $\Gamma$-poset gives a projective
$\B[\Gamma]$-semimodule. At a fixed point, the label must be projective
over $R[\Gamma]$. Ordinary projectivity over $R$ alone is insufficient.
Theorem~\ref{thm:finite-group-projectives} evaluates the general support
obstruction for this class; it does not classify arbitrary pairs $(Q,F)$.

The general projectivity problem admits a presentation-independent support
condition. For a left $Q$-semimodule $L$, let $\mathcal A_L$ be the
preadditive action category of Section~\ref{sec:ringoid}: its objects are
$l\in L$, with transition arrows $t_{l,k}$ and scalar arrows
$a_{q,x,l}:l\to ql$, subject to the module relations displayed there.
Call a $Q$-linear map $\phi:L\to Q$ \emph{admissible at $l$} if
$\phi(k)l\leq k$ for every $k\in L$. Define
\[
 \mathcal J_L(l,k)=
 \left\langle
 t_{\phi(k)l,k}a_{\phi(k),x,l}:
 \phi\text{ admissible at }l,\ x\in F(\phi(k))
 \right\rangle_{\Z}.
\]
\begin{theorem}[canonical projectivity ideal]\label{thm:canonical-ideal}
For every unital inverse semiring $T=(Q,F)$, these groups form a
two-sided ideal of $\mathcal A_L$. A $T$-semimodule with support $L$ and
amplitude $G$ is finitely generated projective if and only if
\begin{gather*}
 L\text{ is finitely generated projective over }Q,\qquad
 G\text{ is finitely generated projective over }\mathcal A_L,\\
 \mathcal J_LG=G.
\end{gather*}
Equivalently, the last condition is
$(\mathcal A_L/\mathcal J_L)\otimes_{\mathcal A_L}G=0$.
For a fixed identified support $L$, its projectives are thus exactly the
finite projective $\mathcal A_L$-modules generated by this ideal.
Support isomorphisms transport the category, ideal and amplitude data.
No commutativity, finiteness or injectivity assumption on $(Q,F)$ is required.
\end{theorem}
The ideal depends only on the scalar diagram and the support. It replaces
the existential support-section test of Theorem~\ref{thm:ringoid-projective}
by ordinary module projectivity and generation by one specified ideal.
Classifying arbitrary projective $Q$-supports and projectives over the
resulting ordinary action categories remains a separate problem.

Let $Q$ be a nontrivial unital additively idempotent semiring such that
$qq'=0$ implies $q=0$ or $q'=0$. Such a semiring is called \emph{entire}.
The map
\[
 \chi:Q\longrightarrow\B,\qquad \chi(q)=0\ \Longleftrightarrow\ q=0,
\]
is a unital semiring homomorphism: a join is zero exactly when all its
summands are zero, and the product assertion is the entirety hypothesis.
For a unital ring $R$, possibly noncommutative or zero, put
\[
 T=Q\times_{\B}R_0.
\]
Thus $T$ has one element over $0\in Q$ and the additive group $R$ over
each nonzero $q$. All transitions between nonzero fibers are identities.
For a left $Q$-semimodule $L$ write
\[
 D=\B\otimes_Q L,\qquad b:L\longrightarrow D,\quad b(l)=1\otimes l.
\]
Here $\B$ is a right $Q$-semimodule through $\chi$. Equivalently, $D$ is
the join-semilattice quotient imposing $ql=l$ for every nonzero $q$.
Every element of $D$ is $b(l)$, and $b(ql)=b(l)$ when $q\ne0$.

\begin{theorem}[Entire support semirings]\label{thm:entire-support-projectives}
The finitely generated projective left $T$-semimodules are classified by
a finitely generated projective left $Q$-semimodule $L$ and a finitely
generated projective left $R$-module $V_p$ for each join-irreducible $p$
of $D=\B\otimes_Q L$. The lattice $D$ is finite and distributive,
although $L$ need not be finite or distributive. Zero labels are allowed.
The corresponding semimodule has support $L$ and groups
\[
 M_l=\bigoplus_{p\le b(l)}V_p,
\]
with the inclusion transitions. A nonzero scalar $(q,r)$ acts by left
multiplication by $r$ into $M_{ql}$, whose displayed summands are the
same as those of $M_l$; the zero scalar acts by the global zero.
Two such data give isomorphic semimodules exactly when a $Q$-isomorphism
of their supports identifies the labels through its induced isomorphism
of Booleanizations.
\end{theorem}
The Booleanization controls the coefficient labels even when the original support is not a distributive lattice. The theorem evaluates this amplitude problem without classifying arbitrary projective $Q$-supports.

Finite support multiplication permits a stronger conclusion than
finiteness of the whole inverse semiring. An ordinary unital ring is
\emph{semiperfect} if its simple modules have projective covers;
equivalently, its identity is a finite sum of orthogonal idempotents
whose corner rings are local.
\begin{theorem}[finite support and semiperfect upper rings]\label{thm:finite-inverse-projectives}
Let $T=(Q,F)$ have finite $Q$, and suppose that every coefficient ring
$F(e)$ with $1\leq e=e^2$ is semiperfect. Let $L$ be a finite left
$Q$-semimodule. For each $l\in L$, let $\delta_l:L\to Q$ be the
pointwise join of all $Q$-linear functionals admissible at $l$.
Then $\delta_l$ is itself admissible and greatest, and $L$ is projective
if and only if
\[
 k=\bigvee_{l\in L}\delta_l(k)l\qquad(k\in L).
\]
When this holds it gives a retraction of $Q^{|L|}$. Put
\[
 r(l,k)=\bigvee\{q\in Q:ql\leq k\}.
\]
The action category and its canonical ideal are
\[
 \mathcal A_L(l,k)=F(r(l,k)),\qquad
 \mathcal J_L(l,k)=\im\bigl(F(\delta_l(k))\longrightarrow F(r(l,k))\bigr),
\]
with the induced coefficient composition and transition map. Thus
\[
 B_L=\bigoplus_{l,k\in L}\mathcal A_L(l,k),\qquad
 I_L=\bigoplus_{l,k\in L}\mathcal J_L(l,k)
\]
is a semiperfect unital ring with a two-sided ideal $I_L$.
Let $S_i$ represent its simple left modules and $P_i$ their projective
covers. If $L$ passes the displayed support test, its finitely generated
projective $T$-semimodules, with identified support $L$, have exactly
the amplitudes
\[
 \bigoplus_{i:\,I_LS_i\ne0}P_i^{\,n_i},\qquad n_i\in\mathbb N.
\]
There are finitely many simple types and the multiplicities are unique.
Forgetting the support identification takes their orbits under
$\operatorname{Aut}_Q(L)$. If the support test fails, no such projective
exists. The zero category ring has no simple types and gives the sole
zero-amplitude vector, still retaining the specified support $L$.
These data classify all finitely generated projectives over $T$.
\end{theorem}
The coefficient groups may be infinite, and the coefficient rings need not
be commutative. Every finite
inverse semiring satisfies the hypotheses; in that case $B_L$ is a
finite ring and is determined by finite tables. There is no asserted
uniform bound on support sizes, nor an algorithm for arbitrary ordinary
ring presentations. The semiperfect and Krull--Schmidt inputs are
classical \cite[Proposition~4.1 and Theorem~4.2]{Krause}.

The action category in the general criterion has an explicit description
directly from the coefficient diagram, without a word presentation.
\begin{theorem}[coefficient colimits]\label{thm:action-colimits}
Let $T=(Q,F)$ be any unital inverse semiring and $L$ any left
$Q$-semimodule. For $l,k\in L$, put
\[
 K(l,k)=\{q\in Q:ql\leq k\}.
\]
This is a nonempty directed poset, and there are canonical identifications
\[
 \mathcal A_L(l,k)=\varinjlim_{q\in K(l,k)}F(q).
\]
The composite of representatives $x\in F(q)$ at $l\to h$ and
$y\in F(p)$ at $h\to k$ is represented by
$\mu_{p,q}(y,x)\in F(pq)$; object identities are represented by the
global scalar unit. Let $D_l$ be the directed poset of $Q$-linear
functionals admissible at $l$, ordered pointwise. Then
\[
 \mathcal J_L(l,k)=\im\left(
  \varinjlim_{\phi\in D_l}F(\phi(k))
  \longrightarrow\varinjlim_{q\in K(l,k)}F(q)\right).
\]
Every ideal element thus has a single admissible functional and
coefficient representative. The displayed map need not be injective.

If $Q$ is finite, with no finiteness assumption on $L$ or the groups,
put $r(l,k)=\bigvee K(l,k)$. The greatest admissible functional
$\delta_l=\bigvee_{\phi\in D_l}\phi$ also exists. In that case
\[
 \mathcal A_L(l,k)=F(r(l,k)),\qquad
 \mathcal J_L(l,k)=
 \im\bigl(F(\delta_l(k))\longrightarrow F(r(l,k))\bigr).
\]
Composition first uses $\mu$ at $r(h,k)r(l,h)$ and then the coefficient
transition to $r(l,k)$. For finite $L$, these are the exact blocks and
multiplication of the ordinary category ring, even when its groups are
infinite.
\end{theorem}
The unrestricted ordinary category and its canonical ideal are therefore
known explicitly from $(Q,F)$ and the support action. Classifying their
finite projectives and the allowed projective supports remains the
residual problem; the theorem does not assert a finite algorithm for
arbitrary infinite data.

For a finite projective support, the functionals in the preceding theorem
can themselves be eliminated from the formula. A \emph{finite support
frame} is a family $l_1,\ldots,l_n\in L$ and left $Q$-linear maps
$\lambda_i:L\to Q$ such that
\begin{equation}\label{eq:support-frame}
 k=\bigvee_{i=1}^n\lambda_i(k)l_i\qquad(k\in L).
\end{equation}
Such a frame exists for every finitely generated projective support.
The empty frame is allowed when $L$ is the zero semimodule.

\begin{theorem}[finite-frame evaluation of the canonical ideal]
\label{thm:finite-frame-ideal}
Let $T=(Q,F)$ be any unital inverse semiring and let $L$ be a finitely
generated projective left $Q$-semimodule with a finite support frame.
$Q$ need not be finite or commutative.
For $l,k\in L$, define the downward-closed set
\begin{equation}\label{eq:frame-C}
 C(l,k)=\mathord\downarrow\left\{
  \bigvee_{i=1}^n\lambda_i(k)a_i:
  a_i\in Q,\ a_i l\leq l_i\text{ for every }i\right\}.
\end{equation}
This is a nonempty finite-join-closed subset of $K(l,k)$, independent of
the chosen frame. The canonical ideal is exactly
\begin{equation}\label{eq:frame-J}
 \mathcal J_L(l,k)=\im\left(
  \varinjlim_{q\in C(l,k)}F(q)
  \longrightarrow\varinjlim_{q\in K(l,k)}F(q)\right).
\end{equation}
Thus, once a finite projective support frame is supplied, evaluating
the ideal requires only the scalar inequalities $a_i l\leq l_i$ and the
given coefficient diagram.

More precisely, every functional admissible at $l$ has the form
\[
 \psi_a(k)=\bigvee_i\lambda_i(k)a_i,
 \qquad a_i l\leq l_i,
\]
and this parametrization is final for directed colimits. If each
$K(l,l_i)$ has a greatest element $r(l,l_i)$, then the greatest
admissible functional is
\begin{equation}\label{eq:frame-delta}
 \delta_l(k)=\bigvee_i\lambda_i(k)r(l,l_i).
\end{equation}
In particular, this holds whenever $Q$ is finite. The order of the
factors in these formulas is part of the assertion.
\end{theorem}

The finite frame is the usual semimodule dual basis
\cite[Dual Basis Lemma, pp.~161--163]{DeorePatil}; the proof below derives
it directly from a finite free retraction. The new formula does not
classify arbitrary projective $Q$-supports or ordinary projectives over
their coefficient categories. It removes the additional quantification
over functions once the support frame is fixed.

There is also a coefficient reduction with no restriction on the support
semiring. A \emph{coefficient ideal} of $(Q,F)$ is a family of additive
subgroups of its coefficient groups, preserved by transitions and by
coefficient multiplication on either side. Write $t_{a,b}$ for the
transition $F(a)\to F(b)$ when $a\leq b$.

\begin{theorem}[an unrestricted coefficient quotient]
\label{thm:upper-unit-quotient}
Let $T=(Q,F)$ be any unital inverse semiring. Define
\[
 K_q=\{x\in F(q):t_{pqr,pqr\vee1}(uxv)=0
 \text{ for every }p,r\in Q,\ u\in F(p),\ v\in F(r)\}.
\]
Then $K$ is the largest coefficient ideal whose group at every
$s\geq1$ is zero. The support-preserving quotient
$T'=(Q,F/K)$ leaves those coefficient groups unchanged and has no
nonzero coefficient ideal with that property.

For every idempotent $n\times n$ support matrix $E$, $n\geq1$, the
kernel $N_E$ of the induced map on additive coefficient ring fibers
$A_E(T)=\prod_{i,j}F(E_{ij})$
satisfies $N_E^{n+1}=0$. Extension of scalars therefore induces a
direct-sum-preserving bijection on isomorphism classes of finitely
generated projective left semimodules and isomorphisms
\[
 V(T)\simeq V(T'),\qquad K_0(T)\simeq K_0(T').
\]
No finiteness, commutativity, integrality or order-height condition is
imposed on $Q$ or its coefficient groups. The assertion concerns
projective isomorphism classes, not equivalence of the full projective
categories.
\end{theorem}

Thus the unrestricted classification can first discard this largest
coefficient ideal and work with pairs having no such remaining ideal.
The context tests may be infinite. The quotient is invariant under
pair isomorphisms; no functoriality under arbitrary homomorphisms or
decision algorithm is asserted. The integral-support reduction below
identifies this quotient by ordinary ideals in one ring.

The preceding reduction can be strengthened using only the upper
multiplicative idempotents of the support semiring. Put
\[
 U(Q)=\{e\in Q:1\leq e,\ e^2=e\}.
\]
The coefficient group $F(e)$ at such an $e$ is an ordinary unital ring,
with unit the image of $1$ under $F(1)\to F(e)$.

\begin{theorem}[detection at upper idempotents]
\label{thm:upper-idempotent-quotient}
For any unital inverse semiring $T=(Q,F)$, define
\[
 K^{\mathrm{id}}_q=\{x\in F(q):t_{pqr,e}(uxv)=0
 \text{ whenever }e\in U(Q),\ pqr\leq e,
 \ u\in F(p),\ v\in F(r)\}.
\]
This is the largest coefficient ideal that vanishes at every support
in $U(Q)$. The quotient $T^{\mathrm{id}}=(Q,F/K^{\mathrm{id}})$
retains all supports and has no nonzero coefficient ideal with that
vanishing property.

For every idempotent $n\times n$ support matrix $E$, the kernel of
$A_E(T)\to A_E(T^{\mathrm{id}})$ has $(n+1)$st power zero.
Extension of scalars gives a direct-sum-preserving bijection on
isomorphism classes of finitely generated projective left semimodules,
and hence
\[
 V(T)\simeq V(T^{\mathrm{id}}),\qquad
 K_0(T)\simeq K_0(T^{\mathrm{id}}).
\]
The ideal $K^{\mathrm{id}}$ contains the ideal of
Theorem~\ref{thm:upper-unit-quotient}, and the containment can be strict.
For integral $Q$, both are $\ker(F(q)\to F(1))$ in each fiber.
\end{theorem}

This is again unrestricted in the support and coefficient data. The
rings $F(e)$ are unital, whereas the additive matrix fibers $A_E$ may
not be. The global coefficient ideal need not be nilpotent: the bound
is inside each fixed idempotent matrix fiber. Isomorphism classes,
rather than the full projective categories, are identified. The
remaining classification may therefore be restricted further to pairs
whose coefficient ideals are detected by these upper idempotent rings.

A further reduction allows ordinary locally nilpotent ideals in the
upper idempotent rings. For a unital ring $A$, let $\operatorname{Lev}(A)$
be its largest locally nilpotent two-sided ideal: every finite subset
of this ideal generates a nilpotent, possibly nonunital, subring.

\begin{theorem}[a compatible locally nilpotent quotient]
\label{thm:upper-levitzki-quotient}
Let $T=(Q,F)$ be any unital inverse semiring. Define $H_q$ by the tests
\[
 x\in H_q\quad\Longleftrightarrow\quad
 t_{pqr,e}(uxv)\in\operatorname{Lev}(F(e))
\]
for every $u\in F(p)$, $v\in F(r)$ and $e\in U(Q)$ with $pqr\leq e$.
Then $H$ is the largest coefficient ideal whose fiber at each
$e\in U(Q)$ is locally nilpotent. It contains $K^{\mathrm{id}}$, and
$T/H=(Q,F/H)$ has no nonzero coefficient ideal with that same property.

For every finite idempotent support matrix $E$, the kernel of
$A_E(T)\to A_E(T/H)$ is locally nilpotent. Extension of scalars gives
a direct-sum-preserving bijection of finitely generated projective
left-semimodule isomorphism classes, and hence
\[
 V(T)\simeq V(T/H),\qquad K_0(T)\simeq K_0(T/H).
\]
For every $m\geq1$, this coefficient ideal commutes with finite matrices:
\[
 H(M_m(T))_P=\prod_{i,j}H(T)_{P_{ij}}
 \qquad(P\in M_m(Q)).
\]
Consequently $M_m(T)/H(M_m(T))\simeq M_m(T/H(T))$ canonically.
\end{theorem}

The nilpotence assertion is local inside each ordinary matrix fiber;
there is no bound depending only on matrix size. For an ordinary ring
$T=R$, the quotient is the familiar $R/\operatorname{Lev}(R)$, so the
ordinary radical and nil-ideal lifting are known inputs. The new
coefficient condition concerns compatible ideals. It does not make
every individual upper idempotent ring radical-free. Exact supports
are retained, and no equivalence of full projective categories,
arbitrary-map functoriality or arbitrary Morita-equivalence assertion
is made. The unrestricted support and amplitude classifications
remain to be evaluated after this further coefficient reduction.


Proposition~\ref{prop:upper-transition-radical} below evaluates this
same ideal using only ordinary-ring radicals and upward transitions.
Every coefficient group uses two-by-two upper-idempotent fiber rings;
for finite $Q$ the intersection is finite, even when its coefficient
groups are infinite. This is a formula, not a general radical algorithm.

The quotients of Theorems~\ref{thm:upper-unit-quotient} and~\ref{thm:upper-idempotent-quotient} coincide and have a particularly
explicit form in the following broad class. Call an idempotent support semiring \emph{integral} here if its
multiplicative unit is its greatest element. This use of the word does
not mean absence of zero divisors.

\begin{theorem}[removing transition kernels]
\label{thm:integral-support-reduction}
Let $T=(Q,F)$ be a unital inverse semiring with integral support $Q$.
Put $R=F(1)$ and $I_q=\im(F(q)\to R)$. Then the $I_q$ form an
increasing diagram of two-sided ideals of the unital ring $R$, with
$I_1=R$ and $I_qI_p\subseteq I_{qp}$. Their inclusions and ordinary
ring products define an inverse semiring $T_{\mathrm{sep}}=(Q,I)$ and
a surjective homomorphism $T\to T_{\mathrm{sep}}$ that preserves every
support. This is precisely either of these two zero-detection quotients.

For every idempotent $n\times n$ support matrix $E$, $n\geq1$, the
additive matrix fiber $A_E=\prod_{i,j}F(E_{ij})$ is a possibly
nonunital ring. The kernel $N_E$ of its map to
$\prod_{i,j}I_{E_{ij}}$ satisfies
\[
 N_E^{\,n+1}=0.
\]
Consequently extension of scalars induces a bijection on isomorphism
classes of finitely generated projective left semimodules, preserving
direct sums. In particular,
\[
 V(T)\simeq V(T_{\mathrm{sep}}),\qquad
 K_0(T)\simeq K_0(T_{\mathrm{sep}}).
\]
There is no finiteness, height or commutativity hypothesis on $Q$, and
no finiteness or commutativity hypothesis on the coefficient groups or
the unit ring. The assertion concerns projective isomorphism classes;
it does not assert equivalence of the full projective categories.
\end{theorem}

For integral support, transition kernels therefore create no additional
finite-projective classes. The residual coefficient data are actual
ideals in one ordinary ring. Projectives over that ideal diagram still
require classification; the preceding theorem retains its unrestricted scope.

When the unit ring is local, the ordinary projective data in this
reduction have a complete multiplicity description. Here a local ring
is nonzero and may be noncommutative.

\begin{theorem}[local unit rings]
\label{thm:local-unit-projectives}
Suppose $T=(Q,F)$ has integral support and $R=F(1)$ is local. Put
$I_q=\im(F(q)\to R)$ and $W=\{q:I_q=R\}$. For a finitely generated
projective left $Q$-support $L$, define
\[
 l\preceq_W k\quad\Longleftrightarrow\quad
 ql\leq k\text{ for some }q\in W,
\]
and let $l\sim_W k$ mean comparability in both directions. A class
$[l]$ is \emph{allowed} if a functional $\psi$ admissible at $l$
satisfies $\psi(l)\in W$, equivalently if $C(l,l)\cap W\ne\varnothing$.
This condition is independent of its representative.

Finite-projective left $T$-semimodules are classified by $L$ and a
finitely supported function
\[
 m:\{\text{allowed }\sim_W\text{ classes in }L\}\longrightarrow\N,
\]
up to $Q$-support isomorphisms transporting these multiplicities.
For the separated diagram of Theorem~\ref{thm:integral-support-reduction},
choose representatives of the classes with positive multiplicity. Its
group over $k$ is
\begin{equation}\label{eq:local-unit-fibers}
 G_k=\bigoplus_{[l]}
 \left(\bigcup_{ql\leq k}I_q\right)^{m([l])}.
\end{equation}
Transitions are inclusions and scalars act by left multiplication in
$R$. The original $T$-class is recovered by lifting an idempotent; it
is not obtained by an assertion that restriction of scalars preserves
projectivity.
\end{theorem}

The theorem evaluates the ordinary amplitude classification for every
integral support with local unit ring, including infinite supports and
infinite coefficient groups. Classifying the projective $Q$-supports
themselves remains a separate problem. Zero multiplicities are allowed;
the resulting zero-amplitude module still has its specified support.

The finite-support hypothesis in Theorem~\ref{thm:finite-inverse-projectives}
can be removed if one describes the ordinary projective types by local
idempotents in the canonical action ideal. A unital ring is
\emph{semiperfect} if its identity is a finite sum of orthogonal
idempotents with local corner rings; the zero ring is allowed.

\begin{theorem}[amplitudes with semiperfect upper coefficients]
\label{thm:upper-semiperfect-amplitudes}
Let $T=(Q,F)$ be a unital inverse semiring such that $F(e)$ is
semiperfect for every $e\in U(Q)$. For any $Q$-semimodule $L$, put
$\mathcal A=\mathcal A_L$ and $\mathcal J=\mathcal J_L$. Consider pairs
$(l,f)$ with
\[
 0\ne f=f^2\in\mathcal J(l,l),\qquad
 f\mathcal A(l,l)f\text{ a local ring}.
\]
Declare $(l,f)\sim(k,g)$ if there are arrows $u:l\to k$ and
$v:k\to l$ in $\mathcal A$ satisfying
\[
 u=guf,\qquad v=fvg,\qquad vu=f,\qquad uv=g.
\]
Every finitely generated projective $\mathcal A$-module $G$ with
$\mathcal JG=G$ is a finite direct sum of
\[
 P_{l,f}=\mathcal A(l,-)f.
\]
Its multiplicities on the displayed equivalence classes are unique,
and every such finite sum occurs. Consequently finitely generated
projective $T$-semimodules are classified by a finitely generated
projective $Q$-support $L$ and these finite multiplicities, modulo
$\operatorname{Aut}_Q(L)$.

If every upper coefficient ring is local or zero, the permitted pairs
reduce to objects $l$ with $\mathcal A(l,l)\ne0$ and
$1_l\in\mathcal J(l,l)$, with $f=1_l$. The multiplicities are then
indexed by their isomorphism classes as objects of $\mathcal A$.
\end{theorem}

There is no finiteness assumption on $Q$, $L$ or the coefficient
groups, and no commutativity assumption. The empty sum is zero
amplitude over the given support. The theorem determines the
amplitude parameter in this class; it leaves the classification of
projective $Q$-supports and the evaluation of the displayed
idempotents and intertwiners as separate questions. In particular,
it does not assert that every action endomorphism ring is
semiperfect; Example~\ref{ex:upper-local-action-boundary} shows why.


The next quotient has an unconditional isomorphism-lifting property.
Lifting objects requires an additional hypothesis.

\begin{theorem}[a compatible Jacobson quotient]
\label{thm:upper-jacobson-quotient}
For any unital inverse semiring $T=(Q,F)$, define
\[
 x\in C_q\quad\Longleftrightarrow\quad
 t_{pqr,e}(uxv)\in\operatorname{Jac}(F(e))
\]
for every $u\in F(p)$, $v\in F(r)$ and $e\in U(Q)$ with
$pqr\leq e$. Then $C$ is the largest compatible coefficient ideal
with $C_e\subseteq\operatorname{Jac}(F(e))$ at every upper
idempotent. It contains $H$, and $T/C$ has no nonzero compatible
coefficient ideal with this property. It commutes entrywise with
every finite matrix construction.

For each finite idempotent support matrix $E$, the kernel $N_E$ of
$A_E(T)\to A_E(T/C)$ is a Jacobson radical, possibly nonunital,
ring. For any idempotent $P\in A_E(T)$,
\[
 P N_E P\subseteq\operatorname{Jac}(P A_E(T)P).
\]
Every isomorphism between reductions of two finitely generated
projective $T$-semimodules lifts to an isomorphism. Thus
$V(T)\to V(T/C)$ is injective, and a quotient projective has at
most one lifting isomorphism class.

If all $F(e)$, $e\in U(Q)$, are semiperfect, every quotient
projector lifts at its exact support. In that case
\[
 V(T)\simeq V(T/C),\qquad K_0(T)\simeq K_0(T/C).
\]
There is no finiteness or commutativity hypothesis on $Q$ or $F$.
\end{theorem}

The unrestricted assertion gives neither existence of all lifting
objects nor an injection on group completions. The ordinary
semilocal domain in Example~\ref{ex:jacobson-lifting-boundary}
already prevents unconditional object lifting. This is a quotient
of coefficient groups retaining $Q$; it is not the Bourne
quotient by either of the usual global semiring Jacobson radicals.


Corollary~\ref{cor:jacobson-transition-tests} evaluates this same
coefficient ideal by upward ordinary-ring transitions and explicit
two-by-two fiber rings, with finitely many tests when $Q$ is finite.

The action-colimit description also admits an exact finite model for
each specified family of projective amplitudes, without a
semiperfectness hypothesis.

\begin{theorem}[finite coefficient models for projective amplitudes]
\label{thm:finite-coefficient-models}
Let $T=(Q,F)$ be any unital inverse semiring, $L$ any
$Q$-semimodule, $\mathcal A=\mathcal A_L$ and
$\mathcal J=\mathcal J_L$. Every finite family of finitely
generated projective $\mathcal A$-modules $G_s$ with
$\mathcal JG_s=G_s$ has a simultaneous model as follows.

Choose finitely many $l_i\in L$ and admissible functionals
$\phi_j$ at $l_j$, and put
\[
 q_{ij}=\phi_j(l_i),\qquad
 d_{ii}=1\vee q_{ii},\qquad d_{ij}=q_{ij}\quad(i\ne j).
\]
Then $D=(d_{ij})$ is an upper idempotent support matrix, and its
ordinary coefficient ring $C=A_D(T)$ maps unitally to the full
action-arrow matrix ring $B_{ij}=\mathcal A(l_j,l_i)$.
Let $N\subseteq C$ have arbitrary off-diagonal coefficients in
$F(q_{ij})$ and diagonal coefficients in the image of
$F(q_{ii})\to F(d_{ii})$. It is an ideal, its image lies in the
canonical arrow ideal, and, for the indicated map $\pi$,
\[
 (\ker\pi)N=0.
\]
There are idempotents $e_s\in N$ whose amplitude images are
$G_s$. For any idempotents $e,f\in N$, including these, the map
\[
 eCf\xrightarrow{\ \sim\ }\pi(e)B\pi(f)
\]
is an additive isomorphism compatible with composition. For
$e=f$ it is a unital ring isomorphism. Conversely every such
idempotent yields a finite projective $\mathcal J$-generated
amplitude.
\end{theorem}

Thus all morphisms and isomorphisms in any specified finite family
are computed in an actual ordinary coefficient ring, with the
usual opposite convention for covariant representables. The
chosen ring depends on that family. No universal ring, algorithm
for finding its admissible data, or classification of arbitrary
ordinary projectives or projective $Q$-supports is asserted.
Example~\ref{ex:no-fixed-finite-objects} shows that even a fixed
finite set of action objects need not detect all projective
amplitudes on one free support.


A unital ring $B$ is an \emph{exchange ring} if, for every
$b\in B$, there is an idempotent $p\in Bb$ with
$1-p\in B(1-b)$. This is the ordinary exchange condition
\cite[Proposition~1.1 and Theorem~2.1]{NicholsonExchange}; it
includes semiperfect rings. It permits a stronger object-lifting
statement than the Jacobson case alone.

\begin{theorem}[coefficient quotients with exchange upper rings]
\label{thm:exchange-coefficient-lifting}
Let $T=(Q,F)$ be any unital inverse semiring such that every
$F(e)$, $e\in U(Q)$, is an exchange ring. Let $\mathcal I$
be any compatible coefficient ideal, and let $T/\mathcal I$
retain the support $Q$.

For every finite idempotent support matrix $E$, every idempotent
of $A_E(T/\mathcal I)$ lifts to an idempotent of $A_E(T)$.
Thus the lift has exactly the prescribed support, even when
that additive matrix fiber is nonunital. Every finitely generated
projective over $T/\mathcal I$ lifts, and the maps
\[
 V(T)\longrightarrow V(T/\mathcal I),\qquad
 K_0(T)\longrightarrow K_0(T/\mathcal I)
\]
are surjective. For the compatible Jacobson ideal $C(T)$ of
Theorem~\ref{thm:upper-jacobson-quotient}, both maps are
isomorphisms.
\end{theorem}

No nilpotence or radical condition is imposed on $\mathcal I$.
Injectivity for an arbitrary coefficient ideal is not asserted.
The local-idempotent multiplicity classification in
Theorem~\ref{thm:upper-semiperfect-amplitudes} retains its
semiperfectness hypothesis.

The finite-family construction can be made uniform. An
\emph{admissible functional at $l$} is a $Q$-linear map
$\phi:L\to Q$ satisfying $\phi(k)l\leq k$ for every $k$.
Let $X_L$ consist of all pairs $x=(l_x,\phi_x)$ of this kind.
For $x,y\in X_L$, set
\[
 q_{yx}=\phi_x(l_y),\qquad
 d_{yx}=\begin{cases}q_{yx},&x\ne y,\\1\vee q_{xx},&x=y.\end{cases}
\]

\begin{theorem}[a uniform category of coefficient fibers]
\label{thm:uniform-coefficient-category}
For every unital inverse semiring $T=(Q,F)$ and every
$Q$-semimodule $L$, the groups
\[
 \mathcal C_L(x,y)=F(d_{yx})\qquad(x,y\in X_L)
\]
form a small unital preadditive category. Composition is
coefficient multiplication followed by the indicated upward
transition. It has the two-sided ideal
\[
 \mathcal N_L(x,y)=
 \begin{cases}
 F(q_{yx}),&x\ne y,\\
 \im\bigl(F(q_{xx})\to F(1\vee q_{xx})\bigr),&x=y.
 \end{cases}
\]
The quotient category has zero arrows between distinct objects
and diagonal rings
\[
 F(1\vee q_{xx})/\im F(q_{xx}).
\]
There is a unital coefficient-to-action functor
$\pi:\mathcal C_L\to\mathcal A_L$, $x\mapsto l_x$.
Its left Kan extension gives an equivalence
\[
 \left\{P\in\operatorname{proj}(\mathcal C_L):
                  \mathcal N_LP=P\right\}
 \simeq
 \left\{G\in\operatorname{proj}(\mathcal A_L):
                  \mathcal J_LG=G\right\},
\]
where $\operatorname{proj}$ denotes finitely generated
projective covariant additive modules. The source is equivalently
the kernel, on finite projectives, of extension to
$\mathcal C_L/\mathcal N_L$.

The equivalence preserves all morphisms and is equivariant for
$u\in\operatorname{Aut}_Q(L)$ acting by
$(l,\phi)\mapsto(u(l),\phi u^{-1})$. If $L$ is finitely
generated projective over $Q$, these data classify finite
projective $T$-semimodules with specified support $L$;
unmarked classes take the indicated automorphism orbits.
\end{theorem}

This category is fixed before choosing a projective amplitude.
Every morphism group is one coefficient group, and its
endomorphism rings are upper coefficient rings. Neither its
object set nor its ideal is assumed finite; the ideal need
not be idempotent or have local units. The remaining problems
are projective $Q$-supports and ordinary projectives in the
explicit kernel above, with no general algorithm asserted.


\subsection*{Credit and neighboring results}
The construction $R_0$ is \cite[Definition~3.6]{FGJM}; the inverse additive structure of every module is \cite[Proposition~6.1]{FGJM}. Their Corollary~4.12 proves that Construction~3.5 exhausts the inverse semirings with two additive idempotents. Their Proposition~5.1 identifies $\Z_0$ as the initial inverse semiring, and their Proposition~5.2 describes free commutative inverse semirings. The description by groups over a semilattice is Clifford's structure theorem \cite{Clifford}. Caviglia, Faul, Manuell and Mesiti recover its categorical form in \cite[\S\S4--5]{CFMM}, including the equivalences for fixed and varying semilattices, and their Theorem~8.5 adds semiring multiplication. Their Remark~8.7 identifies the two-idempotent case with \cite[Construction~3.5]{FGJM}. We use their covariant join convention from \S8. The gluing in Theorem~\ref{thm:twolevelprojectives} is analogous to projective patching for ring fiber products \cite[Lemma~15.6.9, Tag~0D2G]{Stacks}. Its proof below constructs the required semimodule retraction directly. The nilpotent correction is the standard idempotent-lifting polynomial \cite[Lemma~10.32.7, Tag~05BU]{Stacks}, applied inside the matrix group fiber; the proof checks that this is a ring with a square-zero kernel.

Sen, Bhuniya and Maity \cite[Abstract]{SBM} characterize Clifford semimodules as strong semilattices of modules and credit Yusuf's 1966 inverse-semimodule work. This is a substantial antecedent to the diagram viewpoint. For fields, Dokuchaev, Johari and Vilca-Rodr\'iguez \cite[Definitions~2.2--2.3 and Lemma~2.4, pp.~1641--1643]{DJV} give a precise neighboring framework. Their inverse semivector spaces whose addition is a monoid, with zero-preserving linear maps, are exactly ordinary $F_0$-semimodules: restrict the action to $F$ in one direction; in the other, let $\tau$ act by the global additive zero. Their Lemma~2.4 fixes every additive idempotent under old scalars, so this extension satisfies the zero and mixed scalar laws. Their presheaf reconstruction and fixed-support equivalence for semilattices of algebras \cite[Propositions~3.10 and~3.12, pp.~1656--1657]{DJV} provide further context. Their meet convention reverses the join order used here.

The public foundation note \cite[Theorems~1--3 and Remark~5.4]{FoundationPublic} contains the zero-adjoined-ring spectrum, module and projective core, with the corrected ordinary-module interpretation. The broader classifications and general ideal criterion are stated separately here.

The finite Boolean-semifield case of projectivity is the projective-semilattice theorem of Horn and Kimura \cite[Theorem~5.3]{HK}. Macpherson \cite[Theorem~3.9, pp.~8--9]{Macpherson} also describes Boolean projectives as modules freely generated by lower-finite posets; in the finite case this is the order-ideal description used here. For commutative regular semigroups, Bulman-Fleming and McDowell \cite[Theorem~3.3, pp.~255--258]{BM} already give a projective classification by the semilattice and its component groups. Their maps need not preserve an identity. The zero least group here records the extra identity-preserving convention. Our ring-coefficient proof is included because this distinction and the labelled-poset form are needed for the Grothendieck monoid. Katsov, Nam and Zumbr\"agel \cite[Remark~6.6(3), Proposition~6.7]{KNZ} prove that $K_0(\B)$ contains a free abelian subgroup of countable rank; Theorem~\ref{thm:Kzero} describes its whole canonical basis by connected finite posets. Di Nola, Lenzi and Nam \cite[Abstract]{DLN} describe strong projectives over finite products of matrix semirings over commutative chain semirings, and use their ordered $SK_0$-groups. Theorem~\ref{thm:iteratedprojective} also treats ordinary projectives over the finite max--min chains and over their ring-valued top fibers. Rowen's quasi-zeros \cite{Rowen} are precisely the additive idempotents $m+(-m)$ used below. Connes and Consani's characteristic-one homological algebra \cite{CC} supplies the context for the generic $\B$-fiber. The sheaf statement uses ordinary semimodules and does not identify that fiber with their arithmetic site.

\section{Proof overview}
The old zero $e$ acts on every semimodule as a projection to its additive idempotents. The fiber of this projection over an idempotent is an abelian group, and the old ring acts linearly there. Addition moves two vectors to the group indexed by the join of their supports. This proves the diagram description directly, while the same distinction between old and new zero gives the ideals and congruences.

For projectives, a retraction from a finite free semimodule has two parts. Its action on idempotents retracts a finite Boolean lattice and produces a finite distributive lattice. Its linear actions on group fibers have split injections along the order. Coordinate restriction proves both the sum and intersection laws for these fibers. Splitting at join-irreducible elements then gives the projective labels. Conversely an explicit idempotent block matrix realizes every labelled poset. Direct sums become disjoint unions of posets, so connected components determine the Grothendieck monoid without further relations.

Repeated zero adjunction changes the Boolean support lattice of a free module to a product of finite chains. In a retract, truncation to each chain level is intersection with a fixed order ideal. A coordinate detecting a join-irreducible also detects all of its positive truncations, forcing the greatest-element condition in Theorem~\ref{thm:iteratedprojective}. Explicit section and retraction maps prove that the condition is sufficient. Only the top colour can carry a nonzero group, so the same splitting proof then supplies all ring-module labels.

For the two-idempotent construction, first replace the lower bimodule by its image ideal. A free fiber then embeds in its top ring module and is the inverse image of a coordinate submodule modulo that ideal. The labelled-poset classification over the quotient ring determines all fibers. To prove sufficiency, lift a quotient splitting and correct its error, which lands in the ideal, by adding free coordinates whose section has bottom support. The kernel of the original bimodule map is then removed by square-zero lifting within each idempotent-support matrix group. This proves existence and uniqueness of projective isomorphism classes without requiring a direct-sum complement for a general projective semimodule.

For sheaves, the crucial topological fact is that every nonempty open meets the open point $\eta$. A support label must therefore be the same in every piece of a compatible family. Only the ring-module sections remain to be glued. Their usual localization and affine-scheme theory prove the result. The general inverse-semiring statement retains this fiberwise construction and replaces scalar multiplication by natural bilinear maps indexed by the idempotent action.

For an arbitrary inverse semiring, a support frame separates the support retraction from ordinary module projectivity over an action category. The canonical ideal records exactly which amplitude generators can be supplied by a free semimodule. To replace action-arrow colimits by coefficient groups, pair each support point with an admissible scalar functional. These pairs form a single category whose off-diagonal arrows are already coefficient groups; adjoining the support unit only to diagonal groups supplies identities. The kernel of its action functor annihilates its distinguished ideal. This one-sided annihilation both lifts projectors by squaring and makes every mixed projective corner exact. Consequently the category models all amplitudes on the chosen support, without resolving the remaining ordinary projective classification.

\section{Ideals, congruences and functoriality}\label{sec:spectrum}
This section gives the algebraic facts that determine the spectrum and its structure sheaf.

\begin{proposition}\label{prop:ideal}
The ideals and primes are those in Theorem~\ref{thm:main}(1), and every ideal is subtractive. The semiring $R_0$ is zerosumfree and has no zero divisors relative to its new zero.
\end{proposition}
\begin{proof}
A sum or product of old elements lies in $R$, and therefore never equals $\tau$. This proves the last assertions and shows that $\{\tau\}$ is prime. If an ideal $I$ contains $r\in R$, it contains $er=e$ and $(-1)r=-r$. Hence $I\cap R$ is a ring ideal and $I=(I\cap R)\cup\{\tau\}$. The converse follows directly from the operations. For an old ring ideal $J$, primality of $J\cup\{\tau\}$ is equivalent to primality of $J$, because the only products still to consider have both factors in $R$.

If $a\in I$ and $a+b\in I$, the cases $a=\tau$ or $b=\tau$ are immediate. Otherwise $a,b\in R$ and subtraction in $I\cap R$ gives $b\in I$. The ideal $\{\tau\}$ is subtractive by zerosumfreeness.
\end{proof}

\begin{proposition}\label{prop:localization}
For $s\in R$ one has $S_s\simeq(R_s)_0$. At an old prime $\mathfrak p$ this gives $S_{\mathfrak p\cup\{\tau\}}\simeq(R_{\mathfrak p})_0$. The localization $S_{\eta}$ is $\B$. The nonempty opens are exactly $\{\eta\}\cup U$, and
\[
 \mathcal O_X(\{\eta\}\cup U)\simeq\Gamma(U,\mathcal O_Y)_0,
\]
where $\Gamma(\varnothing,\mathcal O_Y)$ is the zero ring. If $R\ne0$ has finite dimension, $\dim S=\dim R+1$.
\end{proposition}
\begin{proof}
In a fraction with an old denominator, a new-zero numerator stays the new zero. Two old numerators are equivalent exactly when they are equivalent in ring localization, and no old numerator becomes equivalent to $\tau$: the multiplier and denominator in the defining fraction equivalence are old. This proves the localization formulas.

At $\eta$, all old elements are inverted, including $e$. From $e^2=e$ one obtains $e=1$, and from $er=e$ every old $r$ becomes $1$. The new zero remains distinct, so the resulting semiring is $\B$. For $s\in R$, the distinguished open is $\{\eta\}\cup D_R(s)$, and $D_S(e)=\{\eta\}$. This proves the description of opens. On distinguished opens the asserted section formula is the localization formula. It defines a sheaf on all opens: any compatible family has one common choice between its new-zero section and an old section, because every two nonempty members meet at $\eta$; the old sections glue uniquely in $\mathcal O_Y$. It is consequently the structure sheaf. Finally every prime chain in $R$ can be preceded by $\eta$, and these are all possible chains.
\end{proof}

\begin{proposition}\label{prop:congruences}
The congruences of $S$ are exactly the two families in Theorem~\ref{thm:main}(2). The map $\pi:S\to R$ that fixes $R$ and sends $\tau$ to $e$ is universal among maps from $S$ to rings.
\end{proposition}
\begin{proof}
Restrict a congruence to the old ring. It is a ring congruence, since multiplication by $-1$ respects it, and thus has the form $r\equiv r'$ if and only if $r-r'\in J$ for a unique ideal $J$. If $\tau$ is a singleton class, this gives exactly $(R/J)_0$. If $\tau$ is congruent to an old $r$, multiplying by $e$ gives $\tau\equiv e$. The quotient then factors through $R$ and is $R/J$. Conversely both specified relations respect the operations. Their kernels on $R$ recover $J$, and the class of $\tau$ distinguishes the families.

For a map to a ring, the additive idempotent $e$ must map to zero. It therefore has the same image as $\tau$, which proves the universal property. This is also the ring reflection of \cite[Proposition~4.13]{FGJM} in this case.
\end{proof}

\begin{proposition}\label{prop:functor}
Zero adjunction is fully faithful on commutative unital rings. For an injective ring homomorphism $R\hookrightarrow K$,
\[
 R_0\simeq R\times_K K_0.
\]
For arbitrary homomorphisms of commutative unital rings $A\to C\leftarrow B$,
\[
 (A\times_C B)_0\simeq A_0\times_{C_0}B_0.
\]
In particular, taking $C=0$ gives $(A\times B)_0\simeq A_0\times_{\B}B_0$ using the support maps.
The common fixed points of multiplication by $R^\times$ on $R_0$ are
\[
 \Ann_R\bigl((u-1:u\in R^\times)\bigr)\sqcup\{\tau\}.
\]
They are exactly $\{e,\tau\}$ if and only if this annihilator is zero.
\end{proposition}
\begin{proof}
Let $h:R_0\to T_0$ preserve zero and unit. If $h(e)=\tau$, then $h(1)+h(-1)=\tau$, contrary to zerosumfreeness and $h(1)=1$. Thus $h(e)$ is the old zero, the only other additive idempotent. No old $r$ can map to $\tau$, because $er=e$ would give $\tau=h(e)$. Hence the restriction is a ring homomorphism. Every ring homomorphism extends uniquely, proving full faithfulness and, in particular, $\operatorname{Aut}(R_0)=\operatorname{Aut}(R)$.

In the first pullback an old $r$ corresponds to $(r,r)$, and the new zero to $(e,\tau)$. Injectivity of $R\to K$ ensures that these are all pairs, and their operations agree. For the second pullback, an old element maps to an old element of $C_0$, whereas a new zero maps to its new zero. Equality of the images therefore forces both coordinates to be old or both to be new zero. The old pairs are exactly $A\times_C B$, and the only remaining pair is $(\tau,\tau)$. Componentwise operations give the displayed isomorphism. This argument includes the zero-ring cases. For fixed points, an old $r$ is fixed by every unit exactly when $(u-1)r=0$ for every unit $u$. The new zero is always fixed.
\end{proof}

The two-element set $\{e,\tau\}$ has Boolean operations, but its multiplicative identity is $e$; its inclusion in $R_0$ is not unital unless $R=0$. The zero ring gives $(0)_0=\B$, and the support map $R_0\to\B$ is the image of $R\to0$ under zero adjunction. Even powers of a unit need not lie in its fixed-point set: in $\Z_0$, one has $(-1)^2=1$, but $(-1)\cdot1=-1\ne1$.

\section{Semimodules as diagrams}\label{sec:diagrams}
This section proves Theorem~\ref{thm:main}(3) and identifies the localization seen by the new point.

\begin{lemma}\label{lem:fibers}
For an $S$-semimodule $M$, put $-m=(-1)m$ and $\epsilon(m)=em$. Its additive monoid is inverse, its additive idempotents are $L=\epsilon(M)$, and $M_l=\epsilon^{-1}(l)$ is an $R$-module with zero $l$.
\end{lemma}
\begin{proof}
Scalar distributivity gives
\[
 m+(-m)+m=m,\qquad (-m)+m+(-m)=-m.
\]
For completeness, inverses in a commutative regular monoid are unique. If $b,c$ are inverses of $a$, then $a+b$ and $a+c$ are idempotents, and
\[
 b=b+a+b=b+(a+c+a)+b=(b+a)+(a+b)+c=a+b+c.
\]
Similarly $c=a+b+c$, so $b=c$. This proves the inverse assertion, as in \cite[Proposition~6.1]{FGJM}. The inverse of an idempotent is itself. Consequently an additive idempotent $l$ satisfies $\epsilon(l)=l+(-l)=l$, while every $\epsilon(m)$ is idempotent. Addition on $L$ is a join operation with least element $0_M$.

If $\epsilon(m)=l$, then $m+l=m$ and $m+(-m)=l$. The fiber is closed under addition and inverse, so it is an abelian group with zero $l$. For $r\in R$, $\epsilon(rm)=erm=em=l$. The semimodule axioms restricted to the fiber give the ring-module axioms, with $em=l$ as the zero-scalar identity.
\end{proof}

\begin{proof}[Proof of the diagram equivalence]
For $l\le k$ define $\rho_{l,k}(m)=m+k$. This is linear: $rk=k$ for every old scalar $r$, because $k$ is the zero of its own group. It preserves the fiber zeros, respects composition, and is the identity when $l=k$. If $m\in M_l$ and $n\in M_k$, their sum lies in $M_{l\vee k}$ and is the sum of their images there.

Conversely, from a diagram $A:L\to R\text{-}\Mod$ form the disjoint union of its underlying groups, writing its elements as $(l,a)$. Set
\[
 (l,a)+(k,b)=\bigl(l\vee k,\rho_{l,l\vee k}(a)+\rho_{k,l\vee k}(b)\bigr).
\]
The zero is $(0_L,0)$, old $r$ acts as $(l,a)\mapsto(l,ra)$, and $\tau$ sends everything to $(0_L,0)$. Coherence of the diagram makes addition associative; the other additive axioms are immediate. Ring linearity proves distributivity for old scalars, and the cases involving $\tau$ follow from the definition. Each group zero $(l,0)$ is an idempotent and $e(l,a)=(l,0)$, so the recovered diagram is the original one. A semimodule map preserves $em$, hence induces a join map and linear fiber maps; conversely such compatible maps preserve the displayed formulas. The two constructions are mutually inverse up to their evident natural identifications.
\end{proof}

\begin{proposition}\label{prop:genericmodule}
For $s\in R$, localization of an $R_0$-semimodule with diagram $(L,A)$ preserves $L$ and replaces $A_l$ by $(A_l)_s$. At $\eta$ it is exactly $L$, viewed as a $\B$-semimodule. The cancellative $R_0$-semimodules are the ordinary $R$-modules restricted along $\pi:R_0\to R$.
\end{proposition}
\begin{proof}
Multiplication by an old scalar preserves the support label. Fraction equivalence therefore never identifies different labels, and within one fiber it is the usual module fraction equivalence. Inverting $e$ replaces every vector by its group zero and hence gives $L$. Finally an additive idempotent in a cancellative monoid must be its global zero, so the diagram has one fiber and is a ring module. A ring module is cancellative and defines an $R_0$-semimodule through $\pi$.
\end{proof}

For $R=\Z$, every commutative inverse monoid acquires the unique scalar action given by repeated addition, its inverse operation, and the new zero. Thus the diagram equivalence specializes to Clifford's description of commutative inverse monoids. The uniqueness follows because these operations determine the actions of every integer and of $\tau$.

\begin{proposition}[normalized derivations]\label{prop:derivations}
Let $A\to R$ be a unital homomorphism of commutative rings, and let $M$ be an $R_0$-semimodule with least group $M_{0_L}$. A derivation is required to preserve the additive zero, be additive, satisfy the Leibniz rule, and have $d(1)=0_M$. Then naturally in $M$,
\[
 \operatorname{Der}_{A_0}(R_0,M)
 \simeq \operatorname{Der}_{A}(R,M_{0_L}).
\]
The relative differential semimodule is therefore the ordinary module $\Omega^1_{R/A}$ regarded as an $R_0$-semimodule through ring reflection. Every derivation into a projective $R_0$-semimodule is zero; in particular every derivation $R_0\to R_0$ is zero.
\end{proposition}
\begin{proof}
The identity $1+e=1$ and the normalization give $d(e)=0_M$. Hence
\[
 d(r)+d(-r)=d(e)=0_M
\]
for each $r\in R$. Thus every $d(r)$ belongs to the group with identity the global zero, namely $M_{0_L}$. The restriction to $R$ is an ordinary $A$-derivation. Conversely an ordinary derivation into this group extends by $d(\tau)=0_M$; the laws involving $\tau$ follow because it acts by the global zero. These operations are inverse.

For any ordinary $R$-module $V$, an $R_0$-linear map $V\to M$ has image in $M_{0_L}$, since it preserves additive inverses relative to the global zero. Conversely an $R$-linear map $V\to M_{0_L}$ is $R_0$-linear. Applying this observation to the universal property of $\Omega^1_{R/A}$ proves the representing-object assertion. The least group of a free $R_0$-semimodule is zero, and the same is true of any retract, proving the projective-target assertion.
\end{proof}

For example, with $A=\Z$ and $R=\Z[x]$, the differential module is the nonzero ordinary $R$-module $R\,dx$, with one support point. Thus vanishing of derivations into the scalar semiring does not imply vanishing of the universal differential module. The single support point here has a nonzero group, as permitted in Theorem~\ref{thm:main}(3).

\section{Finite projectives and their Grothendieck monoids}\label{sec:projectives}
This section proves the labelled-poset classification and its consequences for $K_0$.

\begin{lemma}\label{lem:lattice}
A retract of a finite Boolean lattice by zero-and-join preserving maps is a finite distributive lattice. Conversely every finite distributive lattice is such a retract.
\end{lemma}
\begin{proof}
Let $i:L\to2^{[n]}$ and $r:2^{[n]}\to L$ have $ri=1$. A finite join-semilattice with zero has meets, obtained by joining all common lower bounds. If $p$ is join-irreducible, then
\[
 p=r(i(p))=\bigvee_{j\in i(p)}r(\{j\}),
\]
so $p=r(\{j\})$ for some $j\in i(p)$. If $p\le a\vee b$, this $j$ belongs to $i(a)\cup i(b)$, and therefore $p\le a$ or $p\le b$. Every join-irreducible is consequently join-prime. Let $P$ be their poset. The map $l\mapsto\{p\in P:p\le l\}$ is injective because every element of a finite lattice is a join of join-irreducibles. It is onto $\Id(P)$: for an ideal $I$, primality implies that the join-irreducibles below $\bigvee I$ are exactly those in $I$. Hence $L\simeq\Id(P)$ is distributive. Conversely $\Id(P)\subseteq2^P$ is retracted by taking downward closure. This is the finite zero-preserving form of \cite[Theorem~5.3]{HK}.
\end{proof}

\begin{proof}[Proof of Theorem~\ref{thm:main}(4)]
Suppose $M$ is a retract of $F=S^n$, with maps $i:M\to F$ and $r:F\to M$. Its idempotent semilattice $L$ is a retract of $2^{[n]}$, so Lemma~\ref{lem:lattice} identifies it with $\Id(P)$ for a finite poset $P$. Write $i_0,r_0$ for the maps on idempotents. The fiber of $F$ over $A\subseteq[n]$ is $R^A$. Its maps to higher fibers extend vectors by old zero.

Each group $M_I$ is a retract of $R^{i_0(I)}$, hence a finitely generated projective $R$-module. Each transition $M_I\to M_J$, $I\subseteq J$, is a split injection: compose $i$ with coordinate projection from $i_0(J)$ to $i_0(I)$ and then with $r$. Naturality shows that this composite is a left inverse. The least fiber is a retract of $R^{\varnothing}=0$ and is therefore zero. We identify the fibers, by these injections, with submodules of the top group.

The submodules satisfy
\begin{equation}\label{eq:latticegroups}
 M_{I\cup J}=M_I+M_J,\qquad M_I\cap M_J=M_{I\cap J}.
\end{equation}
For the sum, the coordinate set of $i_0(I\cup J)$ is $i_0(I)\cup i_0(J)$. Split an embedded vector into one vector on $i_0(I)$ and a second on $i_0(J)$, assigning overlap coordinates to the first and old zeros to the second. Applying $r$ proves the asserted equality. For the intersection, a common vector in the top free group has zero coordinates outside $i_0(I)\cap i_0(J)$. Restrict it to this intersection and apply $r$. Its support is at most $I\cap J$, since
\[
 r_0(i_0(I)\cap i_0(J))\le I,J.
\]
Embedding back into the top group recovers the original vector. The reverse inclusion is automatic. This proves \eqref{eq:latticegroups} without assuming that $i_0$ preserves meets.

For $p\in P$, put $I_p=\{q:q\le p\}$ and $I_p^-=I_p\setminus\{p\}$. The injection $M_{I_p^-}\to M_{I_p}$ splits, so its quotient $P_p$ is finitely generated projective. Choose a section $P_p\to M_{I_p}$. For any ideal $I$, list its elements in a linear extension. Adding one element $p$ uses
\[
 I_{\rm old}\cup I_p=I_{\rm new},\qquad I_{\rm old}\cap I_p=I_p^-.
\]
The two identities in \eqref{eq:latticegroups} show that the chosen section of $P_p$ complements $M_{I_{\rm old}}$ inside $M_{I_{\rm new}}$. Induction gives
\[
 M_I=\bigoplus_{p\in I}P_p.
\]
The sections were chosen once for each $p$, so these identifications commute with every transition. This proves necessity.

For sufficiency choose an idempotent matrix $A_p\in M_{n_p}(R)$ with image $P_p$, taking $n_p\ge1$ even if $P_p=0$. Use coordinate blocks indexed by $P$ and define a matrix $E$ over $S$ by
\[
 E_{pp}=A_p,\qquad
 E_{pq}=\begin{cases}\text{the all-}e\text{ block},&p<q,\\
                         \text{the all-}\tau\text{ block},&p\not\le q.
          \end{cases}
\]
The diagonal blocks of $E^2$ are $A_p^2=A_p$. If $p<q$, every old contribution to the $(p,q)$ block is an old-zero block, and contributions through the diagonal blocks exist because their sizes are positive; thus the result is again the all-$e$ block. If $p\not\le q$, every contribution is the new zero. Therefore $E^2=E$.

The support of $Ex$ is the downward closure of the set of coordinate blocks containing an old entry of $x$. Over an ideal $I$, its ring vector has component in $\im A_p=P_p$ for each $p\in I$. Conversely choose representatives in $R^{n_p}$ for these vectors, supported on every coordinate of blocks in $I$, and put $\tau$ in other blocks. Their image is the prescribed element. Thus $\im E=\mathcal P(P,(P_p))$, proving projectivity.

An isomorphism preserves the lattice of idempotents and its join-irreducibles. It also preserves the quotient
\[
 P_p\simeq M_{I_p}/M_{I_p^-}.
\]
Hence it preserves the labels up to a poset isomorphism. The converse follows directly from \eqref{eq:labelled}.
\end{proof}

\begin{example}
Let $P=\{a<b\}$ and label both points by $R$. The diagram has groups $0,R,R^2$ over the three ideals $\varnothing,\{a\},\{a,b\}$. It is the image of
\[
 \begin{pmatrix}1&e\\ \tau&1\end{pmatrix}:S^2\longrightarrow S^2.
\]
Its three-element idempotent lattice shows that it is not free. A single point with label zero gives $\B$, a projective $S$-semimodule through the support action: $S\to\B$ is split by $1\mapsto e$ as a semimodule map.
\end{example}

\begin{proof}[Proof of Theorem~\ref{thm:Kzero}]
The product of the diagrams for $P$ and $Q$ has idempotent lattice
\[
 \Id(P)\times\Id(Q)=\Id(P\sqcup Q),
\]
and its fibers have exactly the original labels. Thus addition in $V(S)$ is disjoint union of labelled posets. Every finite poset is the disjoint union of the connected components of its comparability graph, uniquely up to their order and isomorphism. The classification proves that there are no other relations. This gives $V(S)=\N^{(\mathcal C_R)}$, and group completion gives the free abelian group on the same basis.

The free semimodule $S^n$ corresponds to $n$ isolated points, each labelled $R$. If $P\oplus Q\simeq S^n$, the labelled posets of $P$ and $Q$ form a disjoint union of precisely these points. Therefore $P\simeq S^k$ for some $k$. Distinct $k$ give different numbers $2^k$ of idempotents, proving $SV(S)\simeq\N$ and $SK_0(S)\simeq\Z$.

Finally any homomorphism from a finite-support diagram to an ordinary ring module kills the idempotents and factors uniquely through the top group. For \eqref{eq:labelled} that group is $\bigoplus_{p\in P}P_p$. This is the universal property of extension of scalars along $S\to R$.
\end{proof}

\begin{remark}
For $R=0$, all labels are zero and $K_0(\B)$ is free abelian on nonempty connected finite posets. For a field $k$, labels are nonnegative integers, their dimensions. Already one-point posets with different dimensions are different basis elements. For general rings, whenever distinct projective modules have the same class in $K_0(R)$, the one-point labelled posets still distinguish their isomorphism classes in $V(R_0)$ and $K_0(R_0)$.
\end{remark}

\subsection{Tensor products and rank characters}
The following uses the ordinary tensor product of semimodules over the
commutative semiring $R_0$. Write $\operatorname{Pic}$ for the group of
isomorphism classes of tensor-invertible semimodules.

\begin{proposition}\label{prop:tensor-picard}
For a commutative unital ring $R$, including the zero ring,
\[
 \mathcal P(P,(M_p))\otimes_{R_0}\mathcal P(Q,(N_q))
 \simeq\mathcal P(P\times Q,(M_p\otimes_R N_q)).
\]
Zero labels are allowed. Let $\mathcal C_R$ be the commutative monoid of
isomorphism classes of nonempty connected finite $V(R)$-labelled posets,
with this product and unit the singleton labelled $R$. Then, as a
semiring and a ring respectively,
\[
 V(R_0)\simeq\N[\mathcal C_R],\qquad
 K_0(R_0)\simeq\Z[\mathcal C_R],\qquad
 \operatorname{Pic}(R_0)\simeq\operatorname{Pic}(R).
\]
If $\mathfrak p\in\Spec R$, let $r(M)$ be the rank of $M$ at
$\mathfrak p$. For every $s\in\mathbb C$ the formula
\[
 D_s^+(\mathcal P(P,(M_p)))
   =\sum_{\substack{p\in P\\r(M_p)>0}}r(M_p)^{-s}
\]
extends to a ring homomorphism $K_0(R_0)\to\mathbb C$.
For actual projectives, $D_{-1}^+\ge D_0^+$.
\end{proposition}
\begin{proof}
Choose a split free presentation of each $M_p$, with projector $a_p$
on $R^{n_p}$ and $n_p\ge1$; use $n_p=1$ and $a_p=0$ for a zero label.
The construction in the proof of Theorem~\ref{thm:main}(4) presents its
semimodule by a matrix with full occupied support block from $q$ to $p$
exactly when $p\le q$, and ring-amplitude matrix block diagonal with
blocks $a_p$. Thus the other occupied entries have old-zero amplitude.
Tensor the two split presentations. The free basis is the Cartesian
product of the original bases, and the projector is their Kronecker
product. Its support order is the product order, while its diagonal
ring block is $a_p\otimes b_q$, with image $M_p\otimes_R N_q$.
This proves the tensor formula, including zero labels and empty posets.
A product of two nonempty connected posets is connected: move along a
comparability path in one coordinate with the other fixed, and then in
the other coordinate. Direct sum is disjoint union by
Theorem~\ref{thm:Kzero}. These observations give the monoid semiring
and its ring completion.

Every tensor-invertible semimodule $X$ is dualizable, with its tensor
inverse $Y$. Write the coevaluation of $1$ as a finite sum
$\sum_i x_i\otimes y_i$. The triangle identity gives
$x=\sum_i x_i\operatorname{ev}(y_i\otimes x)$. Hence the coefficient
maps $x\mapsto\operatorname{ev}(y_i\otimes x)$ split the surjection
$R_0^n\to X$ defined by the $x_i$: $X$ is finite projective.
The same holds for $Y$. The tensor formula now forces their support
posets to be singletons, since their product is the singleton poset
of $R_0$. Their labels $M,N$ satisfy $M\otimes_R N\simeq R$.
Conversely such invertible ring modules give inverse singleton-labelled
semimodules. This proves the Picard comparison. These projectives have
a zero least group; they are not the ordinary ring modules obtained
by restriction along the ring reflection.

Local rank is a nonnegative integer, is additive on ring direct sums,
and satisfies $r(M\otimes_R N)=r(M)r(N)$. A product rank is positive
exactly when both ranks are positive. The tensor formula therefore
factors the displayed sum, while disjoint union makes it additive.
The unit has rank one, so group completion gives the asserted unital
ring homomorphism. Finally each retained positive rank is at least one,
proving the inequality on actual projectives.
\end{proof}
The restriction to positive ranks is necessary. A singleton with zero
label has total rank zero but one support vertex; counting that vertex
at $s=0$ would give a false inequality. At positive real $s$ its
unrestricted expression $0^{-s}$ is undefined. Even a nonzero module
may have rank zero at a chosen prime of a disconnected ring. No rank
character is asserted for the zero ring, whose prime spectrum is empty.


\section{Quasi-coherent sheaves}\label{sec:sheaves}
This section proves the affine equivalence in the semiring-scheme convention fixed in Section~1.

\begin{proof}[Proof of Theorem~\ref{thm:QC}]
Let $L\to\Qcoh(Y)$ be a diagram. Define a presheaf on $X$ by \eqref{eq:sheaf} on nonempty opens and by the one-element module on the empty open. Its restrictions keep the label and restrict the ordinary module section. The action of $\mathcal O_X(\{\eta\}\cup U)=\mathcal O_Y(U)_0$ is the ring action on each fiber, with the new zero sending everything to the global zero.

Consider a cover of a nonempty open by nonempty opens and a compatible family of sections. Each member restricts at $\eta$ to its label in $L$. Thus all labels coincide, and the family is a compatible family of sections of one $\mathcal A_l$. It glues uniquely because $\mathcal A_l$ is a sheaf. This proves the sheaf axiom. Since $Y$ is affine, write $\mathcal A_l=\widetilde{A_l}$ for the module $A_l=\Gamma(Y,\mathcal A_l)$; naturality identifies the diagram maps as localizations of module maps. Proposition~\ref{prop:genericmodule} shows on every distinguished open that the sheaf just constructed is $\widetilde{M}$ for the diagram semimodule $M$. It is therefore quasi-coherent.

Conversely let $\mathcal M$ be quasi-coherent on $X$. Its restriction to the open point $\eta$ is a $\B$-semimodule, hence a join-semilattice $L$. On a distinguished open $V=\{\eta\}\cup D_R(s)$ where $\mathcal M$ is associated to a semimodule, Proposition~\ref{prop:genericmodule} identifies its support semilattice with this same $L$. The inverse operation and multiplication by $e$ are defined locally and commute with restrictions. The idempotent section for each $l\in L$ is unique locally and therefore glues. The sections restricting to $l$ at $\eta$ form an abelian group with this idempotent as identity. Their old scalar action is an $\mathcal O_Y$-module action. On overlaps these groups and actions agree, because both are defined from $\mathcal M$ itself. They consequently define sheaves $\mathcal A_l$ on $Y$, with linear maps obtained by adding the idempotent for a larger label.

On the indicated distinguished cover these are the sheaves associated to the localized ring-module fibers. Hence they are quasi-coherent $\mathcal O_Y$-modules. Every section of $\mathcal M$ on a nonempty open has one label, and the preceding construction identifies its group fiber with $\Gamma(U,\mathcal A_l)$. Thus it recovers \eqref{eq:sheaf}. Morphisms preserve the global generic labels and restrict to compatible linear sheaf maps; conversely those data preserve the sheaf formulas. This proves the equivalence of categories. Taking $U=Y$ in \eqref{eq:sheaf} and using the ordinary affine equivalence proves that association and global sections are inverse.
\end{proof}

If $R=0$, then $Y$ is empty and $X$ is the single Boolean point. The diagram on $Y$ has a zero group at every label, so the assertion reduces to the equivalence between $\B$-semimodules and join-semilattices with zero. This also checks the empty-open convention in the theorem.

\section{Modules over an arbitrary inverse semiring}\label{sec:general}
This section writes the scalar action in the same fiberwise language as the inverse-semiring structure theorem of Caviglia, Faul, Manuell and Mesiti.

Let $Q$ be an idempotent semiring, ordered by $q\le q'$ if $q+q'=q'$. Let $F:Q\to\Ab$ be lax monoidal for multiplication, with maps
\[
 \mu_{q,r}:F(q)\otimes_{\Z}F(r)\longrightarrow F(qr),
 \qquad \iota:\Z\longrightarrow F(1_Q).
\]
Let $T$ be the inverse semiring of \cite[Theorem~8.5]{CFMM}: its elements are $(q,a)$, its addition transports to joins, its product is $(qr,\mu_{q,r}(a,b))$, its additive zero is $(0_Q,0)$, and its multiplicative unit is $(1_Q,\iota(1))$. Multiplication need not be commutative in this section.

\begin{definition}\label{def:actiondata}
Module data over $(Q,F)$ consist of an idempotent left $Q$-semimodule $L$, a functor $G:L\to\Ab$, and bilinear maps
\[
 \alpha_{q,l}:F(q)\otimes_{\Z}G(l)\longrightarrow G(ql).
\]
These maps are natural in both variables: for $q\le q'$ and $l\le l'$, transporting the input or the output gives the same map to $G(q'l')$. They satisfy
\begin{align*}
 \alpha_{1_Q,l}(\iota(1),m)&=m,\\
 \alpha_{qr,l}(\mu_{q,r}(a,b),m)
 &=\alpha_{q,rl}\bigl(a,\alpha_{r,l}(b,m)\bigr).
\end{align*}
A morphism is a $Q$-linear map $u:L\to L'$ with natural group maps $G(l)\to G'(u(l))$ commuting with the action maps.
\end{definition}

\begin{theorem}[module form of the structure theorem]\label{thm:general}
Left $T$-semimodules are equivalent to the module data of Definition~\ref{def:actiondata}. There is no requirement that $G(0_L)$ be zero.
\end{theorem}
\begin{proof}
By \cite[Proposition~6.1]{FGJM}, or by the inverse identities in Lemma~\ref{lem:fibers}, a $T$-semimodule is additively a commutative inverse monoid. Let $L$ be its idempotents and $G(l)$ its group fibers. Clifford's construction gives the functor $G$, with transitions $m\mapsto m+l'$.

If $a\in T$ has idempotent $q=a+(-a)$ and $l$ is an idempotent of the module, then $al$ is idempotent, and $(-a)l=al$ by uniqueness of its additive inverse. Hence $ql=al+(-a)l=al$. Also, for $m\in G(l)$, the idempotent of $am$ is $q l$. Thus the action on module idempotents depends only on $q$, and makes $L$ a $Q$-semimodule. The scalar action restricts to the displayed group-bilinear maps.

The two naturality identities follow by distributing across the addition of a scalar idempotent or of a module idempotent. The unit and associative laws are inherited from the scalar action. This constructs the stated data.

Conversely form the disjoint union of the $G(l)$, add by transport to joins, and define
\[
 (q,a)(l,m)=(ql,\alpha_{q,l}(a,m)).
\]
The support of each distributive law is the corresponding distributive law for the $Q$-action on $L$. Once transported to that common support, naturality and bilinearity prove the equality of group components. Unit and associativity are exactly the two displayed axioms. The additive zero acts by zero because its group component is zero and $0_Q l=0_L$; every scalar sends the global module zero to itself because $q0_L=0_L$ and the action is bilinear. The reconstruction therefore is a semimodule. Its group fibers, transports and scalar action recover the original data. The same formulas prove the assertion on morphisms.
\end{proof}

\subsection{The two-idempotent construction}
Suppose now that $R$ is any unital ring, $A$ an $R$--$R$-bimodule, and $f:A\to R$ is a bimodule map with $f(a)b=af(b)$. The semiring $I(f)=R\sqcup A$ of \cite[Construction~3.5]{FGJM} has the operations given in Section~1. All modules in this subsection are left modules.

\begin{corollary}\label{cor:two}
An $I(f)$-semimodule is equivalently a diagram of $R$-modules $G:L\to R\text{-}\Mod$ and natural $R$-linear maps
\[
 \theta_l:A\otimes_R G(l)\longrightarrow G(0_L)
\]
such that, for $\phi_l=G(0_L\le l)$,
\[
 \phi_l\bigl(\theta_l(a\otimes m)\bigr)=f(a)m.
\]
Naturality here means $\theta_{l'}(a\otimes\rho_{l,l'}m)=\theta_l(a\otimes m)$.
\end{corollary}
\begin{proof}
In Theorem~\ref{thm:general}, take $Q=\B$, $F(0)=A$, $F(1)=R$ and the transition $f$. Every $\B$-semimodule is a join-semilattice with zero. The upper action makes the groups $R$-modules. The lower action gives the maps $\theta_l$, whose target is $G(0_L)$. Naturality in $0\le1$ is precisely the displayed identity. Naturality in $l$ is the other stated condition. At $l=0_L$ the identity says $\theta_0(a,m)=f(a)m$. Associativity of two lower scalars follows because
\[
 \theta_0(a,\theta_l(b,m))=f(a)\theta_l(b,m)
 =\theta_l(f(a)b,m).
\]
The mixed associative laws are left $R$-linearity and balancing over the right $R$-action on $A$. Thus these conditions are both necessary and sufficient.
\end{proof}

For $A=0$ this is the diagram theorem for $R_0$. For $A=R$ and $f=1$, the semiring is $R\times\B$ and its semimodule category is $R\text{-}\Mod\times\B\text{-}\Mod$. Indeed the orthogonal multiplicative idempotents $(1,0)$ and $(0,1)$ sum to the unit and decompose every module into their images. Its finitely generated projectives are therefore pairs of finitely generated projective $R$-modules and finite distributive lattices. This illustrates why labels that are invariant for $R_0$ cannot be carried unchanged to every two-idempotent inverse semiring.

\subsection{Finite projectives in fiberwise matrix form}
The following complete description is useful for concrete computations; it is not an intrinsic normal form for arbitrary $Q$. Matrices in this subsection act on free right $T$-semimodules. Replacing $T$ by its opposite gives the corresponding description for left semimodules; the distinction disappears in the commutative cases.

\begin{proposition}\label{prop:matrix}
Let $E=(a_{ij})\in M_n(T)$ be idempotent, write $a_{ij}=(b_{ij},v_{ij})$, and define
\[
 \beta(q)_i=\bigvee_j b_{ij}q_j\qquad(q\in Q^n),\qquad
 H(q)=\bigoplus_j F(q_j).
\]
Let $E_q:H(q)\to H(\beta q)$ be the group map whose $i$th component is the sum of $\mu_{b_{ij},q_j}(v_{ij},x_j)$ transported to $F((\beta q)_i)$. Then $\beta^2=\beta$, and $\im E$ has semilattice $\operatorname{Fix}(\beta)$ and group
\[
 G(l)=\im(E_l:H(l)\to H(l))\quad(l\in\operatorname{Fix}(\beta)),
\]
with the inherited transitions and actions. Every finitely generated projective is of this form. Two idempotents $E\in M_n(T)$ and $E'\in M_m(T)$ give isomorphic projectives exactly when there are matrices $U,V$ with
\[
 U=E'UE,\quad V=EVE',\quad VU=E,\quad UV=E'.
\]
\end{proposition}
\begin{proof}
The free semimodule $T^n$ has precisely the semilattice and fibers $Q^n,H(q)$. The formula for matrix multiplication is the stated formula for its fiber map. Thus $E^2=E$ gives $\beta^2=\beta$ and $E_{\beta q}E_q=E_q$. The idempotents in its image are exactly the fixed points of $\beta$. If an image vector has support $l$, applying $E$ again fixes it and so it belongs to $\im E_l$. Conversely every vector in $\im E_l$ is an image vector with that support. Naturality of the free-module maps gives the inherited transitions and actions. Every finite projective is the image of an idempotent because it is a retract of a finite free semimodule. Finally an isomorphism between two images extends to the free modules by the retractions and inclusions, giving $U,V$ as written. Conversely those equations restrict to mutually inverse maps on the images.
\end{proof}

In Construction~3.5 the free group over $I\subseteq[n]$ is $A^{[n]\setminus I}\oplus R^I$. Thus Proposition~\ref{prop:matrix} is an explicit description in terms of the given ring, module, and map $f$; it uses no unknown semimodule operations.

\subsection{An additive-category reduction and its support obstruction}
\label{sec:ringoid}

Proposition~\ref{prop:matrix} puts support and amplitude in one idempotent
equation. We now separate an ordinary additive-category projectivity problem
from a precise obstruction to lifting its support. This first criterion does not eliminate the search for a support section.
Theorem~\ref{thm:canonical-ideal} below replaces that search by a canonical ideal condition, while leaving arbitrary projective $Q$-supports unclassified.

We return to left semimodules. Fix a $Q$-semimodule $L$. Define a preadditive category $\mathcal A_L$
with objects $l\in L$, generated by arrows
\[
 t_{l,k}:l\longrightarrow k\quad(l\leq k),\qquad
 a_{q,r,l}:l\longrightarrow ql\quad(r\in F(q)).
\]
Write $F_{q,q'}$ for a transition of $F$. The relations are
\begin{align*}
 t_{l,l}&=1_l,& t_{k,h}t_{l,k}&=t_{l,h},\\
 a_{q,r+r',l}&=a_{q,r,l}+a_{q,r',l},&
 a_{q',F_{q,q'}(r),l}&=t_{ql,q'l}a_{q,r,l}\quad(q\leq q'),\\
 a_{q,r,k}t_{l,k}&=t_{ql,qk}a_{q,r,l}\quad(l\leq k),&
 a_{q,r,pl}a_{p,s,l}&=a_{qp,\mu_{q,p}(r,s),l},\\
 a_{1_Q,\iota(1),l}&=1_l.&&
\end{align*}
A left $\mathcal A_L$-module is an additive covariant functor
$\mathcal A_L\to\Ab$. This is the usual rings-with-several-objects
formalism \cite{Mitchell,Murfet}; all of the projectivity arguments needed below
are given explicitly. We use projectivity with respect to surjections,
as in \cite[Definition~7.1, p.~18]{BorgerJun}.

\begin{proposition}[fixed support]\label{prop:ringoid}
The category of $T$-semimodules with identified support $L$, and maps
inducing $1_L$, is equivalent to $\mathcal A_L$-modules.
\end{proposition}
\begin{proof}
The transition arrows give the group diagram, and the scalar arrows give
the action homomorphisms in Definition~\ref{def:actiondata}. The transition and naturality relations give the group diagram and its
compatibility with the action; the other relations give scalar
additivity, associativity, and the unit. Scalar addition across
different supports follows by transporting both summands to their join
and then using additivity there. Module addition across different fibers
follows in the same way from transition naturality. Thus
Theorem~\ref{thm:general} reconstructs the semimodule, and its morphism
statement gives the asserted equivalence. The zero object of this
additive category has the zero group at every $l$; its underlying
$T$-semimodule is $L$, which need not be the zero semimodule.
\end{proof}

\begin{proposition}[finite generation]\label{prop:ringoid-finite}
A $T$-semimodule $M$, with support $L$ and amplitude functor $G$, is
finitely generated if and only if $L$ is finitely generated over $Q$
and $G$ is finitely generated over $\mathcal A_L$.
\end{proposition}
\begin{proof}
The supports of finite semimodule generators generate $L$. The same
elements generate $G$: each scalar multiple is a scalar arrow, and a
sum is formed by transitions to a common fiber followed by group addition.

Conversely choose finite $\mathcal A_L$-generators of $G$ and finite
$Q$-generators $l_i$ of $L$. Use the former elements and the fiber zeros
$0_{l_i}$ as semimodule generators. All fiber zeros are generated, because
zero amplitudes in $F(q)$ move $0_{l_i}$ to $0_{ql_i}$ and addition joins
their supports. A transition to a larger fiber is realized by adding its
fiber zero. The scalar $-1_T$ gives the group inverse at the same
support. Therefore every additive-category expression in the selected
generators is an expression in the generated semimodule.
\end{proof}

Choose any finite free cover
\[
 f:T^n\twoheadrightarrow M,\qquad
 (s_j)\longmapsto\sum_j s_jm_j,
 \qquad m_j\in G(l_j).
\]
Its support map is
\[
 \rho:Q^n\twoheadrightarrow L,\qquad
 (q_j)\longmapsto\bigvee_jq_jl_j.
\]
For a $Q$-linear section $\sigma:L\to Q^n$, let
\[
 N_\sigma(l)=\bigoplus_{j=1}^nF\bigl(\sigma(l)_j\bigr).
\]
The free-module transitions and actions restrict along $\sigma$ to make
$N_\sigma$ an $\mathcal A_L$-module. Restriction of $f$ is an
$\mathcal A_L$-linear map $c_\sigma:N_\sigma\to G$. Define
\[
 C_\sigma=\operatorname{coker}(c_\sigma)
 \quad\text{in }\mathcal A_L\text{-}\Mod.
\]

\begin{theorem}[support-section criterion]\label{thm:ringoid-projective}
A finitely generated $T$-semimodule $M$ is projective if and only if
\begin{enumerate}
 \item[\textup{(i)}] $L$ is a finite projective $Q$-semimodule;
 \item[\textup{(ii)}] $G$ is a finitely generated projective $\mathcal A_L$-module;
 \item[\textup{(iii)}] some $Q$-linear section $\sigma$ of $\rho$ satisfies $C_\sigma=0$.
\end{enumerate}
The criterion is independent of the selected finite free cover. For a
fixed $\sigma$, condition \textup{(iii)} is equivalent to the finitely
many linear image conditions
\[
 m_j\in\im\bigl(c_\sigma(l_j)\bigr)\qquad(1\leq j\leq n).
\]
The existence of $\sigma$ already implies \textup{(i)}; listing it
separately isolates the first support-theoretic test.
\end{theorem}
\begin{proof}
Suppose $M$ is projective. Any finite free cover $f$ splits. The support
$\sigma$ of a splitting satisfies $\rho\sigma=1_L$, so $L$ is a
finite free retract. On group fibers the splitting lands in
$N_\sigma(l)$ and splits $c_\sigma$. Hence $C_\sigma=0$.

Let $E\twoheadrightarrow G$ be an epimorphism of $\mathcal A_L$-modules.
It is pointwise surjective, so reconstruction gives a surjective
$T$-map onto $M$ with support $1_L$. Projectivity of $M$ splits it.
The support of this splitting must also be $1_L$, and therefore the
splitting is $\mathcal A_L$-linear. Every epimorphism onto $G$ splits,
which proves projectivity of $G$. Its finite generation follows from
Proposition~\ref{prop:ringoid-finite}.

Conversely suppose \textup{(ii)} and \textup{(iii)} hold.
Projectivity of $G$ splits the pointwise epimorphism $c_\sigma$ by an
$\mathcal A_L$-map $u:G\to N_\sigma$. Reconstruct this as a $T$-map
with support $1_L$, and then include each restricted free fiber into
$T^n$ with support map $\sigma$. The resulting $h:M\to T^n$
satisfies $fh=1_M$ on groups and supports. Thus $M$ is a finite free
retract. Finally, $\im c_\sigma$ is an $\mathcal A_L$-submodule and
the $m_j$ generate $G$, proving the finite image test. Projectivity of
$G$, not arbitrary choices of pointwise lifts, supplies the coherence
of the section.
\end{proof}

\begin{proposition}[a nonzero support obstruction]\label{prop:ringoid-obstruction}
For $T=R_0$ with $R\ne0$, projectivity of the support and of its
amplitude module does not imply projectivity over $T$.
\end{proposition}
\begin{proof}
Take $L=\B^2$ with atoms $a,b$ and top $t=a\vee b$. Here
$\mathcal A_L$ is the $R$-linear incidence category. The diagram
$G(t)=R$, $G(l)=0$ for $l\ne t$, is its projective representable
born at $t$. Its semimodule is generated by $0_a,0_b$ and $1\in G(t)$.

For any finite free cover and any support section,
$\sigma(t)=\sigma(a)\vee\sigma(b)$. Every active free coordinate
at $\sigma(t)$ is therefore active at $\sigma(a)$ or $\sigma(b)$.
The free group at $\sigma(t)$ is the sum of the transition images
of these two lower groups. Their images under $f$ are zero, since
$G(a)=G(b)=0$. Thus $c_\sigma(t)=0$ and
$C_\sigma(t)=R\ne0$ for every $\sigma$. Theorem~\ref{thm:ringoid-projective}
excludes projectivity.
\end{proof}

For $G_{\operatorname{Id}(C)}(R)$, the action category is again the
$R$-linear incidence category of its support: all proper scalar group
fibers vanish. Its finite projective diagrams may have birth labels at
arbitrary objects. Theorem~\ref{thm:finite-lattice-projectives}
identifies exactly which births pass the additional support test,
namely the admissible ideals. The next subsection replaces the presentation-dependent support-section obstruction by a canonical ideal for arbitrary $(Q,F)$.

\subsection{The action category as coefficient colimits}
\label{sec:action-colimits}
We prove Theorem~\ref{thm:action-colimits}. All colimits in this subsection
are taken in abelian groups.
\begin{proof}
The set $K(l,k)$ contains $0_Q$, is downward closed, and is closed under
finite joins: if $ql,q'l\leq k$, then
$(q\vee q')l=ql\vee q'l\leq k$. Hence it is directed.
Define a category with objects $L$ and prospective arrow groups
\[
 H(l,k)=\varinjlim_{q\in K(l,k)}F(q).
\]
Every element has one coefficient representative $[q,x]$, because a
finite sum can be transported to the join of its indices and added
there. Two representatives agree exactly when their images agree at a
common later index; this follows directly from the directed-colimit
relations.

Define composition by
\[
 [p,y]\,[q,x]=[pq,\mu_{p,q}(y,x)].
\]
Its index is valid: $ql\leq h$ and $ph\leq k$ imply $pql\leq k$.
Naturality of $\mu$ makes the definition independent of transported
representatives, and bilinearity follows after moving finite sums to
common indices. Associativity and the global unit follow from the lax
monoidal axioms. The unit coefficient in $F(1_Q)$ has valid index at
$(l,l)$ since $1_Ql=l$. If its class is zero, the resulting object is
correctly a zero object. Thus $H$ is a preadditive category.

Send a scalar generator $a_{q,x,l}$ to $[q,x]$, and a transition
$t_{l,k}$ to the class of the global unit at support $1_Q$.
Every defining relation of $\mathcal A_L$ holds: transitions compose
by the unit law, coefficient transitions are the colimit relations,
scalar-transition naturality uses the same coefficient at the same
index, and scalar products use the displayed composition. Additivity
at a fixed scalar support is group addition; at different supports it
follows by transport to their join. This gives a functor
$\mathcal A_L\to H$.

Conversely assign
\[
 [q,x]\longmapsto t_{ql,k}a_{q,x,l}\quad\text{in }\mathcal A_L(l,k).
\]
For $q\leq q'$ in $K(l,k)$, the coefficient-transition relation gives
\[
 t_{q'l,k}a_{q',F_{q,q'}(x),l}
 =t_{q'l,k}t_{ql,q'l}a_{q,x,l}
 =t_{ql,k}a_{q,x,l}.
\]
Thus the map descends to the colimit. Scalar-transition naturality and
scalar composition show that it preserves composition, with order
$pq$ and coefficient order $y,x$, without commuting them. The two
functors are inverse on all coefficient representatives and on the
scalar and transition generators. This proves the exact hom-group
identification, rather than only a surjection.

The poset $D_l$ contains zero and is directed by pointwise finite
joins. Such joins preserve $Q$-linearity and admissibility by
distributivity. Evaluation $\phi\mapsto\phi(k)$ maps it into
$K(l,k)$. Under the preceding identification, the image of a coefficient
$x\in F(\phi(k))$ is precisely the defining canonical generator
$t_{\phi(k)l,k}a_{\phi(k),x,l}$. The image of the displayed colimit
map is an additive subgroup containing all these generators and no
others beyond their sums, so it equals $\mathcal J_L(l,k)$.
Directedness combines every finite sum at one admissible functional.
No injectivity of any coefficient map is used.

If $Q$ is finite, $K(l,k)$ has the greatest element $r(l,k)$, and its
colimit is $F(r(l,k))$. The inequality
$r(h,k)r(l,h)\leq r(l,k)$ follows from its action on $l$, and yields
the stated multiplication formula.

For arbitrary $L$, the pointwise join of all admissible functionals
still exists when $Q$ is finite. At each argument there are only
finitely many distinct values. Hence finite distributivity proves that
this join preserves zero, binary joins and the left $Q$-action.
For its admissibility at a fixed $k$, select finitely many functionals
realizing those distinct values; their joined value acting on $l$ is
at most $k$. The resulting map is therefore the greatest member
$\delta_l$ of $D_l$. No finite global subfamily of functionals is
asserted. The second colimit is consequently $F(\delta_l(k))$, and
its map to the first is the coefficient transition to $F(r(l,k))$.
This proves the finite-$Q$ formulas with arbitrary $L$ and $F$.
\end{proof}
The maximum $r(l,k)$ is the scalar-action residual. For the classical
residuation convention, compare Cohen, Gaubert and Quadrat
\cite[Sections~2.1--2.2, equations~(2)--(3) and~(13)--(14) of the
preprint]{CGQ}; their displayed semimodule convention is right-sided.
Here it is the left-action residual. Completeness is not assumed in the
general theorem: the colimits replace unavailable maxima.

\begin{corollary}[ordinary group-valued modules]
\label{cor:group-projectives-general}
Let
\[
 R(T)=\varinjlim_{q\in Q}F(q),\qquad
 I_0=\im\bigl(F(0_Q)\longrightarrow R(T)\bigr).
\]
Coefficient multiplication makes $R(T)$ the universal ordinary ring
reflection of $T$. Semimodules with singleton additive support are
exactly ordinary $R(T)$-modules. Such a module $P$ is finitely generated
projective over $T$ if and only if it is finitely generated projective
over $R(T)$ and $I_0P=P$.
\end{corollary}
\begin{proof}
The poset $Q$ is directed by finite joins. Naturality and associativity
of $\mu$ make its group colimit a unital ring. The map sending
$(q,x)\in T$ to the class of $x$ preserves both operations, zero and
unit. Every semiring map from $T$ to a ring sends additive idempotents
to zero, so its restrictions to the groups $F(q)$ are compatible.
They induce a unique unital ring map from this colimit, proving the
universal property.

For the singleton-zero support $L=\{o\}$, one has $K(o,o)=Q$, so
$\mathcal A_L=R(T)$. Its only $Q$-linear functional to $Q$ is zero;
therefore $\mathcal J_L=I_0$. The fixed-support equivalence and
Theorem~\ref{thm:canonical-ideal} prove the remaining statements,
since $L$ is a finite projective $Q$-support. This includes $R(T)=0$.
\end{proof}
For $T=R_0$ the ideal $I_0$ vanishes, recovering
the one-idempotent obstruction in Section~\ref{sec:projectives}; for the two-level ideal
construction it recovers Proposition~\ref{prop:canonical-arbitrary-ideal}.

\subsection{Finite support frames and the scalar ideal}
\label{subsec:finite-frame-ideal}

We prove Theorem~\ref{thm:finite-frame-ideal} and distinguish its scalar
ideal from the ordinary coefficient ideal. The finite dual-basis
argument is classical for semimodules \cite{DeorePatil}. For commutative
semirings, finite-projective dualizability is also described by
Niefield and Wood \cite[Corollary~2.2 and its proof, pp.~1223--1225]{NW}.
Niefield's quantale theorem \cite[Theorem~4.4, p.~848]{Niefield} relates
projectivity to total continuity in a category whose modules and maps
preserve arbitrary joins. Here only finite joins are required, and $Q$
need not be a quantale. We use no transfer of projectivity between these
two categories.

\begin{proof}[Proof of Theorem~\ref{thm:finite-frame-ideal}]
Choose a splitting $L\xrightarrow{s}Q^n\xrightarrow{r}L$ with $rs=1$.
Let $\lambda_i$ be the $i$th coordinate of $s$ and let $l_i=r(e_i)$.
Then $k=r(s(k))=\bigvee_i\lambda_i(k)l_i$, proving the frame identity.
In particular, each $\lambda_i$ is admissible at $l_i$.

Fix $l$. For a tuple $a_i\in K(l,l_i)$, the function
$\psi_a(k)=\bigvee_i\lambda_i(k)a_i$ is left $Q$-linear: right
multiplication by $a_i$ commutes with the left scalar action, and finite
joins of linear functions remain linear. Moreover,
\[
 \psi_a(k)l
 =\bigvee_i\lambda_i(k)(a_i l)
 \leq\bigvee_i\lambda_i(k)l_i=k.
\]
Thus $\psi_a\in D_l$. Conversely, if $\psi\in D_l$, put
$a_i=\psi(l_i)$. Admissibility gives $a_i l\leq l_i$, and applying
$\psi$ to the frame identity gives
\[
 \psi(k)=\bigvee_i\lambda_i(k)\psi(l_i)=\psi_a(k).
\]
This proves the parametrization without interchanging any scalar
factors.

Each $K(l,l_i)$ contains zero and is closed under finite joins. Their
finite product is a nonempty directed poset, and the map
$a\mapsto\psi_a$ is monotone and preserves finite joins. For any
$\psi\in D_l$, the poset of tuples with $\psi\leq\psi_a$ is nonempty
by the displayed representation. Coordinatewise joins give upper
bounds inside that same poset. Its upper comma categories are therefore
nonempty and directed, so the parametrization is final.

It follows in particular that
\[
 C(l,k)=\mathord\downarrow\{\psi(k):\psi\in D_l\}.
\]
The right side is intrinsic. Since $D_l$ is closed under pointwise
finite joins, its value set is nonempty and finite-join closed; its
downward closure has the same properties. Admissibility gives
$\psi(k)l\leq k$, proving $C(l,k)\subseteq K(l,k)$.

Evaluation $D_l\to C(l,k)$ is final as well. For each $q\in C(l,k)$
there is a functional with $q\leq\psi(k)$, and pointwise joins give
upper bounds for any two such functionals. Consequently
\[
 \varinjlim_{\psi\in D_l}F(\psi(k))
 \simeq\varinjlim_{q\in C(l,k)}F(q),
\]
compatibly with both maps into $\varinjlim_{K(l,k)}F$.
Theorem~\ref{thm:action-colimits} now gives \eqref{eq:frame-J}.
This is an equality of images, with no injectivity assertion.

If each scalar constraint set has a greatest element, its product has
the greatest tuple $(r(l,l_i))_i$. The resulting functional dominates
every $\psi_a$, proving \eqref{eq:frame-delta}. For finite $Q$, the
finite join of all elements of $K(l,l_i)$ supplies that maximum. With
the empty frame, $L$ is zero, the parameter product is a singleton, the
only functional is zero, and $C(l,k)=\{0\}$. All formulas and finality
arguments remain valid.
\end{proof}

The complete-join hypothesis in the quantale comparison cannot be
removed by taking ideal completions and then restricting a splitting.
Here $\operatorname{Idl}$ denotes downward-closed, finite-join-closed
subsets containing zero, with joins given by ideal closure of unions.

\begin{proposition}[ideal completion does not reflect projectivity]
\label{prop:completion-projectivity-boundary}
Let
\[
 Q=\{0\}\sqcup\{t^n:n\in\Z\},\qquad
 t^m+t^n=t^{\max(m,n)},\qquad t^m t^n=t^{m+n},
\]
where $0$ is the additive zero and multiplicative absorber. Give
$L=\B$ the action in which every nonzero scalar acts as the identity.
Then $L$ is cyclic and finitely presented but is not projective over
$Q$. Nevertheless $\operatorname{Idl}(L)$ is a retract of the rank-one
free complete module over the quantale $\operatorname{Idl}(Q)$.
\end{proposition}
\begin{proof}
The presentation with one generator $x$ and relation $tx=x$ identifies
all $t^n x$, by multiplication with positive and negative powers of $t$.
It gives exactly $L$, since the support map $Q\to\B$ preserves the
relation and distinguishes zero from the common nonzero class. If that
surjection had a $Q$-linear section, the image of $x$ would be $t^m$ for
some integer $m$, and linearity would imply
$t^{m+1}=t\,s(x)=s(tx)=s(x)=t^m$, a contradiction.

The ideals of $Q$ are precisely $\{0\}$, the principal ideals
$\mathord\downarrow t^n$, and $Q$ itself: a nonempty bounded-above
downward subset of integer exponents has a greatest element. Write
$\widehat Q=\operatorname{Idl}(Q)$ and $E=Q\in\widehat Q$. Its unit is
$\mathord\downarrow t^0$, whereas $E$ is an idempotent satisfying
$IE=E$ for every nonbottom ideal $I$ and $\{0\}E=\{0\}$. Indeed, any
nonzero $t^n\in I$ satisfies $t^n Q=Q$.

The completed module $\operatorname{Idl}(L)$ is again $\B$, with each
nonbottom ideal acting as the identity. The support map
$\widehat Q\to\B$ has the section $0\mapsto\{0\}$,
$1\mapsto E$. Both maps preserve arbitrary joins, and the displayed
product identities prove their $\widehat Q$-linearity. This is the
required rank-one retraction.

Finally, $E$ is the directed join of the proper principal ideals
$\mathord\downarrow t^n$ and is not compact. The completed section
therefore sends the compact nonzero element of $\B$ to a noncompact
ideal. It does not restrict to the original coefficient semiring.
Finite presentation of $L$ does not repair this failure.
\end{proof}

There is also a structural property of the intrinsic scalar sets.
Regard $K(l,k)$ as the hom set of a category with objects $L$, identity
$1_Q$ at every object, and composition $b\circ a=ba$. Its hom sets are
join-semilattices with zero; composition preserves finite joins. An
\emph{order ideal} of this category means downward-closed,
finite-join-closed subsets of its hom sets stable under composition on
both sides. Products of order ideals use downward finite-join closure
of composites. Arbitrary joins of scalar arrows are not assumed.

\begin{proposition}[idempotence of the scalar support ideal]
\label{prop:scalar-ideal-idempotent}
For every finitely generated projective support $L$, the sets $C(l,k)$
form a two-sided order ideal and satisfy $C^2=C$. Explicitly,
\[
 C(l,k)=\mathord\downarrow\left\{
  \bigvee_{j=1}^m b_j a_j:
  a_j\in C(l,h_j),\ b_j\in C(h_j,k),\ h_j\in L\right\},
\]
where $m\geq0$. This does not imply
$\mathcal J_L^2=\mathcal J_L$.
\end{proposition}
\begin{proof}
Take $a\in C(l,h)$ and $b\in K(h,k)$. Choose $\psi\in D_l$ with
$a\leq\psi(h)$. Linearity and monotonicity give
\[
 ba\leq b\psi(h)=\psi(bh)\leq\psi(k),
\]
so $ba\in C(l,k)$. For the other side, take $a\in K(l,h)$ and
$b\in C(h,k)$, and choose $\phi\in D_h$ with $b\leq\phi(k)$.
The function $\theta(x)=\phi(x)a$ is left linear and satisfies
$\theta(x)l=\phi(x)(al)\leq\phi(x)h\leq x$.
Thus $\theta\in D_l$ and $ba\leq\theta(k)$, proving the other ideal
condition. It follows that $C^2\subseteq C$.

Conversely, for $q\in C(l,k)$ choose $\psi\in D_l$ with
$q\leq\psi(k)$. The frame expands this value as
\[
 q\leq\psi(k)=\bigvee_i\lambda_i(k)\psi(l_i).
\]
Here $\psi(l_i)\in C(l,l_i)$. Since $\lambda_i$ is admissible at
$l_i$, we also have $\lambda_i(k)\in C(l_i,k)$. Each term is an
allowed composite for $C^2$, proving the reverse inclusion. Empty joins
cover the zero support.

The coefficient groups introduce a separate issue. A scalar inequality
$q\leq\bigvee_i b_i a_i$ need not express every coefficient from $F(q)$
as a sum of products from $F(b_i)$ and $F(a_i)$. No such generation or
surjectivity is among the coefficient-diagram axioms. More concretely,
Proposition~\ref{prop:canonical-arbitrary-ideal} realizes any two-sided
ordinary ideal $I$ as $\mathcal J_L$ at singleton support, while the
scalar ideal is $C=\{0\}$. Choosing $I=2\Z$ gives
$\mathcal J_L^2=4\Z\ne2\Z=\mathcal J_L$, although $C^2=C$.
\end{proof}

\subsection{The canonical projectivity ideal}
\label{sec:canonical-ideal}
We prove Theorem~\ref{thm:canonical-ideal}. This removes the choice of a
support section in the preceding criterion. Put $\mathcal A=\mathcal A_L$.
For a $Q$-linear $\phi:L\to Q$, the groups
\[
 N_\phi(k)=F(\phi(k))
\]
form an $\mathcal A$-module by pulling back the rank-one free module.
There is no assertion that $N_\phi$ is projective. For an admissible
pair $(l,\phi)$ there is a natural map to a representable,
\[
 \theta_{l,\phi}:N_\phi\longrightarrow\mathcal A(l,-),\qquad
 x\longmapsto t_{\phi(k)l,k}a_{\phi(k),x,l}
 \quad(x\in N_\phi(k)).
\]
The transition and scalar naturality relations defining $\mathcal A$
prove its naturality.

\begin{proof}[Proof of Theorem~\ref{thm:canonical-ideal}]
Postcomposition preserves the displayed spans by naturality of
$\theta_{l,\phi}$. Precomposition by $t_{h,l}$, for $h\leq l$, uses the
same $\phi$, now admissible at $h$. The resulting arrow is
$\theta_{h,\phi}(k)(x)$.
For precomposition by $a_{q,u,h}$ with $qh=l$, put
$\psi(k)=\phi(k)q$. Right multiplication by the fixed $q$ makes $\psi$
a left $Q$-linear map even when $Q$ is noncommutative. Moreover
$\psi(k)h=\phi(k)l\leq k$. Scalar composition gives
\[
 \theta_{l,\phi}(k)(x)a_{q,u,h}
 =\theta_{h,\psi}(k)\bigl(\mu_{\phi(k),q}(x,u)\bigr).
\]
Thus precomposition also preserves the spans. Additivity and composition
of generators prove that $\mathcal J_L$ is already a two-sided ideal.
This includes $F(0)\ne0$: the zero functional is admissible and its
coefficient group is retained.

Suppose $M$ is a retract of $T^n$, with support maps
$\sigma:L\to Q^n$, $\rho:Q^n\to L$, and $\rho\sigma=1$.
Let $l_i=\rho(e_i)$ and $\phi_i=\sigma_i$. Then
\[
 k=\bigvee_i\phi_i(k)l_i,
\]
so every $(l_i,\phi_i)$ is admissible. Let $m_i\in G(l_i)$ be the image
of the $i$th free generator. On amplitudes the free retraction is the
sum of the evaluation maps
\[
 N_{\phi_i}\xrightarrow{\theta_{l_i,\phi_i}}
 \mathcal A(l_i,-)\xrightarrow{m_i}G.
\]
Its composite with the section is $1_G$, and every such evaluation has
image in $\mathcal J_LG$. Hence $\mathcal J_LG=G$. The support is a
finite free retract. Finite generation of $G$ follows from
Proposition~\ref{prop:ringoid-finite}; its ordinary projectivity follows
by lifting a pointwise amplitude epimorphism with identity support, as in
the proof of Theorem~\ref{thm:ringoid-projective}.

Conversely suppose all three conditions hold. Choose finitely many
$\mathcal A$-generators of $G$. The equality $\mathcal J_LG=G$ expresses
each as a finite sum of values of the defining arrows on elements
$m_i\in G(l_i)$. Collect the finitely many admissible pairs and elements
occurring in these expressions. Their evaluation maps give a pointwise
epimorphism
\[
 c:\bigoplus_{i=1}^nN_{\phi_i}\longrightarrow G.
\]
Indeed its image is a submodule containing all the chosen generators.
Projectivity of $G$ supplies an $\mathcal A$-linear section $s$.
The corresponding support maps satisfy only
\[
 \rho_1\phi(k)=\bigvee_i\phi_i(k)l_i\leq k.
\]
Choose independently a finite free support retraction
$\sigma_0:L\to Q^m$, $\rho_0:Q^m\to L$. Appending its coordinates gives
\[
 \sigma'=(\phi,\sigma_0),\qquad
 \rho'(u,v)=\rho_1(u)\vee\rho_0(v),\qquad
 \rho'\sigma'=1_L.
\]
Map the first $n$ free generators of $T^{n+m}$ to the $m_i$, and the
remaining generators to the group zeros over $\rho_0(e_j)$.
These zeros have the required supports although their amplitudes vanish.
The resulting $T$-map has support $\rho'$. Give its section support
$\sigma'$ and amplitude $(s,0)$. The fixed-support equivalence shows
that it is $T$-linear. Its composite has identity support and amplitude
$cs=1_G$. Hence $M$ is a finite free retract.

Finally, extension along the quotient category is the pointwise quotient
$G/\mathcal J_LG$: its universal property is precisely that maps annihilate
the action of $\mathcal J_L$. This proves the tensor formulation.
The fixed-support equivalence gives the claimed classification there.
A $Q$-isomorphism of supports transports admissible functionals by
precomposition and transports scalar and transition arrows, hence the
ideal and all the stated data.
\end{proof}

\begin{corollary}[idempotents in the canonical ideal]
\label{cor:canonical-ideal-matrices}
The amplitudes in Theorem~\ref{thm:canonical-ideal} are exactly the images
of idempotent endomorphisms $E$ of finite sums
$P=\bigoplus_i\mathcal A_L(l_i,-)$ whose matrix entries lie in
$\mathcal J_L$. Invariantly, the condition is $E(P)\subseteq\mathcal J_LP$.
The usual rectangular intertwiners classify these images up to isomorphism.
\end{corollary}
\begin{proof}
At an object $k$, the submodule $\mathcal J_LP$ is
$\bigoplus_i\mathcal J_L(l_i,k)$, by the ideal property and the identity
generators of the representables. If $G=\im E$ and $\mathcal J_LG=G$,
its inclusion into $P$ gives $E(P)\subseteq\mathcal J_LP$, exactly the
entry condition. Conversely, if $g\in G\subseteq\mathcal J_LP$, write
$g=\sum_a j_ap_a$ and apply $E$. Then
$g=\sum_a j_aE(p_a)\in\mathcal J_LG$.
Every finite projective has such a representable presentation.
Extending an isomorphism of images by inclusion and projection gives
the rectangular intertwiners, and restricting those identities recovers
the isomorphism.
\end{proof}
The matrix argument is the ordinary trace-ideal argument; compare
Herbera and P\v r\'ihoda \cite[p.~2 of the preprint]{HerberaPrihoda}.
The additional datum here is the explicitly defined ideal
$\mathcal J_L$. For the top-born diagram in
Proposition~\ref{prop:ringoid-obstruction}, admissibility at the top
forces $\phi(a)=\phi(b)=0$, hence $\phi(a\vee b)=0$; therefore
$\mathcal J_LG=0\ne G$. The canonical ideal detects the same obstruction.

There is no claim that these categories are left rooted. The projective
functor classification of Cort\'es-Izurdiaga and Torrecillas
\cite[Theorem~1, pp.~320--322]{CortesTorrecillas} assumes that hypothesis;
their Remark~3 and Example~1 describe the cycle issue. The action
categories here can have scalar loops, so that result does not replace
ordinary projectivity in Theorem~\ref{thm:canonical-ideal}. Nor does the
theorem give a finite procedure for listing all admissible functionals.

\begin{proposition}[the canonical ideal need not be idempotent]
\label{prop:canonical-nonidempotent}
There is a finite distributive, commutative, entire idempotent semiring
$Q$ and a finite free $Q$-semimodule $L$ for which the canonical ideal
of $G_Q(\Z)$ satisfies $\mathcal J_L^2\ne\mathcal J_L$.
\end{proposition}
\begin{proof}
Let $Q=\{0,a,b,c,1\}$, ordered by
$0<a<b<1$ and $0<a<c<1$, with $b,c$ incomparable.
This is the ideal lattice of a three-point poset with one least point
and two incomparable larger points. Addition is join and multiplication
is meet. Thus $b\vee c=1$, $b\wedge c=a$, and every two nonzero
members have nonzero meet. In $T=G_Q(\Z)$ only the top coefficient
group is nonzero, and it is $\Z$.

Take $L=Q^2$, $o=(0,0)$ and $k=(b,c)$. Here $\mathcal A_L$ is exactly
the $\Z$-linear incidence category of $L$: every proper-support scalar
arrow is zero, the top scalar $n$ is $n1_l$, and the remaining relations
are those of the transition arrows. Conversely these assignments in the
incidence category satisfy every defining relation, so no further
relations are imposed.

Every $Q$-linear functional has the form
\[
 \phi(u,v)=(u\wedge d_1)\vee(v\wedge d_2).
\]
A nonzero canonical generator at target $k$ requires $\phi(k)=1$.
Both $d_1,d_2$ must then be nonzero, because its individual summands
are bounded by $b<1$ and $c<1$. If $\phi$ is admissible at
$l=(l_1,l_2)$, evaluation at $(1,0)$ and $(0,1)$ gives
$d_1\wedge l_2=d_2\wedge l_1=0$. Entirety forces $l=o$.
Consequently $\mathcal J_L(l,k)=0$ for $l\ne o$.
At $o$, the functional $\phi(u,v)=u\vee v$ is admissible and has
$\phi(k)=1$, so its coefficient $1$ gives
$\mathcal J_L(o,k)=\Z t_{o,k}$. On the other hand
$\mathcal J_L(o,o)=0$, because every functional preserves zero and
$F(0)=0$. In every composite contributing to
$\mathcal J_L^2(o,k)$, the intermediate object must therefore be $o$,
and the remaining factor is zero. Hence
\[
 \mathcal J_L(o,k)=\Z,\qquad \mathcal J_L^2(o,k)=0.
\]
This proves the assertion. It does not contradict the criterion:
$\mathcal J_LG=G$ implies $\mathcal J_L^nG=G$ for every $n\geq1$,
without requiring equality of the ideals themselves.
\end{proof}

\begin{proposition}[every ordinary ideal occurs]
\label{prop:canonical-arbitrary-ideal}
Let $R$ be any unital ring and $I$ any two-sided ideal, without a
unitality assumption on $I$. There is a two-idempotent inverse semiring
$T$ and a finite projective support $L$ for which
$\mathcal A_L=R$ and $\mathcal J_L=I$.
A semimodule on this support is finite projective over $T$ exactly when
its underlying $R$-module $P$ is finite projective and $IP=P$.
In particular, powers of the canonical ideal need not stabilize.
\end{proposition}
\begin{proof}
Take the disjoint union of a lower copy of $I$ and an upper copy of $R$.
Add within each group in the usual way; a mixed sum lies in the upper
copy and uses $I\hookrightarrow R$. A product lies in the lower copy
if either factor does, and otherwise in the upper copy, with its
coefficient given by ring multiplication. The two-sided ideal property
makes this well defined. Associativity and distributivity hold on
coefficients by the ring identities and on tags by the meet and join
laws of $\B$. The lower zero is the absorbing additive identity, the
upper $1_R$ is the multiplicative identity, and additive inverse is
negation within each group. Thus this is the injective two-idempotent
construction with $F(0)=I$, $F(1)=R$ and transition the inclusion.

Let $L=\{o\}$ be the zero $\B$-semimodule, which is finite projective.
Its action category has one object. The upper scalar arrows give $R$;
the coefficient transition relation identifies each lower scalar
$i\in I$ with that same element of $R$. Conversely the ordinary
$R$-action satisfies all defining relations. Hence there are no further
relations and $\mathcal A_L$ is the one-object category $R$.
The only $\B$-linear functional $L\to\B$ is zero. It is admissible,
and its canonical generators are precisely $F(0)=I$. Therefore
$\mathcal J_L=I$. The projectivity assertion is now
Theorem~\ref{thm:canonical-ideal}. Taking $R=\Z$ and $I=2\Z$ gives
$\mathcal J_L^n=2^n\Z$, which is strictly decreasing for every $n$.
\end{proof}
These examples separate the canonical support ideal from the trace
ideal of a particular finite projective. They also show that arbitrary
ordinary ideal-generated projective modules occur in the residual
classification problem.

The generation condition can also be restricted to an idempotent ideal.
The largest-idempotent-subideal construction is ordinary ideal theory;
for a unital ring, see Emmanouil \cite[Lemma~1.3(ii)]{Emmanouil}.
We give the small-category argument to specify its application here.

\begin{proposition}[the idempotent part of the canonical ideal]
\label{prop:canonical-ideal-core}
For a small preadditive category $\mathcal A$ and a two-sided ideal
$J$, there is a largest idempotent two-sided ideal $C(J)\subseteq J$.
It is the stable value of the transfinite sequence
\[
 J_0=J,\qquad J_{\alpha+1}=J_\alpha^2,\qquad
 J_\lambda=\bigcap_{\alpha<\lambda}J_\alpha
 \quad(\lambda\text{ a limit ordinal}).
\]
For every finitely generated projective covariant $\mathcal A$-module
$P$,
\[
 JP=P\quad\Longleftrightarrow\quad C(J)P=P.
\]
Every finite idempotent matrix with typed entries in $J$ already has
all its entries in $C(J)$. Thus Theorem~\ref{thm:canonical-ideal}
remains unchanged if $\mathcal J_L$ is replaced by
$C(\mathcal J_L)$.
\end{proposition}

\begin{proof}
Products of ideals consist of finite sums of typed composites;
arbitrary intersections and sums are ideals. Smallness ensures that
the family of ideals is a set. Let $C$ be the sum of all idempotent
ideals $I\subseteq J$. Each $I=I^2$ lies in $C^2$, so $C\subseteq
C^2$. The reverse inclusion follows because $C$ is an ideal.
Therefore $C$ is the largest idempotent ideal contained in $J$.

The displayed transfinite sequence is descending. Every idempotent
$I\subseteq J$ lies in each term: use $I=I^2$ at successors and
intersection at limits. The sequence stabilizes because the total
typed arrow collection of $\mathcal A$ is a set. Otherwise each
successor before its successor cardinal would remove a distinct
arrow, which is impossible. At a stable stage the ideal is idempotent,
so its maximality property identifies it with $C$. No finite-stage
stabilization has been assumed.

Let $E$ be a finite idempotent matrix with entries in $J$, and let
$I_E$ be the two-sided ideal generated by its entries. The equations
$E_{ij}=\sum_k E_{ik}E_{kj}$ put every generator in $I_E^2$.
Since $I_E^2$ is an ideal, $I_E\subseteq I_E^2$; hence
$I_E^2=I_E$. Maximality gives $I_E\subseteq C$, proving the
entrywise assertion.

For completeness, choose a finite representable sum $N$ and split
maps $P\xrightarrow{i}N\xrightarrow{r}P$ with $ri=1_P$.
If $JP=P$, then $i(P)=J i(P)\subseteq JN$, so all entries of the
idempotent $E=ir$ lie in $J$. They therefore lie in $C$, giving
$E(N)\subseteq CN$. Applying $r$ gives $P=rE(N)\subseteq CP$,
and consequently $CP=P$. The converse follows from $C\subseteq J$.
Apply this equivalence to the action category and its canonical ideal.
\end{proof}

This is square descent; it is not the iteration
$J_{\alpha+1}=J J_\alpha$ with a fixed left factor, whose stable
value need not be idempotent. The distinction is explicit in
\cite[Lemmas~1.1 and~1.3]{Emmanouil}. The already realized case
$J=2\mathbb Z$ has $C(J)=0$: its finite square-descending terms are
$2^{2^n}\mathbb Z$, with zero intersection at $\omega$. Thus the
reduction can be strict and need not stabilize at a finite stage.
If $C(\mathcal J_L)=0$, every allowed projective amplitude vanishes,
while its support $L$ is retained. No local-unit property, decision
procedure or classification of projectives over arbitrary idempotent
ideals is asserted.

\subsection{Exact finite coefficient models}
\begin{proof}[Proof of Theorem~\ref{thm:finite-coefficient-models}]
By Corollary~\ref{cor:canonical-ideal-matrices}, present each
$G_s$ by an idempotent with entries in $\mathcal J$. Combine their
finite object lists and pad by zeros, giving idempotents $E_s$ in
one full action matrix ring $B$. For each source $j$, choose
admissible functionals representing all entries in column $j$ of
all these projectors. Their finite pointwise join $\phi_j$ is
still admissible, and all the coefficients transport to
$q_{ij}=\phi_j(l_i)$. An empty column uses the zero functional.

Admissibility and $Q$-linearity give, in the stated order,
\begin{equation}\label{eq:finite-coefficient-supports}
 q_{ij}q_{jk}
 =\phi_k\bigl(\phi_j(l_i)l_j\bigr)\leq q_{ik}.
\end{equation}
If $Q_0=(q_{ij})$, then $Q_0^2\leq Q_0$ and
$D=I\vee Q_0$ satisfies
\[
 D^2=D,\qquad DQ_0=Q_0=Q_0D.
\]
Its diagonal entries are upper idempotents. The ordinary fiber
$C=A_D(T)$ is unital, with diagonal corners $F(d_{ii})$; the
images of the global unit in those corners are its diagonal
identity entries. Since $d_{ij}l_j\leq l_i$, coefficient-to-action
maps give the asserted unital $\pi:C\to B$.

The two displayed support products show that the subgroup $N$
is an ideal. In each product with an $N$ entry, choose its
preimage in $F(q_{jk})$ or $F(q_{ij})$ and transport the product
through $F(q_{ik})$ to the target group of $C$. Naturality
identifies this with the original product, irrespective of
injectivity of the diagonal transitions. Its image is in the
canonical ideal by admissibility of the source functional.

Put $K=\ker\pi$. If an entry $x\in F(d_{ij})$ of an element of
$K$ vanishes in $\mathcal A(l_j,l_i)$, the directed-colimit zero
criterion supplies $r\geq d_{ij}$ with $rl_j\leq l_i$ and
$t_{d_{ij},r}(x)=0$. For a representative
$y\in F(q_{jk})$ of an $N$ entry, one has
\[
 rq_{jk}=\phi_k(rl_j)\leq\phi_k(l_i)=q_{ik}.
\]
The product $xy$ can therefore first transport $x$ to $F(r)$,
where it is zero, before transporting through $F(q_{ik})$ to
$F(d_{ik})$. Every summand of the matrix product vanishes, so
$KN=0$. A symmetric annihilator assertion is not needed.

By the choice of the functionals, $E_s\in\pi(N)$. Choose a lift
$a_s\in N$. Its defect $a_s^2-a_s$ lies in $K\cap N$ and
annihilates $a_s$ on the right. Thus $a_s^3=a_s^2$, and
$e_s=a_s^2$ is an idempotent in $N$ with image $E_s$. This
exact lift needs no semiperfectness or general lifting theorem.

For any $X,Y\in\pi(N)$ and $b\in B$, an entry of $b_{jk}$
has a representative at $r$ with $rl_k\leq l_j$. The supports
of its sandwich terms satisfy
\[
 rq_{km}=\phi_m(rl_k)\leq q_{jm},\qquad
 q_{ij}rq_{km}\leq q_{ij}q_{jm}\leq q_{im}.
\]
Thus $XbY\in\pi(N)$, even for arbitrary full action arrows
$b$. If $e,f\in N$ are idempotents, choose a preimage $c\in C$
of any element of $\pi(e)B\pi(f)$ and replace it by $ecf$.
This proves surjectivity of the corner map. If $z\in eCf$ is
in $K$, then $z=zf=0$ by $KN=0$, proving injectivity.
Composition and the corner identities are preserved by $\pi$.

Idempotent matrix calculus on finite sums of representables now
identifies these mixed corners with all amplitude morphisms,
with the customary opposite convention. Conversely an idempotent
in $N$ has canonical-ideal entries after applying $\pi$, so its
amplitude is finite projective and $\mathcal J$-generated. Zero
projectors handle zero amplitudes; an empty family may use the
zero matrix ring. None of these choices collapses the specified
support. This proves all assertions.
\end{proof}

\begin{example}[no fixed finite set of action objects]
\label{ex:no-fixed-finite-objects}
The finite coefficient model in Theorem~\ref{thm:finite-coefficient-models}
cannot in general be obtained by restricting to one fixed finite
set of action objects, even for one free support and semisimple
upper coefficient rings.

Let $k$ be a nonzero field and let
$M=\langle a,b\mid ba=1\rangle$ be the bicyclic monoid. Put
$e_n=a^nb^n$ for $n\geq1$, and let $Q$ be the semiring of finite
subsets of $M$, with union and setwise multiplication. Form the ring
\[
 R=k1\oplus\bigoplus_{n\geq1}kz_n,
 \qquad z_n^2=z_n,\quad z_nz_m=0\quad(n\ne m).
\]
Give it the $M$-grading $R_1=k1$, $R_{e_n}=kz_n$, with all
other components zero. Set $F(S)=\bigoplus_{m\in S}R_m$, with
inclusion transitions and the restricted multiplication of $R$,
and reconstruct $T=(Q,F)$. Every upper multiplicative idempotent
$E$ of $Q$ has $F(E)$ a finite product of copies of $k$.

Nevertheless, on the single free support $L=Q$ there are pairwise
nonisomorphic finitely generated projective $T$-semimodules with
amplitudes
\[
 G_n(S)=
 \begin{cases}k,&a^n\in S,\\0,&a^n\notin S,\end{cases}
 \qquad n\geq1,
\]
whose nonzero transitions are identities and whose amplitude
endomorphism rings are $k$. For any finite collection of objects
$S_1,\ldots,S_r\in L$, some $G_n$ restricts to the zero functor
on their full action subcategory, just as the zero-amplitude
projective on $L$ does.
\end{example}
\begin{proof}
Every element of $M$ has a normal form $a^ib^j$. The action on
$\mathbb N$ given by $a(t)=t+1$ and $b(t)=\max(t-1,0)$ sends
this form to $i+\max(t-j,0)$ and distinguishes the forms.
Consequently the degree $i-j$ is additive, and
$e_ne_m=e_{\max(n,m)}$. The displayed grading is multiplicative:
different nonidentity homogeneous components have zero product.
This verifies the coefficient system and hence the reconstruction.

An upper idempotent $E$ is a finite submonoid containing $1$.
An element of nonzero integer degree would have infinitely many
distinct powers. Thus $E=\{1\}\cup\{e_n:n\in A\}$ for a
finite set $A$ of positive integers. Conversely every such set
is a submonoid. The ring isomorphism
\[
 F(E)\longrightarrow k^{1+|A|},\qquad
 c_0 1+\sum_{n\in A}c_nz_n\longmapsto
       \bigl(c_0,(c_0+c_n)_{n\in A}\bigr)
\]
proves the assertion about upper coefficient rings.

The element $t_n=(\{e_n\},z_n)$ is multiplicatively idempotent.
Right multiplication by $t_n$ is an idempotent endomorphism of
the regular left $T$-semimodule, so $P_n=Tt_n$ is finitely
generated projective. Its support is $Q\{e_n\}$, identified
with $Q$ by the mutually inverse $Q$-linear maps
\[
 S\longmapsto S\{b^n\},\qquad
 U\longmapsto U\{a^n\}.
\]
Here $U\{e_n\}=U$, and the inverse identities follow from
$b^na^n=1$ and $a^nb^n=e_n$.

The fixed coefficients of right multiplication by $t_n$ at
such a support $U$ are $\{x\in F(U):xz_n=x\}$. Linear
independence of $1,z_1,z_2,\ldots$ makes this group $kz_n$
if $e_n\in U$, and zero otherwise. Normal forms show that
$e_n\in S\{b^n\}$ exactly when $a^n\in S$. This proves
the asserted amplitude, including its action and transition
maps inherited from the retract. Every amplitude endomorphism
is $k$-linear. The identity transitions from $\{a^n\}$ to
every nonzero fiber force a common scalar, and every scalar
commutes with the action. Thus its endomorphism ring is $k$.

There is no nonidentity $Q$-linear automorphism of the regular
support $Q$. Such an automorphism is right multiplication by
a finite set $U$ with a finite inverse $V$, so
$UV=VU=\{1\}$. For every $u\in U$ and $v\in V$ one has
$uv=vu=1$. Normal forms imply that $M$ has only the unit $1$:
if $a^ib^j a^rb^s=1$, then $i=s=0$ and $j=r$, and the
reverse product is $1$ only when $j=0$. Hence
$U=V=\{1\}$. The fibers at $\{a^n\}$ now distinguish
the $P_n$ even without a chosen support identification.

Right multiplication by $(\{1\},0)$ is the zero-amplitude
projector with support $Q$. Finally, a finite union
$S_1\cup\cdots\cup S_r$ omits some $a^n$. All values of
$G_n$ on these objects vanish, and therefore so do all its
restricted arrows. It remains nonzero at $\{a^n\}$. This
proves the restriction obstruction. It makes no assertion
against every possible universal ordinary-ring construction.
\end{proof}


\subsection{A uniform category of coefficient groups}
Theorem~\ref{thm:uniform-coefficient-category} replaces the
family-dependent coefficient choice by a single category,
retaining the ordinary ideal-generated projective problem.
We use the standard covariant module convention for a
preadditive category; representables, finite projectors and
extension along an additive functor have their usual meanings
\cite{Mitchell}. We use the covariant convention and representable
projectivity described in \cite[\S1, Proposition~3, pp.~1--4]{Murfet};
extension and its adjunctions are given in
\cite[\S4, Theorem~18 and Proposition~19, pp.~23--25]{Murfet}.
The coefficient category and its ideal below specify the form
needed for inverse-semiring amplitudes.

\begin{proof}[Proof of Theorem~\ref{thm:uniform-coefficient-category}]
Admissibility and $Q$-linearity give, in the required order,
\[
 q_{zy}q_{yx}
 =\phi_x\bigl(\phi_y(l_z)l_y\bigr)
 \leq\phi_x(l_z)=q_{zx}.
\]
Hence $q_{xx}^2\leq q_{xx}$ and $1\vee q_{xx}$ is an upper
multiplicative idempotent. Accounting for a diagonal identity
term also gives
\[
 d_{zy}d_{yx}\leq d_{zx},\qquad
 d_{zy}q_{yx}\leq q_{zx},\qquad
 q_{zy}d_{yx}\leq q_{zx}.
\]
For instance the extra term in the first factor occurs only
when $z=y$, and is then $q_{yx}=q_{zx}$.
These inequalities define composition in $\mathcal C_L$.
The coefficient identities, associativity and naturality give
its unital category axioms. Every $\mathcal N_L$-arrow has a
representative at its $q$-support, so the two mixed inequalities
prove both ideal closures. The diagonal quotient follows.
The set of objects is a subset of $L\times Q^L$, hence is small.

Since $d_{yx}l_x\leq l_y$, every coefficient defines an action
arrow, giving the unital functor $\pi$. Its image on
$\mathcal N_L$ lies in $\mathcal J_L$ by admissibility.
Write $K=\ker\pi$. Then
\begin{equation}\label{eq:uniform-annihilator}
 K\mathcal N_L=0.
\end{equation}
Indeed, if $a\in K(y,z)$, the directed-colimit zero criterion
for $\mathcal A_L(y,z)$ supplies $r\geq d_{zy}$ with
$rl_y\leq l_z$ at which $a$ becomes zero. For an
$\mathcal N_L(x,y)$-arrow represented at $q_{yx}$,
\[
 rq_{yx}=\phi_x(rl_y)\leq\phi_x(l_z)=q_{zx}.
\]
Compute the product by first transporting $a$ to that zero
and then transporting the product to $d_{zx}$. This proves
\eqref{eq:uniform-annihilator}. No reverse annihilator,
injectivity of transitions, or vanishing of $F(0)$ is used.

For finite lists of objects, let $e,f$ be idempotent matrices
with entries in $\mathcal N_L$. Matrix calculus uses the common
opposite convention for covariant representables on both sides.
All mixed action corners
lift exactly to coefficient corners. To see surjectivity,
represent an intervening action arrow $b_{jk}$ at $r$ with
$rl_k\leq l_j$. The two outside entries have representatives
at $q_{ij}$ and $q_{km}$, and
\[
 rq_{km}=\phi_m(rl_k)\leq q_{jm},\qquad
 q_{ij}rq_{km}\leq q_{ij}q_{jm}\leq q_{im}.
\]
Thus each term of $\pi(e)b\pi(f)$ has a representative in
$F(q_{im})$, hence in $\mathcal N_L$. Choose a coefficient
matrix lift $c$ and replace it by $ecf$. If a corner element
$z\in e\operatorname{Mat}(\mathcal C_L)f$ maps to zero,
then $z=zf=0$ by \eqref{eq:uniform-annihilator}.
Consequently
\[
 e\operatorname{Mat}(\mathcal C_L)f
 \xrightarrow{\sim}
 \pi(e)\operatorname{Mat}(\mathcal A_L)\pi(f)
\]
is an additive isomorphism compatible with composition.
Object repetitions in these finite lists are allowed.
A matrix entry between two copies of the same categorical
object belongs to $F(1\vee q_{xx})$, and its ideal entry is
still represented at $q_{xx}$; the calculations remain valid.

A finite projective $P$ over a preadditive category satisfies
$\mathcal NP=P$ exactly when it has an idempotent presentation
with entries in $\mathcal N$. For a split presentation
$P\to V\to P$ by finite representables, the section sends
$\mathcal NP$ into $\mathcal NV$, proving necessity.
Conversely an $\mathcal N$-entry projector $e$ satisfies
$eV=e^2V\subseteq e(\mathcal NV)=\mathcal N(eV)$.
Left Kan extension sends the representable at $x$ to the
representable at $l_x$, and therefore sends finite projectors
to their images under $\pi$. The exact corner calculation
proves full faithfulness on the stated subcategory.

For essential surjectivity, present a finite projective
$\mathcal J_L$-generated amplitude by an idempotent $E$ with
entries in $\mathcal J_L$. In each source column $j$, choose
admissible functionals representing its finitely many entries
and replace them by their finite pointwise join $\phi_j$.
Every entry then transports to $F(\phi_j(l_i))$.
Choose the objects $(l_j,\phi_j)$ of $\mathcal C_L$.
The matrix $E$ lifts to a matrix $a$ with entries in
$\mathcal N_L$, even when some of these objects coincide.
Its defect belongs to $K$, so
\[
 (a^2-a)a=0,\qquad a^3=a^2,\qquad (a^2)^2=a^2.
\]
Thus $a^2$ is an exact coefficient projector lifting $E$.
This proves essential surjectivity and the equivalence.
Extension to the quotient category is $P\mapsto
P/\mathcal N_LP$, proving the asserted kernel description.

Finally, $u\in\operatorname{Aut}_Q(L)$ sends an admissible
pair to $(u(l),\phi u^{-1})$, since
\[
 \phi(u^{-1}(k))u(l)=u\bigl(\phi(u^{-1}(k))l\bigr)\leq k.
\]
All $q$-values are unchanged, so this is a relabeling of the
same coefficient groups and products. It commutes with $\pi$
and the action-category relabeling. The canonical projectivity
criterion now gives the fixed-support classification and the
automorphism-orbit statement.
\end{proof}

The coefficient category may still be infinite. In the
singleton-support arbitrary-ideal example of
Proposition~\ref{prop:canonical-arbitrary-ideal}, it is the
single ring $R$ and its ideal is the arbitrary ordinary ideal
$I$. Thus the theorem removes the action-arrow colimits and
family-dependent coefficient choices, while retaining the
ordinary ideal-generated projective problem. It makes no
finite-object assertion and does not contradict
Example~\ref{ex:no-fixed-finite-objects}.

\subsection{A coefficient quotient for arbitrary support}
\label{subsec:upper-unit-quotient}

For distributive lattices of rings, transition images and the
annihilator control of transition kernels already occur in
Bandelt and Petrich \cite[Lemma~4, pp.~159--160]{BandeltPetrich}.
The following argument permits arbitrary support multiplication and
uses a bound inside each additive matrix fiber. The final lifting step
is ordinary nilpotent idempotent lifting
\cite[Lemma~10.32.7, Tag~05BU]{Stacks}.

\begin{proof}[Proof of Theorem~\ref{thm:upper-unit-quotient}]
Each test in the definition of $K_q$ is an additive homomorphism of
$x$, so $K_q$ is a subgroup. If $q\leq q'$, naturality sends its test
to the old test followed by the transition
$pqr\vee1\leq pq'r\vee1$. Thus transitions preserve $K$.
For $a\in F(s)$ and $b\in F(t)$, an arbitrary context around $axb$
is
\[
 u(axb)v=(ua)x(bv),\qquad
 ua\in F(ps),\quad bv\in F(tr).
\]
Both tests have output support $psqtr\vee1$. Hence $axb\in K_{sqt}$.
This proves the two-sided ideal property with the noncommutative
order preserved. The quotient groups inherit all products,
transitions and the unit, giving the asserted inverse semiring $T'$.

If $s\geq1$, use the unit coefficient on both sides of $x\in K_s$.
The required transition is the identity of $F(s)$, so $x=0$.
Conversely, let $H$ be a coefficient ideal vanishing at all such
supports. For $x\in H_q$, every $uxv$ belongs to $H_{pqr}$ and
transports into $H_{pqr\vee1}=0$. Therefore $H\subseteq K$.
This proves maximality. Any ideal of $F/K$ vanishing above $1$ has
an inverse image containing $K$ and vanishing there, since those
groups were unchanged. Maximality forces that inverse image to equal
$K$. There is consequently no residual ideal of the stated kind.

Fix $E\in M_n(Q)$ with $E^2=E$. Its additive fiber
\[
 A_E=\prod_{i,j}F(E_{ij})
\]
is a possibly nonunital ring: each coefficient product is transported
along $E_{ik}E_{kj}\leq E_{ij}$ and the resulting contributions are
added in the target group. The quotient ring map has kernel
$N_E=\prod_{i,j}K_{E_{ij}}$.

Put $e=E_{vv}$ and $b=E_{vw}$. Matrix idempotence gives $eb\leq b$,
and hence
\[
 (e\vee1)b=eb\vee b=b.
\]
For $z\in K_e$, the unit context makes $t_{e,e\vee1}(z)=0$.
Multiplication naturality therefore gives, for every $y\in F(b)$,
\begin{equation}\label{eq:upper-unit-cycle}
 t_{eb,b}(zy)=0.
\end{equation}
This uses the displayed matrix relation, not an assumption $e\leq1$.

Expand a product of $n+1$ matrices from $N_E$ into path
contributions. A path has vertices $i_0,\ldots,i_{n+1}$.
Among its first $n+1$ vertices two coincide, giving a nonempty cycle
with a following edge $v\to w$. The cycle coefficient belongs to $K$
at its scalar product support. Transport it to $F(E_{vv})$; this is
allowed because every positive power of $E$ is $E$, and it remains
in $K_{E_{vv}}$. Equation~\eqref{eq:upper-unit-cycle} annihilates the
cycle followed by the next edge after transport to $F(E_{vw})$.
Replacing this segment by the single $v\to w$ edge leaves a
nonempty valid path with the original endpoints. All its scalar
supports remain bounded by the final entry of $E$. Naturality
identifies the original contribution with the resulting zero
contribution. Thus $N_E^{n+1}=0$. Only the matrix index set was
required to be finite.

To lift an idempotent over $T'$, lift its entries at the same support
$E$ to $a\in A_E$. Its defect $d=a^2-a$ belongs to $N_E$.
Replace $a$ by
\[
 a'=3a^2-2a^3=a+d-2ad.
\]
The new defect is $d^2(4d-3)$, so iteration reaches an idempotent
without changing the quotient. The correction has no constant term
and remains in the possibly nonunital ring $A_E$.

For uniqueness, let lifted idempotents $P,Q$ have isomorphic quotient
images. Lift inverse rectangular maps at their exact supports, then
replace them by $X=PXQ$ and $Y=QYP$. The support equations are
unchanged. We have $XY=P+\delta$ and $YX=Q+\epsilon$, with the
defects in the respective nilpotent coefficient-fiber ideals. In
the corner rings with identities $P,Q$, both products are invertible
by finite geometric series. The identity $(YX)Y=Y(XY)$ gives
\[
 Y'=Y(XY)^{-1}=(YX)^{-1}Y,\qquad XY'=P,\quad Y'X=Q.
\]
These are the required lifted inverse maps. Coefficient zero at a
nonzero support is retained throughout; it is not the zero semiring
matrix.

Extension of scalars sends a split idempotent image to its quotient
idempotent image. The existence and uniqueness just proved therefore
give its bijection on finite-projective classes. Block sums are
preserved, yielding the stated monoid and group isomorphisms.
\end{proof}

The largest-ideal characterization makes the construction invariant
under pair isomorphisms. It does not make it functorial under an
arbitrary map to another pair: the target may contain coefficient
contexts unavailable in the source. The theorem reduces the general
coefficient diagram; it does not classify the residual projective
supports or ordinary amplitudes.

\subsection{Upper idempotents detect the necessary coefficients}
The two-sided context construction admits the stronger testing family
in Theorem~\ref{thm:upper-idempotent-quotient}. The lifting step remains
the ordinary nilpotent correction already proved above and attributed
to \cite[Lemma~10.32.7, Tag~05BU]{Stacks}.

\begin{proof}[Proof of Theorem~\ref{thm:upper-idempotent-quotient}]
For $e\in U(Q)$, the coefficient product lands in $F(e^2)=F(e)$.
Associativity and biadditivity make this a ring. Naturality applied to
$1\leq e$ shows that the transitioned unit acts as a two-sided
identity.

Each defining test for $K^{\mathrm{id}}_q$ is an additive homomorphism.
If $q\leq q'$, a test at $q'$ with target $e$ is the same test at $q$
followed by transport, since $pqr\leq pq'r\leq e$. Thus transitions
preserve the subgroups. Multiplication on either side preserves them
because a test of $axb$ can be regrouped as
\[
 u(axb)v=(ua)x(bv),
\]
with the same ordered support bound. The groups therefore form a
coefficient ideal, without a commutativity assumption. The unit test
at $e\in U(Q)$ gives $K^{\mathrm{id}}_e=0$. Conversely, every
coefficient ideal vanishing at these supports sends each displayed
context into its zero group at $e$, so it is contained in
$K^{\mathrm{id}}$. This proves maximality and the existence of the
quotient. The inverse image of a quotient ideal vanishing at every
$e\in U(Q)$ has the same property upstairs, since these fibers were
unchanged. Maximality makes that quotient ideal zero.

Fix $E^2=E$. For a vertex $v$, put $a=E_{vv}$. The diagonal
idempotence equation gives $a^2\leq a$, and consequently
\[
 e_v=a\vee1\in U(Q),\qquad
 e_v^2=a^2\vee a\vee1=a\vee1.
\]
For each edge support $b=E_{vw}$, the same matrix equation gives
\[
 e_vb=(E_{vv}\vee1)E_{vw}=E_{vw}=b.
\]
If $z\in K^{\mathrm{id}}_{E_{vv}}$, its transition to $F(e_v)$ is
zero. Naturality therefore shows that $zy$, transported to $F(E_{vw})$,
is zero for every $y\in F(E_{vw})$.

The additive matrix fiber $A_E$ is the possibly nonunital ring with
entries $F(E_{ij})$ and the prescribed matrix transports. Its kernel
$N_E$ has entries $K^{\mathrm{id}}_{E_{ij}}$. In a product of $n+1$
elements of $N_E$, expand an entry into contributions of paths
$i_0,\ldots,i_{n+1}$. Among the first $n+1$ vertices a vertex repeats,
so there is a nonempty cycle with a following edge. Transport the
cycle coefficient to $F(E_{vv})$, using $E^m=E$ for every positive
$m$. It remains in $K^{\mathrm{id}}_{E_{vv}}$. The preceding
annihilation identity kills that cycle and its following edge after
transport to the edge support. Multiplying the path prefix and suffix
and using naturality kills the original contribution. Thus
$N_E^{n+1}=0$.

For completeness, lift a quotient idempotent at its exact support
$E$ to $a\in A_E$. Its defect $d=a^2-a$ lies in $N_E$. Replacing
$a$ by $3a^2-2a^3$ changes the defect to $d^2(4d-3)$, so finitely
many repetitions produce an idempotent. The polynomial has no constant
term, hence is valid in $A_E$ even when this ring is nonunital.
To lift an isomorphism, take rectangular inverse intertwiners in the
quotient, lift their coefficients at the same supports, and cut them
by the two lifted idempotents $P,Q$. The resulting $X=PXQ$ and
$Y=QYP$ have composites $XY=P+\delta$, $YX=Q+\epsilon$, where the
errors lie in the respective nilpotent fiber ideals. Finite geometric
series invert both composites in their corners. Then
\[
 Y'=Y(XY)^{-1}=(YX)^{-1}Y
\]
satisfies $XY'=P$ and $Y'X=Q$. Existence and uniqueness of lifted
projective classes follow. Extension of scalars preserves the split
idempotent images and their direct sums, proving the assertions on
$V$ and $K_0$. All operations preserve exact support; zero amplitude
at a nonzero support has not become the global zero object.

The old upper-unit ideal vanishes at every upper idempotent, so
maximality includes it in $K^{\mathrm{id}}$. If $1$ is greatest,
$U(Q)=\{1\}$. The unit test and multiplication naturality give
$K^{\mathrm{id}}_q=\ker(F(q)\to F(1))$. As before, pair isomorphisms
transport all tests, but arbitrary-map functoriality has not been
asserted.
\end{proof}

\begin{example}[strictness and the scope of nilpotence]
\label{ex:upper-idempotent-strict}
Let $Q=\{\bot\}\sqcup\mathbb N$ have addition $\max$ and
multiplication given by ordinary addition on $\mathbb N$, with
absorbing zero $\bot$ and multiplicative unit $0\in\mathbb N$.
For a nonzero unital ring $R$, put $F(\bot)=0$ and $F(n)=R$ for
all $n\in\mathbb N$, with identity transitions between nonbottom
fibers and ordinary coefficient products. These operations satisfy the
pair axioms directly. Every nonbottom support is above the unit, so
the ideal of Theorem~\ref{thm:upper-unit-quotient} is zero.

Here $U(Q)=\{0\}$. At support $0$ the unit test gives
$K^{\mathrm{id}}_0=0$. For $n>0$, a context with nonbottom exterior
supports has support $p+n+r>0$ and cannot reach an upper idempotent.
Contexts with a bottom exterior support have zero coefficients.
Therefore $K^{\mathrm{id}}_n=R$ for every $n>0$. This proves strict
containment. The coefficient $1_R$ at support $1\in\mathbb N$ has
nonzero powers at all positive supports, so the global ideal is not
nilpotent. Nevertheless its kernel in each fixed idempotent matrix
fiber has the nilpotence bound just proved.
\end{example}

\subsection{Locally nilpotent upper-fiber ideals}
We use ordinary Levitzki-radical facts; see
\cite[p.~53 and Corollary~1, p.~57]{Rich}. The associative finite-corner
argument needed here is included, without the additional hypotheses
used for alternative rings in that source.

The global semiring Levitzki radical of Olson and Jenkins
\cite[Definitions~1--2, p.~91]{OlsonJenkins} uses products vanishing at
the global zero and Bourne ideal quotients. The present construction
uses quotients of the coefficient groups and retains the support.
For example, the global radical of $R_0$ is $\{\tau\}$, since no
old-ring element has a power equal to $\tau$, whereas our old
coefficient fiber is $\operatorname{Lev}(R)$ and can be nonzero.
Thus the two notions of radical are distinct.

\begin{proof}[Proof of Theorem~\ref{thm:upper-levitzki-quotient}]
First, locally nilpotent rings are closed under extensions by ideals.
If $I\lhd B$ and both $I$ and $B/I$ are locally nilpotent, take a finite
set $X\subset B$. All words of some length $a$ in $X$ lie in $I$.
There are finitely many such words, and they generate a nilpotent
subring, of some exponent $b$. Every word of length $ab$ in $X$
vanishes. This also proves that sums of locally nilpotent ideals are
locally nilpotent: use finite sums and then the finite number of
summands needed for any finite set. Their sum is
$\operatorname{Lev}(B)$, and its quotient has no nonzero locally
nilpotent ideal, by the same extension argument.

The context tests define additive subgroups $H_q$. Transition
stability follows by transporting a test to its later source;
two-sided multiplication stability follows by replacing
$u(axb)v$ by $(ua)x(bv)$. At an upper idempotent $e$, the unit test
gives $H_e\subseteq\operatorname{Lev}(F(e))$. This is an ideal of
$F(e)$ and is locally nilpotent. Conversely, any coefficient ideal
locally nilpotent at all these fibers passes every context test, which
proves maximality. The zero-test ideal $K^{\mathrm{id}}$ is contained
in $H$. For a quotient coefficient ideal with the stated property,
its inverse-image fiber at $e$ is an extension of $H_e$ by a locally
nilpotent ideal. It is locally nilpotent, so maximality makes that
inverse image $H$. This proves the residual assertion. Pair
isomorphisms preserve all tests.

We next prove the ordinary finite-corner fact. Let $B$ be unital with
orthogonal idempotents $c_1,\ldots,c_n$ summing to $1$, and let
$J\lhd B$ have all $c_iJc_i$ locally nilpotent. Then $J$ is locally
nilpotent. Induct on $n$. For a finite subset of $J$, include all its
Peirce components $c_ixc_j$. Put $c=c_1$ and $f=1-c$. Induction
makes $fJf$ locally nilpotent, so words in the finitely many
$f$-to-$f$ generators vanish at some length $a$. In a nonzero block
word, an excursion from $c$ through $f$ and back therefore has the
form
\[
 x_{cf}y_1\cdots y_sx_{fc},\qquad 0\leq s<a.
\]
These finitely many coefficients, together with the direct
$c$-to-$c$ generators, lie in $cJc$ and generate a nilpotent subring,
of some exponent $b$. A nonzero path has fewer than $b$ complete
$c$-to-$c$ pieces. Its initial and terminal $f$ segments and every
excursion have bounded length. Thus sufficiently long words in the
finite component alphabet vanish. This proves the corner fact for
arbitrary associative rings.

Fix $E^2=E$ and put $D=E\vee I_n$, entrywise. Then $D^2=D$,
$D_{ii}=e_i=E_{ii}\vee1\in U(Q)$, and the off-diagonal entries are
unchanged. The ordinary ring $B=A_D(T)$ is unital: its identity has
the unit of $F(e_i)$ on the $i$th diagonal and zero coefficients
elsewhere. Indeed,
\[
 e_iE_{ij}=E_{ij}=E_{ij}e_j.
\]
Its coordinate idempotents $c_i$ have corners
$c_iBc_i\simeq F(e_i)$. These are identities in one fixed ordinary
ring fiber, not a change of the semiring identity support.

Entrywise transition gives a ring map
$\phi:A_E(T)\to A_D(T)$. Its kernel has only diagonal entries, in
$\ker(F(E_{ii})\to F(e_i))$. Such an entry annihilates every entry
in its row or column: the product factors through its zero image in
$F(e_i)$ by the displayed support identities. Therefore
\[
 (\ker\phi)A_E=0=A_E(\ker\phi).
\]
The ideal $N_D\lhd B$ with entries $H_{D_{ij}}$ has corner ideals
$H_{e_i}$, all locally nilpotent. The finite-corner fact makes $N_D$
locally nilpotent. The kernel $N_E$ of $A_E(T)\to A_E(T/H)$ maps
into $N_D$ under $\phi$, and its remaining kernel has square zero.
The extension argument therefore makes $N_E$ locally nilpotent.

Lift a quotient idempotent at its exact support $E$ to $a\in A_E$.
Its defect $d=a^2-a\in N_E$ is nilpotent. The correction
$a\mapsto3a^2-2a^3$ changes $d$ to $d^2(4d-3)$, so iteration gives
an exact idempotent without a constant term or change of support.
For an isomorphism, lift its rectangular inverse intertwiners and cut
by lifted idempotents $P,Q$. The resulting composites are
$XY=P+\delta$ and $YX=Q+\epsilon$. Each error is nilpotent, hence
both composites are invertible by finite geometric series in their
corners. As before, $Y'=Y(XY)^{-1}=(YX)^{-1}Y$ is an exact inverse.
This proves existence and uniqueness of lifted projective classes.
Extension preserves direct sums, proving the $V$ and $K_0$ assertions.

Finally let $S=M_m(T)$. The entrywise family $M_mH(T)$ is a
coefficient ideal. At every upper idempotent support matrix its fiber
is locally nilpotent by what was just proved, so maximality gives
$M_mH(T)\subseteq H(S)$. Conversely take $X=(x_{ij})\in H(S)_P$.
For a scalar test $u\in F(p)$, $v\in F(r)$,
$pP_{ij}r\leq e\in U(Q)$, put $u$ in position $(1,i)$ of $U$ and
$v$ in position $(j,1)$ of $V$, with global semiring zeros elsewhere.
Transport $UXV$ to the upper idempotent support
$D=\operatorname{diag}(e,1,\ldots,1)$. Its coefficient belongs to
$\operatorname{Lev}(A_D(T))$. The first ordinary Peirce corner is
$F(e)$, and the corner of a locally nilpotent ideal is locally
nilpotent. Thus the selected scalar coefficient
$t_{pP_{ij}r,e}(ux_{ij}v)$ belongs to $\operatorname{Lev}(F(e))$.
Every entry passes all tests, giving $H(S)\subseteq M_mH(T)$.
This argument allows $F(0)\ne0$; it does not identify $A_D$ with a
direct product. Entrywise quotient gives the asserted matrix
isomorphism.
\end{proof}

\begin{example}[ordinary boundary and the residual condition]
\label{ex:levitzki-boundary}
For an ordinary ring $R$, its one-point support gives
$H=\operatorname{Lev}(R)$ and $K^{\mathrm{id}}=0$. Dual numbers over
a field give strict containment. The rings $k[t]/(t^a)$ have
$1\times1$ fiber kernels $(t)$ of exponent $a$, so no nilpotence
bound depending only on matrix size exists.

Individual upper fibers need not be radical-free after the reduction.
Take $Q=\{\bot<1<e\}$, with maximum for join and for nonbottom
multiplication. Put $F(\bot)=0$, $F(1)=k[t]/(t^2)$ and
$F(e)=M_2(k)$, with the injective unital transition
$a+bt\mapsto aI+bE_{12}$. Products use this map and ordinary ring
products, so the pair axioms hold. Since $M_2(k)$ has no nonzero
locally nilpotent ideal, $H_e=0$; injectivity of the transition then
gives $H_1=0$. Thus $H=0$, although
$\operatorname{Lev}(F(1))=(t)\ne0$. The residual property is exactly
the absence of a compatible coefficient ideal of the specified kind.
\end{example}

The zero-detection ideal alone need not commute with matrices. In
Example~\ref{ex:upper-idempotent-strict}, the support matrix with
integer-zero diagonal, upper off-diagonal $1\in\mathbb N$ and lower
off-diagonal $\bot$ is an upper idempotent. The entrywise
$K^{\mathrm{id}}$ has a nonzero upper off-diagonal group there,
whereas $K^{\mathrm{id}}(M_2(T))$ vanishes on the entire fiber by its
unit test. The locally nilpotent quotient includes this ordinary
triangular radical and has the matrix compatibility just proved.


\subsection{Evaluating the compatible radical by ring transitions.}
The context tests of Theorem~\ref{thm:upper-levitzki-quotient} can be
replaced by transitions between explicit ordinary rings. This evaluates
the same coefficient ideal; it does not classify the remaining projectives.

\begin{proposition}\label{prop:upper-transition-radical}
For $e\in U(Q)$, the compatible coefficient ideal of
Theorem~\ref{thm:upper-levitzki-quotient} satisfies
\[
 H_e=\bigcap_{\substack{f\in U(Q)\\e\leq f}}
       t_{e,f}^{-1}\bigl(\operatorname{Lev}(F(f))\bigr).
\]
For any $q\in Q$, put
\[
 E(q)=\begin{pmatrix}1&q\\0&1\end{pmatrix},
 \qquad
 j_q:F(q)\longrightarrow A_{E(q)}(T),
\]
where $j_q(x)$ has coefficient $x$ in position $(1,2)$ and zero
amplitude in the other specified groups. Then
\[
 H_q=\bigcap_{\substack{D\in M_2(Q)\\D^2=D,\ E(q)\leq D}}
 (t_{E(q),D}j_q)^{-1}\bigl(\operatorname{Lev}(A_D(T))\bigr).
\]
All target rings are unital. If $Q$ is finite, the second intersection
has at most $|Q|^4$ terms, without any finiteness assumption on $F$.
Computing the ordinary radicals remains a separate problem.
\end{proposition}

\begin{proof}
We first recall the ordinary transfer argument underlying the corner
result of Rich~\cite[Corollary~1, p.~57]{Rich}. In a Morita context with
diagonal rings $A,B$ and off-diagonal bimodules $M,N$, let $I$ be a
locally nilpotent ideal of $B$. The additive group $MIN$ generated by
$mzn$, with $z\in I$, is an ideal of $A$. Expand a finite subset using
finitely many terms $m_i z_i n_i$. Products of these terms have the form
\[
 m_{i_1}z_{i_1}(n_{i_1}m_{i_2})z_{i_2}\cdots
 (n_{i_{s-1}}m_{i_s})z_{i_s}n_{i_s}.
\]
The finitely many elements $z_i(n_i m_j)$ belong to $I$ and generate
a nilpotent subring. Thus sufficiently long products vanish, proving
$M\operatorname{Lev}(B)N\subseteq\operatorname{Lev}(A)$.

Necessity of the first formula follows from transition stability of
$H$ and $H_f\subseteq\operatorname{Lev}(F(f))$. For sufficiency, suppose
$x\in F(e)$ passes every stated upward test. Consider any defining
context $u\in F(p)$, $v\in F(r)$ with $per\leq h\in U(Q)$. Set
\[
 a=hpe,\qquad b=erh,\qquad f=e\vee ba=e\vee erhpe.
\]
We have $ab\leq h$, $e(ba)=ba=(ba)e$, and
$(ba)^2=b(ab)a\leq bha=ba$. Consequently $f^2=f$ and $f\geq e$.
Moreover
\[
 ha=a=af,\qquad fb=b=bh.
\]
It follows that
\[
 D=\begin{pmatrix}h&a\\b&f\end{pmatrix}
\]
is an upper multiplicative idempotent. Its ordinary fiber ring
$A_D(T)$ has diagonal corners $F(h),F(f)$ and off-diagonal groups
$F(a),F(b)$. Transport $u,v$ to these off-diagonal groups; this is
possible because $p\leq a$ and $r\leq b$. By hypothesis
$t_{e,f}(x)\in\operatorname{Lev}(F(f))$. The ordinary transfer argument
therefore places their threefold product in
$\operatorname{Lev}(F(h))$. Multiplicative naturality identifies it
with $t_{per,h}(uxv)$: the support identities $ae=a=af$ and
$eb=b=fb$ justify the transports without changing their order.
Thus every original context test passes. No common upper bound for
all upper idempotents has been assumed.

Finally $E(q)^2=E(q)\geq I_2$. Finite-matrix compatibility in
Theorem~\ref{thm:upper-levitzki-quotient} gives
\[
 x\in H(T)_q\quad\Longleftrightarrow\quad
 j_q(x)\in H(M_2(T))_{E(q)}.
\]
Apply the first formula to $M_2(T)$ and $E(q)$. Its upper fiber rings
and transitions are precisely $A_D(T)$ and $t_{E(q),D}$, giving the
second formula. There are at most $|Q|^4$ support matrices when $Q$
is finite. The group $F(0)$ is retained throughout; the lower-left
entry of $j_q(x)$ is its zero element, not an assertion that the
group is zero.
\end{proof}

\subsection{A compatible Jacobson quotient}
For global semiring radicals, Katsov and Nam
\cite[Examples~3.6--3.7, pp.~14--15 of the preprint]{KatsovNamRadicals}
distinguish the Jacobson--Bourne radical $J$ from the radical $J_s$
defined by all simple semimodules. Their Theorem~5.8 and
Proposition~5.15, pp.~32--36, establish matrix and corner formulas
for these radicals. Those are known global ideal results. Here
$\operatorname{Jac}$ always denotes the ordinary radical of a
coefficient ring, and the quotient retains the support. The
nonunital fixed-fiber and exact-support arguments below are included
to make that distinction explicit.

\begin{lemma}[a finite Peirce test]\label{lem:finite-peirce-jacobson}
Let $B$ be a unital ring with finitely many orthogonal idempotents
$c_i$ summing to $1$. If $N\lhd B$ and
$c_iNc_i\subseteq\operatorname{Jac}(c_iBc_i)$ for every $i$, then
$N\subseteq\operatorname{Jac}(B)$.
\end{lemma}
\begin{proof}
Omit zero corners and induct on their number. Put $c=c_1$ and
$f=1-c$. The induction hypothesis applies to $fNf$. For $x\in N$,
$c-cxc$ is invertible in $cBc$, and the Schur complement of
$1-x$ in $fBf$ is
\[
 f-fxf-fxc(c-cxc)^{-1}cxf.
\]
Its defect from $f$ belongs to $fNf$, so it is invertible.
Two-by-two block elimination now inverts $1-x$ in $B$.
The same argument applies to every multiple of $x$, because $N$
is an ideal. This is the Jacobson criterion.
\end{proof}

\begin{proof}[Proof of Theorem~\ref{thm:upper-jacobson-quotient}]
The defining inverse images are additive subgroups. Transporting a
test upward proves transition stability; replacing $u(axb)v$ by
$(ua)x(bv)$ proves the two-sided coefficient-ideal property.
The unit test at $e$ gives $C_e\subseteq\operatorname{Jac}(F(e))$.
Conversely every compatible ideal with that upper property passes
all tests, proving maximality. A locally nilpotent ideal of a
unital ring is contained in its Jacobson radical: every multiple
of one of its elements is nilpotent, and the finite geometric
series inverts its complement to $1$. Hence $H\subseteq C$.

If $I\subseteq\operatorname{Jac}(A)$, the maximal-left-ideal
correspondence gives
$\operatorname{Jac}(A/I)=\operatorname{Jac}(A)/I$.
A compatible ideal in $T/C$ with the required upper property
therefore pulls back to one with that property in $T$.
Maximality forces the pullback to equal $C$, proving the residual
assertion.

Fix $E^2=E$ and set $D=E\vee I_n$. As in the proof of
Theorem~\ref{thm:upper-levitzki-quotient}, $D^2=D$, $DE=ED=E$,
and $B=A_D(T)$ is unital with diagonal corners
$F(e_i)$, where $e_i=E_{ii}\vee1\in U(Q)$. Write
$A=A_E(T)$ and let $\phi:A\to B$ be entrywise transition.
Its kernel annihilates $A$ on both sides. The same support
identities also give actual left and right $B$-actions on $A$;
transition of these products is multiplication in $B$.
Consequently both $\phi(A)$ and
\[
 M=\phi(N_E)
\]
are two-sided ideals of $B$, the second by compatibility of $C$.
This ideal property is essential; radical containment is not
inherited by arbitrary subrings.

Let $N_D$ be the ideal of $B$ with entries $C_{D_{ij}}$.
Its diagonal corners lie in $\operatorname{Jac}(F(e_i))$, so
Lemma~\ref{lem:finite-peirce-jacobson} gives
$N_D\subseteq\operatorname{Jac}(B)$. In particular
$M\subseteq\operatorname{Jac}(B)$. The quasi-inverse of an
element of $M$ stays in $M$, since the correction to the inverse
of $1-x$ is a product involving $x$ and $M$ is an ideal.
Lift that quasi-inverse to $N_E$. In a unitization, the left and
right inverse defects lie in $\ker\phi$, hence have square zero.
Correcting these defects gives a left and a right inverse of
$1-x$, both with corrections in $N_E$. They agree. Thus $N_E$
is a radical nonunital ring. Compressing a quasi-inverse by an
idempotent $P$ gives a quasi-inverse relative to the identity $P$.
Applying this also to all corner multiples proves
$PN_EP\subseteq\operatorname{Jac}(PAP)$.

For matrix compatibility, the entrywise family $M_mC(T)$ is a
coefficient ideal. At every upper idempotent matrix support,
the preceding finite Peirce test puts its fiber in the ordinary
Jacobson radical; maximality gives
$M_mC(T)\subseteq C(M_m(T))$. Conversely, take an element of the
right side at support $P$. Isolate its $(i,j)$ coefficient with
selector matrices having $u\in F(p)$ in position $(1,i)$ and
$v\in F(r)$ in position $(j,1)$, and global zeros elsewhere.
For $pP_{ij}r\leq e\in U(Q)$, transport the product to
$\operatorname{diag}(e,1,\ldots,1)$. Its first corner is
$t_{pP_{ij}r,e}(ux_{ij}v)$ in $F(e)$. The corner of a Jacobson
radical lies in the corner radical: restrict the inverse of
$1-z$, and apply the same argument to all corner multiples.
Thus every coefficient passes all scalar tests. This proves the
reverse inclusion, including when $F(0)\ne0$.

Given an isomorphism of reduced projectives, lift its rectangular
intertwiners at their specified supports and cut by the original
projectors $P,Q$. Their composites are $P+\delta$ and
$Q+\epsilon$, with errors in the radical ideals of the projector
corners just proved. Both composites are invertible. The usual
correction
$Y'=Y(XY)^{-1}=(YX)^{-1}Y$
is an exact inverse to $X$ and has the prescribed reduction.
This proves lifting of isomorphisms, reflection of isomorphism
and injectivity on $V$.

Now assume the upper coefficient rings are semiperfect.
The diagonal local-idempotent decompositions of the finitely many
$F(e_i)$ assemble to one of $1_B$; hence $B$ is semiperfect.
In a semiperfect ring, idempotents lift modulo any ideal
$M\subseteq\operatorname{Jac}(B)$. Here is the needed reduction
to lifting modulo the radical. Lift the further image in
$B/\operatorname{Jac}(B)$ first. In $B/M$, two idempotents
$a,b$ with the same radical image are conjugate: the element
$u=ab+(1-a)(1-b)$ is a unit and $au=ub$.
Units lift modulo $M$ because lifted two-sided inverses have
invertible defects $1+m$. Conjugating by a lifted $u$ gives
the required exact idempotent lift.

Let $\bar P\in A_E(T/C)$ be a prescribed idempotent, and choose
a coefficient lift $a_0\in A$. Then $b_0=\phi(a_0)$ is
idempotent modulo $M=\phi(N_E)$. Lift it to an idempotent
$b\in B$ with $b-b_0\in M$. Choose $n\in N_E$ with
$\phi(n)=b-b_0$ and put $a=a_0+n$. It has the prescribed
reduction, while $a^2-a$ belongs to the square-zero kernel of
$\phi$. The correction $3a^2-2a^3$ is idempotent, has no
constant term, and still reduces to $\bar P$.
Thus every projector lifts at its exact support. The preceding
uniqueness and preservation of direct sums give the stated
$V$- and $K_0$-isomorphisms.
\end{proof}

\begin{example}[object lifting can fail]
\label{ex:jacobson-lifting-boundary}
Let $R=S^{-1}\mathbb Z$, where $S$ consists of integers coprime
to $6$. Its maximal ideals are $2R$ and $3R$, so
$\operatorname{Jac}(R)=6R$ and
$R/6R\simeq\mathbb F_2\times\mathbb F_3$.
This localization is a principal ideal domain. Every finite
projective is a finitely generated submodule of a finite free
module and is therefore free. Thus reduction on projective
monoids is
\[
 \mathbb N\longrightarrow\mathbb N^2,\qquad n\longmapsto(n,n).
\]
The projective $(\mathbb F_2,0)$ does not lift. Also
$\operatorname{Lev}(R)=0$, since $R$ is a domain, whereas
$C=6R\ne0$. This proves both strict enlargement of the compatible
locally nilpotent ideal and the necessity of a lifting hypothesis.
Monoid injectivity alone supplies no claim of $K_0$-injectivity
for general inverse semirings.
\end{example}

\begin{corollary}[Jacobson tests by upward ring transitions]
\label{cor:jacobson-transition-tests}
For $e\in U(Q)$, the compatible Jacobson family satisfies
\[
 C_e=\bigcap_{\substack{f\in U(Q)\\e\leq f}}
 t_{e,f}^{-1}\bigl(\operatorname{Jac}(F(f))\bigr).
\]
With $E(q)$ and $j_q$ as in
Proposition~\ref{prop:upper-transition-radical}, for any $q\in Q$,
\[
 C_q=\bigcap_{\substack{D\in M_2(Q)\\D^2=D,\ E(q)\leq D}}
 (t_{E(q),D}j_q)^{-1}\bigl(\operatorname{Jac}(A_D(T))\bigr).
\]
All target rings are unital. For finite $Q$, the second formula
has at most $|Q|^4$ terms; computing their ordinary radicals is
a separate problem.
\end{corollary}
\begin{proof}
In any unital ring $B$ with an idempotent $c$, an element
$x\in\operatorname{Jac}(cBc)$, embedded as a single diagonal
block, belongs to $\operatorname{Jac}(B)$. Indeed $1-yx$ is
block triangular for every $y\in B$, with invertible diagonal
blocks $c-cycx$ and $1-c$. The converse corner inclusion follows
by restricting inverses. This is the ordinary ring case of
\cite[Proposition~5.15, p.~36 of the preprint]{KatsovNamRadicals}.
Consequently a sandwich through one radical diagonal corner
lands in the radical of the other diagonal corner.

Necessity of the first formula follows from transition stability.
For sufficiency, use the same context and supports as in the proof
of Proposition~\ref{prop:upper-transition-radical}: for
$per\leq h\in U(Q)$ put
$a=hpe$, $b=erh$ and $f=e\vee ba$. The support identities there
give the upper idempotent matrix
$\left(\begin{smallmatrix}h&a\\b&f\end{smallmatrix}\right)$.
The upward test puts $t_{e,f}(x)$ in the Jacobson radical of its
second corner. The preceding ordinary sandwich fact puts
$t_{per,h}(uxv)$ in the radical of its first corner. Thus every
defining context test passes, in the original noncommutative order.
Finally apply this upper formula to $M_2(T)$ and $E(q)$, using
the matrix compatibility in
Theorem~\ref{thm:upper-jacobson-quotient}. The other coefficients
of $j_q(x)$ are their group zeros, including in $F(0)$, so no
vanishing hypothesis on that group is needed. Counting the
possible matrices gives the finite bound.
\end{proof}


\begin{example}[three different radical conventions]
\label{ex:jacobson-conventions}
Let $k$ be a nonzero field and $Q=\{\bot<1<u\}$, with maximum
for join and nonbottom multiplication. Put $F(\bot)=F(u)=0$
and $F(1)=k$. The unique transitions and ordinary field product
define an inverse semiring $T$. Write $\tau,e,z$ for the zeros
of its three groups. The element $z$ absorbs addition, and
multiplication by a non-$\tau$ element sends $z$ to $z$.

The unit test gives $C_1\subseteq\operatorname{Jac}(k)=0$,
so every coefficient subgroup $C_q$ is zero. Its underlying
union is $Z=\{\tau,e,z\}$, and $T/C=T$. In the right-semiregularity
definition of the Jacobson--Bourne radical
\cite[p.~14 of the preprint]{KatsovNamRadicals}, the whole ideal
$T$ is semiregular: for any $i_1,i_2$, choose $j_1=j_2=z$ in
\[
 i_1+j_1+i_1j_1+i_2j_2=i_2+j_2+i_1j_2+i_2j_1.
\]
Both sides are $z$. Hence $J(T)=T$ and its Bourne quotient is
one point.

In contrast, the homomorphism $T\to\B$ sending $\tau$ to $0$
and every other element to $1$ defines a simple two-element
$T$-semimodule. It has no nontrivial subsemimodule or congruence,
and its annihilator is exactly $\{\tau\}$. Intersecting all
simple-semimodule annihilators therefore gives
$J_s(T)=\{\tau\}$. Finally $1+z=z$ with $z\in Z$ and
$1\notin Z$ proves that $Z$ is not subtractive. Thus all three
selected subsets, and the two indicated quotient conventions,
are different.
\end{example}

\subsection{Lifting through coefficient ideals with exchange upper rings}
We use ordinary exchange theory with its established scope.
The fixed-cardinality exchange property passes to finite direct
sums and direct summands, and a module has the two-summand
exchange property exactly when its endomorphism ring satisfies
the exchange equation. Complete proofs are given in
\cite[Lemmas~2.3--2.4, Proposition~2.7 and Corollary~2.8,
pp.~2--6]{NielsenExchange}. The ideal-lifting implication is
\cite[Proposition~1.1 and Corollary~1.3, p.~270]{NicholsonExchange}.
We include the short calculation needed here.

The terminology also has a global semiring convention.
Das and Kar require absorbing complements for every
multiplicative idempotent in their definition of an exchange
semiring \cite[Section~2 and Definition~3.11, pp.~4 and~10]{DasKarExchange}.
Our hypothesis is the ordinary exchange equation in each
ring $F(e)$; no identification of these conditions is used.

\begin{proof}[Proof of Theorem~\ref{thm:exchange-coefficient-lifting}]
First, if a unital ring $B$ has a finite complete orthogonal
family $p_i$ and every $p_iBp_i$ is exchange, then $B$ is
exchange. Indeed, the right modules $p_iB$ have endomorphism
rings $p_iBp_i$, by left multiplication. The cited module
criterion and finite-sum property show that
$B_B=\bigoplus_i p_iB$ has the two-summand exchange property.
The same criterion applies to its endomorphism ring $B$.

In an exchange ring, let $b^2-b$ belong to a two-sided ideal
$M$. Choose $p^2=p=rb$ and $1-p=s(1-b)$ by the exchange
equation. Then
\[
 p-b=p(1-b)-(1-p)b=(r-s)(b-b^2)\in M.
\]
Thus idempotents lift modulo every two-sided ideal. The
calculation also works for left ideals, though that extra scope
is not needed below.

Fix $E^2=E$ and put $D=I\vee E$. Then
\[
 D^2=D,\qquad DE=E=ED.
\]
Write $A=A_E(T)$ and $B=A_D(T)$. The latter is a unital ring
whose coordinate idempotents have diagonal corners $F(d_{ii})$.
Each $d_{ii}=1\vee e_{ii}$ is an upper multiplicative
idempotent, since $e_{ii}^2\leq e_{ii}$. By the hypothesis
and the preceding finite-corner argument, $B$ is exchange.

The upward transition gives a ring map $\phi:A\to B$.
Its kernel $K$ annihilates $A$ on both sides. For example,
if $\phi(x)=0$ and $y\in A$, transition compatibility and
$DE=E$ identify $xy$ with the mixed product $\phi(x)y$ in
the fiber at $E$, hence with zero. The other order uses
$ED=E$. These are ordinary group zeros at their specified
supports; no vanishing assumption on $F(0)$ is used.

Put $N=\prod_{i,j}\mathcal I_{e_{ij}}\subseteq A$.
Componentwise quotient identifies $A/N$ with
$A_E(T/\mathcal I)$. The image $M=\phi(N)$ is a two-sided
ideal of $B$. Indeed, a mixed product of a $D$-supported
coefficient matrix and an $E$-supported matrix in $N$ again
has support $E$ and coefficients in $\mathcal I$, by
$DE=ED=E$ and compatibility of the coefficient ideal.
Applying $\phi$ proves both ideal closures. This argument
requires no injectivity of the transitions.

Let $\bar a\in A/N$ be idempotent and choose a lift
$a\in A$. Then $\phi(a)$ is idempotent modulo $M$.
The ordinary exchange calculation gives $p^2=p\in B$
with $p-\phi(a)\in M$. Choose $n\in N$ with
$\phi(n)=p-\phi(a)$ and set $a_0=a+n$. Its defect
$a_0^2-a_0$ lies in $K$, so the annihilator property gives
\[
 a_0^3=a_0^2,\qquad (a_0^2)^2=a_0^2.
\]
The idempotent $a_0^2\in A_E(T)$ reduces to
$\bar a^2=\bar a$. All operations stay in that exact
additive matrix fiber, including entries whose coefficients
are group zeros. This proves the projector-lifting assertion.

A finite projective over $T/\mathcal I$ is the split image
of such a projector. Its lift has a finite-projective split
image over $T$, and extension of scalars preserves this split
retract of a finite free semimodule. Reduction respects direct
sums, so it surjects on $V$ and on its group completion.
For $\mathcal I=C(T)$, the unconditional isomorphism-detection
part of Theorem~\ref{thm:upper-jacobson-quotient} also gives
injectivity on $V$. Its bijection is a monoid isomorphism and
therefore induces the asserted $K_0$-isomorphism.
\end{proof}

The ordinary exchange property is used for the whole finite
upper coefficient-matrix ring. Scalar idempotent lifting in
each diagonal ring alone would not justify that step. The
theorem concerns existence of projective lifts; it does not
supply the local decompositions or uniqueness used in the
semiperfect amplitude classification.

\subsection{Removing transition kernels over integral support}
\label{subsec:integral-support-reduction}

\begin{proof}[Proof of Theorem~\ref{thm:integral-support-reduction}]
Write $j_q:F(q)\to R=F(1)$ for the transition, available because
$q\leq1$. Multiplication with $F(1)$ and naturality show that
$I_q=\im j_q$ is an $R$--$R$-ideal. Functoriality gives
$I_q\subseteq I_p$ when $q\leq p$, and multiplicative naturality gives
$I_qI_p\subseteq I_{qp}$. The products are the ordinary products in
$R$ and are independent of representatives. The inclusion diagram
satisfies the unit, associativity and distributivity laws. Proper
ideals are not required to have their own units.

For the ideal $K$ of Theorem~\ref{thm:upper-unit-quotient}, the unit
context gives $K_q\subseteq\ker j_q$. Conversely, if $j_q(x)=0$,
then every product $uxv$, with $u\in F(p)$ and $v\in F(r)$,
transports to $j_p(u)j_q(x)j_r(v)=0$ in $F(1)$. Since
$pqr\vee1=1$, all defining context tests vanish and $x\in K_q$.
Thus $K_q=\ker j_q$, and the quotient $F(q)/K_q$ identifies with
$I_q$, compatibly with all transitions and products.

The matrix-fiber nilpotence bound and the projective-class, monoid
and group isomorphisms now follow from the preceding theorem.
Its full lifting proof preserves exact supports and includes zero
amplitudes at nonzero supports. This is the integral specialization
of the unrestricted reduction, not a claim that arbitrary
coefficient-kernel products vanish before transport.
\end{proof}

\subsection{Multiplicities over a local unit ring}
\label{subsec:local-unit-projectives}

\begin{proof}[Proof of Theorem~\ref{thm:local-unit-projectives}]
First use Theorem~\ref{thm:integral-support-reduction} to replace $F$ by
its image-ideal diagram. Fix a finite-projective support $L$.
Theorem~\ref{thm:action-colimits} identifies its action category with
the ideal category
\begin{equation}\label{eq:local-ideal-category}
 \mathcal A_L(l,k)=H(l,k):=\bigcup_{ql\leq k}I_q\subseteq R.
\end{equation}
The union is directed by finite joins of the scalar indices. Since the
transition maps are inclusions, this is an actual union, and composition
is the ordinary product in $R$. In particular $H(l,l)=R$, because the
index $1$ is available and is greatest.

The set $W$ contains $1$, is upward closed, and is closed under
multiplication: if $I_q=I_p=R$, then
$R=I_qI_p\subseteq I_{qp}$. Hence $\preceq_W$ is a preorder. Moreover
\[
 H(l,k)=R\quad\Longleftrightarrow\quad l\preceq_W k.
\]
Indeed, the union in \eqref{eq:local-ideal-category} contains $1_R$
exactly when one of its constituent ideals contains $1_R$.

Let $P_l=\mathcal A_L(l,-)$ be the covariant representable. Yoneda gives
$\End(P_l)\simeq R^{\mathrm{op}}$, a nonzero local ring, so $P_l$ is
indecomposable. Two representables $P_l,P_k$ are isomorphic exactly
when there are inverse coefficients in $H(l,k)$ and $H(k,l)$.
Inverse coefficients are units of $R$, so both ideals are $R$.
Conversely, if both ideals are $R$, their coefficients $1_R$ give
inverse maps. Thus
\[
 P_l\simeq P_k\quad\Longleftrightarrow\quad l\sim_W k.
\]

Every finite projective amplitude is a retract of a finite direct sum
$P_{l_1}\oplus\cdots\oplus P_{l_s}$. The ordinary Krull--Schmidt
argument applies even though $\mathcal A_L$ may have infinitely many
objects. To see the finiteness explicitly, take the endomorphism ring
$\Gamma$ of this finite sum. The corresponding decomposition of its
regular module has local endomorphism rings. By
\cite[Proposition~4.1, pp.~8--9]{Krause}, $\Gamma$ is semiperfect and
its finite projectives form a Krull--Schmidt category. The additive
Yoneda equivalence for the direct summands of finite sums identifies
the original retract with one of these projectives. Uniqueness and
the description of direct summands in
\cite[Theorem~4.2 and Corollary~4.3, p.~9]{Krause} show that it is a
finite sum of the same representable isomorphism types, with unique
multiplicities. These locators refer to the author's manuscript. No
Artinian or finite-ring hypothesis has been used.

It remains to impose the canonical ideal condition. For any ideal
$\mathcal J$ of an additive category,
\[
 (\mathcal J P_l)(k)=\mathcal J(l,k).
\]
Composition gives one inclusion, and composition with $1_l$ gives the
other. Therefore $\mathcal J P_l=P_l$ exactly when
$1_l\in\mathcal J(l,l)$. For the present inclusion diagram, the
finite-frame formula gives
\[
 \mathcal J_L(l,l)=\bigcup_{q\in C(l,l)}I_q.
\]
This ideal is $R$ exactly when $C(l,l)$ meets $W$. The intrinsic
description $C(l,l)=\mathord\downarrow\{\psi(l):\psi\in D_l\}$ and
upward closure of $W$ give the equivalent admissible-functional
condition. It is invariant under isomorphism of representables because
$\mathcal J_L$ is two-sided, proving independence of the representative.

The ideal acts componentwise on a direct sum. Consequently
$\mathcal J_LG=G$ exactly when all representable types occurring with
positive multiplicity are allowed. Theorem~\ref{thm:canonical-ideal}
now proves necessity and sufficiency of the stated data.

Formula~\eqref{eq:local-unit-fibers} is the resulting sum of
representables. If $x\in I_p$ and $y\in I_q$ with $ql\leq k$, then
$xy\in I_{pq}$ and $(pq)l\leq pk$, so left multiplication has the
asserted target. The transitions, naturality, associativity and unit
law follow from inclusion and ordinary ring multiplication. Support
isomorphisms preserve all scalar inequalities and transport the
allowed classes; conversely a semimodule isomorphism induces a support
isomorphism and, by the uniqueness just proved, the same multiplicities.

Finally lift the finite idempotent for this separated projective along
the coefficient quotient. The preceding theorem gives its unique
original projective isomorphism class. Restriction of the separated
module along the quotient is not the inverse construction. A zero
multiplicity function gives the zero-amplitude module on $L$, which
is the zero semimodule only if $L$ is itself the zero support.
\end{proof}

If $I_0=R$, all support objects are $\sim_W$-equivalent and allowed,
leaving the ordinary free rank over the local ring together with $L$.
If $I_0$ is proper, singleton support admits no nonzero projective
amplitude. If the unit ring is zero, outside the nonzero-local
hypothesis, the unit law forces every coefficient group to be zero,
so only the projective $Q$-support remains.

\subsection{Constant ring fibers over an entire idempotent semiring}
\label{sec:entire-support}

We now evaluate the amplitude condition when all nonzero coefficient fibers
are one ring and their transition maps are identities. The support semiring
may be infinite and need not be commutative. Its own projective modules
remain part of the classification data.

We prove Theorem~\ref{thm:entire-support-projectives}.



\begin{proof}[Proof of Theorem~\ref{thm:entire-support-projectives}]
First suppose that $M$ is the image of an idempotent endomorphism of
$T^n$. For left modules we use row coordinates: if the image of the
$i$th basis vector has $j$th coefficient $e_{ij}$, the endomorphism is
$x\mapsto xE$, where $(xE)_j=\sum_i x_i e_{ij}$. This convention
keeps the coefficient order correct when the semirings are
noncommutative.

Project $E$ to matrices $E_Q$ over $Q$ and $E_0$ over $R_0$, with
common Boolean support matrix $B$. Let $A$ be the ordinary $R$-matrix
obtained from $E_0$ by replacing absent entries by ring zero. Each
matrix is idempotent, and $A_{ij}=0$ wherever $B_{ij}=0$. The converse
need not hold, since an occupied coefficient can have amplitude zero.
The support of $M$ is
\[
 L=\{l\in Q^n:lE_Q=l\}.
\]
It is a finite free retract. Applying $\B\otimes_Q-$ to the support
retraction identifies $D$ with the image of $B$ in $\B^n$, hence with a
finite projective Boolean semimodule. The finite Boolean classification
therefore makes $D$ a distributive lattice.

For $q\in Q^n$ let $I(q)$ be its set of nonzero coordinates.
Entirety and idempotent addition give
\[
 I(qE_Q)=I(q)B.
\]
Indeed, the $j$th coordinate is nonzero precisely when some
$q_i(E_Q)_{ij}$ is nonzero, equivalently when both of its factors are
nonzero. Thus, for $l\in L$, the map $b$ is $l\mapsto I(l)$ under
the preceding identification. Every fixed subset $I$ of $B$ occurs:
apply $E_Q$ to the vector whose coordinates are $1$ on $I$ and $0$
elsewhere. Its occupied set is $IB=I$.

The free group above $q$ is $R^{I(q)}$, with coordinate inclusions as
transitions. Over a fixed support $l\in L$, the group of $M$ is
\[
 \operatorname{im}\bigl(R^{I(l)}\xrightarrow{\,x\mapsto xA\,}
                              R^{I(l)}\bigr).
\]
This expression depends only on $I(l)$. On the other hand,
$N=\operatorname{im}(E_0)\subseteq R_0^n$ is a finite projective
$R_0$-semimodule with support $D$, and its group at $I\in D$ is
exactly the same image of $A$. Consequently the amplitude diagram of
$M$ is the pullback along $b$ of that of $N$. Nonzero $T$-scalars act
through their ring amplitudes. The classification of projectives over
$R_0$ identifies $N$ with the labelled-poset diagram for finite
projective $R$-modules $V_p$, so it gives the asserted description of
$M$.

For the converse, start with the stated data. The displayed groups and
actions form a $T$-semimodule: transitions are inclusions of summands,
$b(ql)=b(l)$ for $q\ne0$, and the remaining identities are the
ordinary left $R$-module laws. We construct a finite free retraction
explicitly.

Choose $Q$-linear maps
\[
 \sigma:L\longrightarrow Q^n,\qquad
 \rho:Q^n\longrightarrow L,\qquad \rho\sigma=1_L.
\]
Put $l_i=\rho(e_i)$ and $\phi_i=\sigma_i$. Then
$\phi_i(l)l_i\le l$ for every $i,l$. Booleanization gives a split
pair $\sigma_B:D\to\B^n$, $\rho_B:\B^n\to D$, with
\[
 \rho_B(e_i)=b(l_i),\qquad
 (\sigma_Bb(l))_i=\chi(\phi_i(l)).
\]
For every join-irreducible $p\in D$, the equality
\[
 p=\bigvee_{i\in\sigma_B(p)}b(l_i)
\]
supplies a coordinate $i(p)$ with $b(l_{i(p)})=p$ and
$i(p)\in\sigma_B(p)$. For this coordinate and every $d\in D$,
\[
 i(p)\in\sigma_B(d)\quad\Longleftrightarrow\quad p\le d.
\]
The forward implication follows by applying $\rho_B$; the reverse
follows from monotonicity. Hence
\begin{equation}\label{eq:entire-occupancy}
 \phi_{i(p)}(l)\ne0\quad\Longleftrightarrow\quad p\le b(l).
\end{equation}
The coordinate is chosen to be active at $p$, not merely to have
support image $p$.

Choose split $R$-linear maps
\[
 V_p\xrightarrow{u_p}R^{m_p}\xrightarrow{v_p}V_p,\qquad
 v_pu_p=1_{V_p},
\]
with $m_p\ge1$, also for a zero label. Take a free $T$-module with
one block of $m_p$ coordinates for each $p$, followed by the original
$n$ support-frame coordinates. Map the basis elements of the
$p$-block to the values of $v_p$ in the $p$-summand of
$M_{l_{i(p)}}$, and map the final basis elements to the group zeros
at $l_i$. These choices define a $T$-linear map $f:T^N\to M$.
Define a support map $\sigma':L\to Q^N$ by repeating
$\phi_{i(p)}$ in its $p$-block and using $\sigma$ in the final
coordinates. If $\rho'$ is the support map of $f$, then
\[
 \rho'\sigma'(l)=
 \bigvee_p\phi_{i(p)}(l)l_{i(p)}\vee\rho\sigma(l)=l.
\]

On the group at $l$, define a section into the free group over
$\sigma'(l)$ by applying $u_p$ to every summand with $p\le b(l)$
and putting zero amplitudes in the support-frame coordinates.
Equation~\eqref{eq:entire-occupancy} says that precisely the required
$p$-blocks have coefficient group $R$; the other blocks have the
zero group. These maps commute with transitions because both sides
use the same summand inclusions. They commute with a nonzero scalar
$(q,r)$ because $q$ preserves the set of active blocks,
$\phi_i(ql)=q\phi_i(l)$, and the maps $u_p$ are left $R$-linear.
They also commute with the zero scalar. Thus they define a
$T$-linear map $s:M\to T^N$ with support $\sigma'$. The composite
$fs$ has identity support and acts on every label as $v_pu_p=1$.
Hence $fs=1_M$, proving finite projectivity.

Finally the labels are invariants. A $T$-isomorphism induces a
$Q$-isomorphism of supports and an $R_0$-isomorphism after extension
of scalars along $T\to R_0$. Extension of scalars preserves the
split free presentation and sends its idempotent to $E_0$, so this
last module is exactly $N$. The unique birth quotients in the
$R_0$ classification recover its labels on $D$. Conversely a
support isomorphism and the indicated label isomorphisms give a
fiberwise direct-sum isomorphism commuting with transitions and
scalar action. This proves the stated isomorphism classification.
\end{proof}

For $L=Q$, the Booleanization is $\B$. Taking its single label to be
$R$ recovers the regular free semimodule $T$; taking the label to be
zero gives the projective support semimodule $Q$. If $R=0$, then
$T\cong Q$ and every label is zero, so the theorem reduces exactly
to the original support classification. For $Q=\B$ it is the
classification over $R_0$. No classification of projective
$Q$-semimodules for an arbitrary $Q$ is asserted by these examples.

\subsection{Finite support and semiperfect upper rings}\label{sec:finite-inverse}
We prove Theorem~\ref{thm:finite-inverse-projectives}. Its ordinary
semiperfect-ring input is the finite local-idempotent characterization
and projective-cover theory, not a finiteness assertion about the
coefficient groups; see Krause \cite[Proposition~4.1 and
Theorem~4.2]{Krause}.

\begin{proof}
\emph{Support and coefficient formulas.}
For fixed $l$, the admissible functionals form a nonempty finite set.
Their pointwise join preserves zero, finite joins and left $Q$-linearity
by distributivity. Its value at $k$, acting on $l$, is at most $k$,
so $\delta_l$ is admissible and greatest. No commutativity is used.
If $L$ is projective, choose a finite free retraction with coordinates
$\phi_i$ and generator images $l_i$. The identity
$k=\bigvee_i\phi_i(k)l_i$ makes $\phi_i$ admissible at $l_i$.
Replacing it by $\delta_{l_i}$ proves the support test, since each term
remains at most $k$. Conversely
\[
 \sigma(k)=(\delta_l(k))_{l\in L},\qquad
 \rho((q_l))=\bigvee_{l\in L}q_ll
\]
are $Q$-linear, and the test says exactly that $\rho\sigma=1_L$.

The set of $q$ with $ql\leq k$ is closed under joins and, since $Q$ is
finite, has greatest element $r(l,k)$. The coefficient-colimit formula
of Theorem~\ref{thm:action-colimits} therefore identifies its colimit
with $F(r(l,k))$. For an admissible $\phi$ at $l$,
$\phi(k)\leq\delta_l(k)$, and the coefficient transition gives
\[
 t_{\phi(k)l,k}a_{\phi(k),x,l}
 =t_{\delta_l(k)l,k}
   a_{\delta_l(k),F_{\phi(k),\delta_l(k)}(x),l}.
\]
Thus every canonical generator belongs to the displayed image from
$F(\delta_l(k))$. The reverse inclusion follows because $\delta_l$
is admissible. These formulas include $F(0)\ne0$.

\emph{Semiperfectness of the action ring.}
Write $e_l=r(l,l)$. Since $1l=l$, one has $1\leq e_l$. Also
$e_ll\leq l$, hence $e_l^2l\leq l$ and $e_l^2\leq e_l$ by
maximality. The reverse inequality follows from $1\leq e_l$, so
$e_l^2=e_l$. With $E_l$ denoting the identity block at the object $l$,
\[
 E_lB_LE_l=F(e_l)
\]
is exactly its ordinary unital coefficient ring. Its multiplication
needs no further support transition because $e_l^2=e_l$; its unit is
the image of the global unit of $T$.

The finitely many $E_l$ are orthogonal and sum to $1_{B_L}$.
Decompose each nonzero corner identity into a finite orthogonal family
of idempotents with local corner rings. If $f\leq E_l$ belongs to
such a family, then
\[
 fB_Lf=f(E_lB_LE_l)f,
\]
so its corner remains local in $B_L$. The combined families give a
finite orthogonal decomposition of $1_{B_L}$ into local idempotents.
The ordinary characterization makes $B_L$ semiperfect. A zero corner
contributes the empty family, and the zero ring is allowed.
No condition on off-diagonal finiteness or on radical preservation by
coefficient transitions has entered the argument.

Covariant additive $\mathcal A_L$-modules correspond to unital left
$B_L$-modules by taking the direct sum of the object groups; the inverse
takes the summands $E_lN$. This preserves finite generation and
projectivity, and identifies the actions of $\mathcal J_L$ and $I_L$.

\emph{Ordinary amplitudes and the ideal filter.}
Let $B$ be a semiperfect ring and let $J$ be its Jacobson radical.
Its simple left-module types $S_i$ have finitely generated
indecomposable projective covers $P_i$, with
$P_i/JP_i\simeq S_i$. Every finitely generated projective $P$ is a
finite sum of these covers, uniquely up to multiplicity. Here is the
usual argument for the last assertion. The top $P/JP$ is a finite
semisimple module. Form the finite sum $P'$ of the covers with those
simple multiplicities and lift its map to $P/JP$ through $P$.
The resulting $P'\to P$ is surjective by Nakayama's lemma and splits.
Its kernel $K$ is finitely generated, and equality of the tops gives
$K/JK=0$. Nakayama gives $K=0$. The multiplicities are then recovered
from the semisimple top. This uses no artinian hypothesis on $B$.

For any two-sided ideal $I$, $IS_i$ is either zero or all of $S_i$.
If $IP_i=P_i$, its image in $S_i$ is all. Conversely $IS_i\ne0$
gives
\[
 IP_i+\ker(P_i\twoheadrightarrow S_i)=P_i.
\]
The cover kernel is superfluous, so $IP_i=P_i$. Consequently a finite
sum $\bigoplus_iP_i^{n_i}$ satisfies $IP=P$ exactly when every present
type has $IS_i\ne0$; the ideal $I$ need not be idempotent or finitely
generated. Apply this with $(B,I)=(B_L,I_L)$ and use
Theorem~\ref{thm:canonical-ideal} to obtain precisely the stated
amplitudes. Zero amplitude retains the elements of $L$.

A support automorphism transports the action category and canonical
ideal, permuting the simple types and their covers. Semimodule
isomorphisms are exactly such support transports together with
amplitude isomorphisms. This gives the asserted orbit set. Finally,
a finitely generated projective $T$-semimodule is a retract of $T^n$,
so its support is a retract of the finite set $Q^n$. Thus every such
projective occurs in the classification.

\emph{The finite-semiring specialization.}
A finite unital ring is semiperfect. For completeness, repeatedly split
its finite regular module into indecomposable projectives. For any
endomorphism $f$ of a finite module, stabilization of kernels and images
gives $P=\ker(f^N)\oplus\im(f^N)$ for large $N$. On an
indecomposable summand, $f$ is therefore invertible or nilpotent;
if it is not invertible, $1-f$ is invertible. Its endomorphism ring
is consequently local. The summands of the regular module give the
required finite orthogonal local-idempotent decomposition.
Thus every finite inverse semiring satisfies the theorem's hypothesis.
Its arrow groups, and hence $B_L$, are finite. The formulas above use
only its finite operation and transition tables. They recover the
finite-semiring classification without imposing a hypothesis on
unitality of the individual lower coefficient fibers.
\end{proof}

\begin{example}[an infinite semiring with three support elements]\label{ex:finite-support-semiperfect}
Let $Q=\{\bot<1<u\}$, with addition maximum, absorbing $\bot$, and
multiplication maximum on the two nonbottom elements. For a field $K$,
take $F(\bot)=0$, $F(1)=F(u)=K$, with identity transition $F(1)\to F(u)$
and ordinary field products. These data define an inverse semiring:
all nonbottom coefficient transitions are identities and all coefficient
products are the field multiplication. Its global unit has support $1$.
Both upper coefficient rings are fields. If $K$ is infinite, the
semiring is infinite, while its support is finite and its unit is not
greatest.

For the regular support $L=Q$, order both the source columns and target
rows as $\bot,1,u$. Directly from $ql\leq k$,
\[
 (r(l,k))_{k,l}=
 \begin{pmatrix}u&\bot&\bot\\u&1&\bot\\u&u&u\end{pmatrix},
 \qquad
 B_L=\begin{pmatrix}K&0&0\\K&K&0\\K&K&K\end{pmatrix}.
\]
This is ordinary lower-triangular matrix multiplication, since every
nonzero coefficient transition and product is the indicated field map.
A $Q$-linear functional $Q\to Q$ is $q\mapsto qa$, and it is
admissible at $l$ exactly when $al\leq1$. Hence $\delta_\bot$ has
values $(\bot,u,u)$, $\delta_1$ is the identity, and $\delta_u$ is
zero. The support test holds already through the $l=1$ term, and
\[
 I_L=\begin{pmatrix}0&0&0\\K&K&0\\K&K&0\end{pmatrix}.
\]
The three diagonal simple types have their projective covers in the
three column ideals. Only the middle simple receives nonzero $I_L$
action. Its cover is the middle column, with groups $0,K,K$ and
identity transition $1\to u$. Thus the projectives with support $Q$
have precisely the groups $0,K^n,K^n$, for one integer $n\geq0$.
Every $Q$-linear support automorphism is multiplication by $a$, and
only $a=1$ is bijective; there is no further orbit identification.
The value $n=0$ retains the nontrivial support with zero amplitude,
and $n=1$ recovers the free rank-one semimodule.
\end{example}

\subsection{Amplitudes over arbitrary supports with semiperfect upper rings}
We use the ordinary semiperfect and finite Krull--Schmidt theorems
in \cite[Proposition~4.1, pp.~8--9, Theorem~4.2 and Corollary~4.3,
p.~9 of the author preprint]{Krause}. Their application below is to
a finite coefficient ring constructed from the projector, not to
the entire action category.

\begin{proof}[Proof of Theorem~\ref{thm:upper-semiperfect-amplitudes}]
Apply Theorem~\ref{thm:finite-coefficient-models} to $G$. It gives
an actual unital coefficient ring $C=A_D(T)$, an ideal $N\lhd C$,
a map $\pi:C\to B$ to a finite full action-arrow ring, and an
idempotent $e\in N$ representing $G$. Every mixed corner between
idempotents in $N$ maps isomorphically to its full action corner.
The diagonal corners of $C$ are $F(d_{ii})$ for upper idempotents
$d_{ii}$. Decompose each of their identities into finitely many
orthogonal idempotents with local corners. Together these form
such a decomposition of $1_C$, so $C$ is semiperfect. This step
makes no assertion about semiperfectness of $B$.

Put $I=CeC\subseteq N$. The finite projective left $C$-module
$Ce$ is $I$-generated. Ordinary semiperfect theory decomposes it
into columns $Cf$ from the chosen diagonal local-idempotent list.
Ideal generation passes to each occurring summand, so
$Cf=ICf\subseteq I$ and in particular $f\in I\subseteq N$.
The exact-corner assertion gives
\[
 fCf\simeq\pi(f)B\pi(f).
\]
Thus $\pi(f)$ is nonzero and has a local full action corner.
It lies in a single diagonal object and in the canonical ideal.
The decomposition of $Ce$ supplies rectangular idempotent
intertwiners over $C$. Applying $\pi$ and composing on
representables identifies $G$ with the corresponding sum of
$\mathcal A(l_i,-)\pi(f)$. These inverse matrices require no
flatness assumption. This proves existence of the asserted
summands using only an ordinary coefficient ring and exact
corners.

Conversely $\mathcal A(l,-)f$ is the image of the idempotent $f$
on one representable. It is $\mathcal J$-generated since its
generator $f$ belongs to $\mathcal J(l,l)$ and $f=f^2$.
Its endomorphism ring is the opposite of the nonzero local ring
$f\mathcal A(l,l)f$. The Yoneda corner calculation identifies its
isomorphisms exactly with the two displayed inverse-arrow
relations. Krull--Schmidt uniqueness for finite sums of objects
with local endomorphism rings proves uniqueness of the
multiplicities. Theorem~\ref{thm:canonical-ideal} now gives the
semimodule classification, and transport of $\mathcal A$ and
$\mathcal J$ gives precisely the support-automorphism orbits.

For the local specialization, let $0\ne f=f^2\in\mathcal J(l,l)$.
Represent $f$ at $q=\phi(l)$ for an admissible $\phi$. As above,
$q^2\leq q$ and $e=1\vee q$ is an upper idempotent stabilizing
$l$. The unital image of $F(e)$ in $\mathcal A(l,l)$ is local
and contains $f$. A nonzero idempotent in a local ring is its
identity, so $f=1_l$. Conversely an allowed identity is a
permitted local idempotent, as the next paragraph verifies.
\end{proof}

\begin{corollary}[the local diagonal alternative]
\label{cor:upper-local-diagonal}
Suppose all upper coefficient rings are local or zero. Put
$A_l=\mathcal A_L(l,l)$ and $I_l=\mathcal J_L(l,l)$.
If $I_l=A_l$, then $A_l$ is local or zero. Otherwise
$I_l\subseteq\operatorname{Jac}(A_l)$.
\end{corollary}
\begin{proof}
An element of $I_l$ is represented at $q=\phi(l)$ for an
admissible $\phi$, and hence lies in the unital local image of
$F(1\vee q)$ in $A_l$. If $I_l$ is proper, such an element
cannot be a unit in that image. Its complement to $1$ is therefore
a unit in $A_l$. Applying this to every element of the ideal
$I_l$ gives the ordinary Jacobson criterion.

If $1_l=[q,x]\in I_l$, represent any $a\in A_l$ by $[p,y]$
with $pl\leq l$. Then $pq=\phi(pl)\leq q$, and
$a=a1_l=[pq,yx]$ transports to $F(q)$. Thus the unital image of
$F(1\vee q)$ is all of $A_l$, proving the assertion.
\end{proof}

\begin{example}[an action ring need not be local]
\label{ex:upper-local-action-boundary}
Let $Q=\{\bot\}\sqcup\mathbb N$, with maximum for join,
$m n=m+n$ on $\mathbb N$, absorbing $\bot$, and unit the integer
$0$. For a field $k$, put $F(\bot)=0$ and let $F(n)$ be the
additive group of polynomials in $k[t]$ of degree at most $n$.
Transitions are inclusions and products are ordinary polynomial
products. These satisfy all pair axioms, with unit $1\in F(0)$.
The sole upper multiplicative idempotent is the integer $0$, whose
ring is $k$.

For the singleton zero $Q$-support, the action endomorphism ring is
$\varinjlim_nF(n)=k[t]$, which is not local: $(t)$ and $(t-1)$ are
distinct maximal ideals. The only admissible functional is the
zero functional with value $\bot$, so the canonical ideal is zero
and only zero amplitude is permitted. The domain $k[t]$ has only the idempotents $0,1$, so its
semiperfectness would imply locality. Thus local upper fibers do
not make all action endomorphism rings local or semiperfect.
\end{example}



\subsection{The intrinsic two-idempotent classification}
This subsection proves Theorem~\ref{thm:twolevelprojectives}, including the case of a noninjective bimodule map.

\begin{lemma}\label{lem:noncomm}
The labelled-poset classification in Theorem~\ref{thm:main}(4) holds for left semimodules over $R_0$ for an arbitrary unital ring $R$.
\end{lemma}
\begin{proof}
The additive inverse and fiber construction use only the scalar $-1$ and the left module laws, so their fibers are left $R$-modules. Every argument in Section~\ref{sec:projectives} using coordinate projection, split injections, submodule sums and intersections, and projective quotients is left linear. For the realizing matrices, use matrices over the opposite ring to represent endomorphisms of free left modules. Their diagonal blocks realize the given left projectives, and their off-diagonal entries are the old or new zero. The idempotence calculation is unchanged because the old zero absorbs every ring product and the poset prevents a nontrivial diagonal cycle. Uniqueness is again recovery of the successive quotient modules. Thus all parts of that proof apply.
\end{proof}

\begin{proof}[Proof of Theorem~\ref{thm:twolevelprojectives}: the injective case]
Suppose first that $f$ is the inclusion of the two-sided ideal $I$. A free $T$-semimodule $T^n$ has idempotents $2^{[n]}$ and group
\[
 F_U=I^{[n]\setminus U}\oplus R^U\subseteq R^n
\]
over $U$. This is the inverse image of $B^U\subseteq B^n$ under reduction modulo $I$. Its group transitions are inclusions.

Let $M$ be a finite free retract. Its top group is a finite projective $R$-module $P$: applying the ring reflection $T\to R$ to the retract gives an idempotent ring matrix with image $P$. More explicitly, the ring amplitudes of the image of an idempotent on $T^n$ form this image in $R^n$, and the amplitude map is injective on every group. The least group is $IP$, since the same idempotent takes $I^n$ onto $IP$. Every group of $M$ therefore embeds into $P$ and contains $IP$.

Reduction $T\to B_0$, sending the lower ideal to the new zero, takes the free fiber $F_U$ modulo $I^n$ to $B^U$. It takes the retract to a finite projective $B_0$-semimodule $N$, with the same support lattice and with groups $N_l=M_l/IP$. Surjectivity on these quotient image groups follows by lifting free representatives; the kernel is $IP$: since $P$ is the image of an idempotent on $R^n$, one has $P\cap I^n=IP$, and every group contains $IP$. By Lemma~\ref{lem:noncomm}, $N$ has the stated labelled-poset form. Its top group is $P/IP$. Each $M_l$ is the entire inverse image of $N_l$, because it already contains $IP$. This proves necessity and formula~\eqref{eq:twolevelfiber}.

Conversely start with a triple and identify $P/IP$ with $N_{\mathrm{top}}$ by $\phi$. The inverse-image groups in~\eqref{eq:twolevelfiber}, with the indicated actions, form a $T$-semimodule $M$. Choose a splitting
\[
 N\longrightarrow (B_0)^m\longrightarrow N.
\]
Its top group maps are $i_B:P/IP\to B^m$ and $r_B:B^m\to P/IP$, with $r_Bi_B=1$. Let $s$ and $\rho$ be its support maps, so $\rho s=1$. Projectivity of $P$ lifts $i_B$ to an $R$-linear map $i:P\to R^m$. Lifting the finitely many basis images lifts $r_B$ to $r:R^m\to P$.

For each support $l$, reduction modulo $I$ places $i(P_l)$ inside the coordinate module $B^{s(l)}$, so $i$ gives a $T$-linear map $M\to T^m$ with support map $s$. The corresponding condition for $r_B$ shows that $r$ gives a $T$-linear map back with support map $\rho$. The maps need not yet split on groups, but their error $d=1_P-ri$ has image in $IP$.

Choose $a:P\to R^n$ and $b:R^n\to P$ with $ba=1_P$. Since $a$ is left linear, $ad$ maps into $I^n$. It defines a $T$-linear map $M\to T^n$ with constant least support. The map $b$ defines a $T$-linear map $T^n\to M$ whose support sends the empty set to the least point and every nonempty set to the top point: on the least free group it lands in $IP$, and elsewhere it lands in $P$. Now
\[
 (i,ad):M\longrightarrow T^{m+n},\qquad
 (r,b):T^{m+n}\longrightarrow M
\]
have composite $ri+bad=ri+d=1_P$ on groups and the identity on supports. They are a retraction. Thus $M$ is finite projective.

Every semimodule isomorphism induces an isomorphism on the top ring module and on the reduced $B_0$-diagram, compatible with their common top quotient. Conversely such maps preserve all the inverse-image groups and hence give a semimodule isomorphism. This proves exactly the equivalence relation on triples, including its full diagram automorphisms.
\end{proof}

\begin{lemma}[square-zero reduction]\label{lem:twolevelnil}
For arbitrary $f$, the map $q:I(f)\to I(I\hookrightarrow R)$, equal to the identity on $R$ and to $f$ on $A$, induces a bijection on isomorphism classes of finite projectives, respecting direct sums.
\end{lemma}
\begin{proof}
Put $K=\ker f$. The Peiffer identity gives $IK=KI=0$, and the product of two lower elements of $K$ is zero. The map $q$ is surjective and preserves the two-element idempotent semiring bijectively.

We use right-module matrices, as in Proposition~\ref{prop:matrix}; applying the same argument to opposite rings gives the left-module assertion. Fix an idempotent Boolean matrix $\beta$. The additive group fiber $C_\beta$ of matrices with this support has entries in $R$ where $\beta_{ij}=1$ and in $A$ elsewhere. Since $\beta^2=\beta$, matrix multiplication stays in this group fiber and is bilinear. It is therefore an associative ring, possibly without a multiplicative identity. Its additive zero has $0_R$ in the upper positions and $0_A$ in the lower positions, and is absorbing within the fiber.

The induced ring map $C_\beta\to C'_\beta$ is surjective. Its kernel $J_\beta$ has zero $R$-entries and entries in $K$ in lower positions. Every product of two kernel matrices is the ring zero: an upper zero times a lower entry is zero, and the product of two lower kernel entries is zero. Thus $J_\beta^2=0$.

Let $a\in C_\beta$ lift an idempotent in $C'_\beta$, and put $d=a^2-a\in J_\beta$. Then $ad=da$ and $d^2=0$. In the unitization of the ring, set
\[
 e=a+(1-2a)d=a+d-2ad.
\]
The expression lies in $C_\beta$, and direct expansion gives $e^2=e$. This is the square-zero instance of the idempotent-lifting polynomial in \cite[Lemma~10.32.7, Tag~05BU]{Stacks}. It lifts every idempotent matrix without changing support, proving surjectivity on finite-projective classes.

For injectivity suppose that idempotents $E,E'$ reduce to isomorphic image projectives. Lift rectangular matrices giving that isomorphism and its inverse, and replace the lifts by $U=E'UE$ and $V=EVE'$. Their reductions still give the inverse pair. Hence
\[
 VU=E+d,\qquad UV=E'+d',
\]
with $d,d'$ in the respective square-zero kernels. They lie in their corners: $d=EdE$ and $d'=E'd'E'$. In the corner rings with identities $E,E'$, the displayed composites have inverses $E-d$ and $E'-d'$. Thus $U$ has both a left and a right inverse on the image projectives, namely $(VU)^{-1}V$ and $V(UV)^{-1}$; they coincide. Proposition~\ref{prop:matrix} gives the required isomorphism before reduction. Block sums commute with $q$, so the bijection respects direct sum.
\end{proof}

\begin{proof}[Completion of the proof of Theorem~\ref{thm:twolevelprojectives}]
Apply Lemma~\ref{lem:twolevelnil} and then the injective-case classification. The ring reflection to $R$ and the reduction to $B_0$ commute with $q$, so the gluing triples are precisely those in the statement. For a finite-support diagram, its universal map to an ordinary ring module is the map to its top group; thus its ring reflection is the indicated projective $P$. The kernel reduction can change morphism groups, which is why only the asserted isomorphism classification follows.
\end{proof}

\begin{example}
Let $R=\Z/4\Z$, $A=R$, and $f(a)=2a$. Then $I=2R$, $B=\Z/2\Z$, and $\ker f=2R$ is nonzero. The finite projectives have the same isomorphism classes as those for $R\sqcup2R$. For an injective-case triple with $P=R$ and a one-point poset labelled by $B$, its two groups are $2R\subset R$. A lift to $I(f)$ is the free rank-one semimodule, whose groups are $A\to R$, with transition multiplication by $2$. Thus the two underlying group diagrams need not be isomorphic even though the quotient map induces the stated bijection of projective classes between the two semirings.
\end{example}

\section{Repeated zero adjunction and monoid semirings}\label{sec:iteration}
This section separates genuine zero adjunction from the non-absorptive construction.

\begin{proposition}\label{prop:iterated}
Let $R^{[k]}$ be obtained from $R$ by adjoining $k\ge1$ new zeros successively. Its idempotent semiring is the chain $C_{k+1}$ under maximum and minimum, with group $R$ at the greatest element and zero groups elsewhere. Its spectrum is $\Spec R$ preceded by a chain of $k$ new points, each new point generic over all points already present. At the point introduced in step $j$, the stalk after step $k$ is $C_{k-j+2}$. At an ordinary prime $\mathfrak p$ it is $(R_{\mathfrak p})^{[k]}$.

Its semimodules are diagrams of $R$-modules over a join-semilattice $L$ equipped with zero-and-join preserving maps $p_0,\ldots,p_k$ satisfying
\[
 p_0=0,\quad p_k=1_L,\quad p_i p_j=p_{\min(i,j)},\quad p_i\le p_j\ (i\le j).
\]
The lower chain scalar $c_i$, $i<k$, sends every vector over $l$ to the group zero over $p_i(l)$. The finite projectives have the intrinsic description of Theorem~\ref{thm:iteratedprojective}.
\end{proposition}
\begin{proof}
The identities for the idempotents follow directly from the successive operations. Only the old ring has a nontrivial additive group. For any semiring $B$, an ideal of $B_0$ other than the new-zero ideal contains the old zero, by multiplying any old member by it. It is consequently $J$ together with the new zero for an ideal $J$ of $B$. Primality is preserved, and the new-zero ideal is prime. Basic opens and localization behave exactly as in Proposition~\ref{prop:localization}; induction proves the spectrum and stalk assertions.

Theorem~\ref{thm:general} now applies to the chain. An action of its idempotents is exactly a family of maps with the displayed identities. The greatest group gives the $R$-module action on every fiber; all the lower groups are zero, so their actions are forced to be the stated group zeros. Naturality and associativity impose no further conditions: the displayed identities give the chain action, $p_i(l)\le l$ gives distributivity with an old ring scalar, and the diagram maps are linear. The projective assertion is proved next.
\end{proof}

\begin{proof}[Proof of Theorem~\ref{thm:iteratedprojective}]
Write $Q=C_{k+1}$, and first consider the idempotents of a retract $M$ of $(R^{[k]})^n$. Its idempotent $Q$-semimodule $L$ is a retract of $Q^n$ by maps $s:L\to Q^n$ and $r:Q^n\to L$ with $rs=1$. Since $Q^n$ is a finite distributive lattice, Lemma~\ref{lem:lattice} makes it a retract of a finite Boolean lattice. Composing retractions proves that $L$ is finite distributive, so write $L=\Id(P)$. The map $s$ is an order embedding. Meets in $L$ obey
\begin{equation}\label{eq:retractmeet}
 x\wedge y=r(s(x)\wedge s(y)).
\end{equation}
Indeed the right side is a lower bound, and a lower bound $z$ has $s(z)\le s(x)\wedge s(y)$, so $z$ is at most that right side.

Let $t$ be the greatest element of $L$ and put $U_i=c_it$. Coordinatewise multiplication by $c_i$ in $Q^n$ is truncation at level $i$. Since $s(x)\le s(t)$, equation~\eqref{eq:retractmeet} gives
\[
 x\wedge U_i=r(s(x)\wedge c_i s(t))=r(c_i s(x))=c_i x.
\]
Thus the scalar action on $\Id(P)$ is intersection with the nested ideals $U_i$, where $U_0=\varnothing$ and $U_k=P$. Setting $h(p)=\min\{i:p\in U_i\}$ gives an order-preserving colour. A principal ideal $\downarrow p$ is fixed by $c_{h(p)}$.

Fix $p\in P$, and identify it temporarily with the join-irreducible $\downarrow p$ of $L$. Decompose $s(p)$ as the join of its coordinate vectors $c_d e_a$, where $e_a$ has level $k$ in coordinate $a$ and zero elsewhere. Applying $r$, join-irreducibility gives a coordinate with
\[
 r(c_d e_a)=p,\qquad d=s(p)_a.
\]
The element $p$ is fixed by $c_d$, so $h(p)\le d$. Since $s(p)=c_{h(p)}s(p)$, also $d\le h(p)$, and hence $d=h(p)$. For $1\le i\le h(p)$ put $x_i=c_ip=r(c_ie_a)$. Its embedded $a$th coordinate is $i$, because $s(x_i)=c_i s(p)$. In particular $x_i\ne0$ and its least fixing level is exactly $i$.

If $x_i=x\vee y$, then $s(x)_a\vee s(y)_a=i$, so one of the two coordinates equals $i$. Suppose it is $s(x)_a$. Then $c_ie_a\le s(x)$, whence $x_i=r(c_ie_a)\le x$ and $x_i=x$. This proves that $x_i$ is join-irreducible. Therefore $x_i=\downarrow p_i$ for a point of colour $i$. Since $x_i=p\wedge U_i$, the point $p_i$ is the greatest point below $p$ of colour at most $i$. This proves the necessary condition on $(P,h)$.

Conversely suppose a finite coloured poset satisfies that condition. The points $p_1<\cdots<p_{h(p)}=p$ form a chain: the greatest-element property makes them comparable, and their distinct colours make the comparisons strict. For $I\in\Id(P)$ define
\[
 \lambda_p(I)=\max\bigl(\{0\}\cup\{i\le h(p):p_i\in I\}\bigr),\qquad
 s(I)=(c_{\lambda_p(I)})_{p\in P}.
\]
Each $\lambda_p$ preserves zero and unions, and intersection with $U_j$ replaces it by $\min(j,\lambda_p)$. Thus $s:\Id(P)\to Q^P$ is $Q$-linear. Define a $Q$-linear map in the other direction by
\[
 r((c_{a_p})_{p\in P})=\bigcup_{p\in P}(\downarrow p\cap U_{a_p}).
\]
Linearity follows because truncation and union commute and $U_i\cap U_j=U_{\min(i,j)}$. For $a\le h(p)$ the ideal $\downarrow p\cap U_a$ is $\downarrow p_a$; for $a=0$ it is empty. Hence every term of $rs(I)$ is contained in $I$. If $p\in I$, then $\lambda_p(I)=h(p)$, so the term indexed by $p$ contains $p$. Consequently $rs(I)=I$. This proves sufficiency for the support module.

Return to the group fibers of a retract of $(R^{[k]})^n$. At $q\in Q^n$ the free group is $R^{J(q)}$, where $J(q)=\{a:q_a=c_k\}$. Transitions extend vectors by zero, and
\[
 J(q\vee q')=J(q)\cup J(q'),\qquad
 J(q\wedge q')=J(q)\cap J(q').
\]
Coordinate projection therefore splits each transition between fibers of the retract, exactly as in the proof of Theorem~\ref{thm:main}(4). The sum proof there uses the first displayed identity. For the intersection proof, restrict a common vector to $J(s(I))\cap J(s(I'))=J(s(I)\wedge s(I'))$ and apply the retraction; equation~\eqref{eq:retractmeet} places its image over $I\cap I'$. Thus both laws in~\eqref{eq:latticegroups} hold. Splitting the successive quotients at principal ideals gives finitely generated projective labels $P_p$ and identifies every fiber with their direct sum. If $h(p)<k$, all coordinates of $s(\downarrow p)$ are below level $k$, so its free group is zero and $P_p=0$. Only the top colour can carry a label, as claimed.

To realize every permitted choice of labels, the support retraction just constructed gives the idempotent $Q$-matrix $B=sr$ whose entries, written as levels, are
\[
 b_{pq}=\lambda_p(\downarrow q)
 =\max\bigl(\{0\}\cup\{i\le h(p):p_i\le q\}\bigr).
\]
At every point of colour $k$, choose an idempotent ring matrix $A_p$ of positive size $n_p$ with image $P_p$. Use a coordinate block of size $n_p$ there and one coordinate at each lower-colour point. Replace each entry of $B$ by a constant block of its indicated size. This larger support matrix is still idempotent: repetitions do not change maxima, and every block has positive size.

Lift it to a matrix $E$ over $R^{[k]}$ as follows. Each entry of level $i<k$ has the unique possible scalar $c_i$. A level-$k$ entry occurs exactly when $h(p)=k$ and $p\le q$; this also forces $h(q)=k$. Put $A_p$ in the diagonal block at such a $p$, and put the all-$0_R$ matrix in every off-diagonal level-$k$ block. The support of $E^2$ equals that of $E$. At a level-$k$ diagonal block, the only products contributing nonzero ring coefficients are those within $A_p$, and they give $A_p^2=A_p$. At a level-$k$ off-diagonal block, every ring product contains an off-diagonal old-zero factor; all lower-support contributions also become old zero when transported to that level. At lower levels the group is zero, so support equality already implies equality of the entries. Therefore $E^2=E$.

Over the support $s(I)$, its group action is the direct sum of the $A_p$ with $p\in I$ of colour $k$: every other ring entry is zero. The image group is consequently the required sum of labels. All group transitions are inclusions. The lower scalars have zero group components and act as stated. This realizes the given data as the image of an idempotent, and hence as a finite projective.

The $Q$-action intrinsically recovers the finite lattice, its join-irreducibles and their colours. The quotients at a principal ideal and its predecessor ideal recover the labels. This proves uniqueness up to the asserted isomorphisms. Direct sum takes disjoint unions of the coloured labelled posets, so unique connected components give the formulas for $V$ and $K_0$. The rank-one free semimodule has a chain of $k$ points, one of each colour, with label $R$ at its top point. Every connected component of the poset for a free module is this chain. A direct summand of a finite free module must therefore be a disjoint union of these chains, and is itself free. Its rank is unique, since its idempotent lattice has $(k+1)^n$ elements in rank $n$. This proves $SK_0(R^{[k]})=\Z$.
\end{proof}

\begin{example}
For two zero adjunctions, the three-element scalar chain itself is obtained by taking $R=0$ and $k=2$. A projective support poset has colours $1$ and $2$, and every colour-$2$ point must have a greatest colour-$1$ point below it. Thus a single colour-$2$ point is forbidden; a single colour-$1$ point is permitted and gives the Boolean module on which both positive scalars act as the identity. A two-point chain of colours $1,2$ gives the free module. A fork with two incomparable colour-$1$ points below a colour-$2$ point fails the greatest-element condition, and its ideal lattice is not projective over this chain despite being projective over $\B$.
\end{example}

\begin{proposition}\label{prop:absorber}
Let $B$ be a commutative semiring and let $A(B)=B\sqcup\{\infty\}$, with
$\infty$ absorbing for both operations. This distributive bimonoid is not
a semiring with absorbing additive zero. There are two ideal conventions.
If an ideal means a nonempty additively closed, multiplication-absorbing
subset, its ideals are $\{\infty\}$ and $J\cup\{\infty\}$ for ideals $J$
of $B$. Its prime spectrum adds one generic open point $\{\infty\}$ below
$\Spec B$. If ideals are instead required to contain the additive identity
$0_B$, its prime spectrum is canonically $\Spec B$.

With the usual unital semimodule axioms, including $a0_M=0_M$ and
$0_Bm=0_M$, the only semimodule is the zero module. Its unit group is
$B^\times$; the extra element $\infty$ is fixed by every unit.
\end{proposition}
\begin{proof}
The equality $0_B\infty=\infty\ne0_B$ violates zero absorption. Every
nonempty multiplication-absorbing subset contains $\infty$. If it also
contains $x\in B$, it contains $0_Bx=0_B$, and its intersection with $B$
is an ideal of $B$. Conversely both displayed types are closed as
required. The ideal $\{\infty\}$ is prime because a product of two elements
of $B$ stays in $B$. The ideal $J\cup\{\infty\}$ is prime exactly when
$J$ is prime. Every nonempty distinguished open contains $\{\infty\}$,
and $D(0_B)=\{\{\infty\}\}$, proving the added-point assertion. Requiring
$0_B$ deletes exactly this ideal and gives the second spectrum.

For a module, associativity gives
\[
 \infty m=(\infty0_B)m=\infty(0_Bm)=\infty0_M=0_M.
\]
Since $\infty+1=\infty$, distributivity gives $0_M+m=0_M$, so $m=0_M$.
A product equal to $1$ cannot involve $\infty$, which proves the unit
assertion; absorption proves that every unit fixes $\infty$.
\end{proof}

For inverse $B$, this is a \emph{zeroless inverse semiring} in the
terminology of \cite[\S3.1]{FGJM}; for rings their Definition~3.8 and
Proposition~3.9 give this construction. The extra prime in the first
convention is therefore valid, but must not be transferred to the second
one. For $B=\Z_0$, the unit-fixed elements are $\tau,e,\infty$, ordered
$\tau<e<\infty$ by addition. Their multiplication is not the meet
operation of that chain: $\tau\infty=\infty$.

\begin{proposition}\label{prop:monoid}
For a commutative monoid $M$, the monoid semiring $R_0[M]$ is the inverse semiring of the pair
\[
 Q=\B[M]=\mathcal P_{\mathrm{fin}}(M),\qquad F(U)=R^U,
\]
where addition in $Q$ is union, multiplication is setwise product, transition maps extend by zero, and multiplication on the group fibers is convolution. In particular a zero ring coefficient at a monomial in $U$ remains distinguished from absence of that monomial.
\end{proposition}
\begin{proof}
An element of $R_0[M]$ has finitely many coefficients different from the new zero. Their set $U$ is its support, and each such coefficient belongs to $R$, including its old zero. Sums have support union because a sum of old coefficients is never the new zero. Products have support $UV$ for the same reason, even when their coefficient sums cancel in $R$. The coefficient formula is precisely convolution, and extending coefficient lists by old zero supplies the transition maps. Convolution is bilinear, natural, associative and unital; hence it is the lax monoidal structure of \cite[Theorem~8.5]{CFMM}. The resulting operations coincide coefficient by coefficient with those of the monoid semiring. For $R=\Z$ and $M=\N^n$ this is the free commutative inverse semiring already described in \cite[Proposition~5.2]{FGJM}.
\end{proof}

\section{Distributive-lattice extensions}\label{sec:latticeextension}
This section replaces the chain of idempotents by a bounded distributive lattice and records the resulting geometry.

Let $R$ again be commutative and let $L$ be a nontrivial bounded distributive lattice. Regard $L$ as a semiring with join and meet. Define
\[
 G_L(R)=\{(0_R,\lambda):\lambda\in L\}\cup\{(r,1):r\in R\}
 \subseteq R\times L,
\]
with coordinatewise operations, and put $z_\lambda=(0_R,\lambda)$. Its zero is $z_0$, its unit is $(1_R,1)$, and its top additive idempotent is $e=z_1$. The pair in Theorem~\ref{thm:general} is $Q=L$, with zero group below $1$ and group $R$ at $1$. For $L=C_{k+1}$ this is $R^{[k]}$.

Write $\Spec_{\mathrm{lat}}L$ for the prime lattice ideals, with basic opens $D_L(\lambda)=\{\mathfrak a:\lambda\notin\mathfrak a\}$. A lattice ideal contains $0$, is downward closed and is closed under finite joins.

\begin{proposition}\label{prop:latticeextension}
The projection $G_L(R)\to R$ is its ring reflection. Semiring maps $G_L(R)\to\B$ are precisely bounded lattice maps $L\to\B$, extended by sending every supported ring element to $1$.

The ideals of $G_L(R)$ are the pure ideals $\{z_\lambda:\lambda\in\mathfrak a\}$ for proper lattice ideals $\mathfrak a$, and
\[
 \{z_\lambda:\lambda\in L\}\cup\{(r,1):r\in J\}
\]
for ring ideals $J$. All these ideals are subtractive. They are prime precisely when $\mathfrak a$ or $J$, respectively, is prime. Thus the spectrum consists of the open stratum $\Spec_{\mathrm{lat}}L$ and the closed stratum $\Spec R$, with every lattice prime below every ring prime. Its basic opens are
\[
 D(z_\lambda)=D_L(\lambda),\qquad
 D((r,1))=\Spec_{\mathrm{lat}}L\cup D_R(r).
\]
Localization at $z_\lambda$ is the semiring $\downarrow\lambda$, with unit $\lambda$. At a lattice prime $\mathfrak a$ the stalk is the lattice-semiring localization $L_{\mathfrak a}$; at a ring prime $\mathfrak p$ it is $G_L(R_{\mathfrak p})$. If $R\ne0$ and both dimensions are finite, then
\[
 \dim G_L(R)=\dim\Spec_{\mathrm{lat}}L+\dim R+1.
\]
For $R=0$, the spectrum and dimension are simply those of $L$.
\end{proposition}
\begin{proof}
The displayed subset is closed under addition and multiplication: a sum with a supported ring element has top support, and a product with a proper-support element has ring coordinate zero. The semiring laws follow from those of $R\times L$. Every $z_\lambda$ is additively idempotent, so a map to a ring kills it. Restriction to the supported ring then proves the reflection property. A map to $\B$ sends $e=1+(-1)$ to $1$; from $er=e$ it follows that every supported $r$ maps to $1$. Its restriction to the $z_\lambda$ is a bounded lattice map, and conversely each such map gives the stated semiring map.

An ideal consisting only of proper-support elements is downward closed by multiplication with $z_\mu$ and join closed by addition, hence is a proper lattice ideal. If it contains a supported ring element $r$, multiplication by $-1$ and addition give $e$. Multiplication by $z_\lambda$ then gives every lattice element, and the supported ring elements in the ideal form a ring ideal. These arguments also prove that every listed set is an ideal. Subtractivity in the pure case follows from downward closure; in the other case, if all relevant elements have top support use ring subtraction, while any element of lower support already lies in the ideal. Primality reduces respectively to meet-primality in $L$ or primality in $R$, because the proper-support elements multiply by meet and the supported ring elements multiply in $R$. This proves the point classification. The formulas for basic opens follow directly and give the topology.

For any multiplicative idempotent $a$ in a commutative semiring, multiplication $x\mapsto ax$ identifies localization at $a$ with $aS$, whose unit is $a$: every invertible idempotent becomes the unit, and every map inverting $a$ factors uniquely through this image. For $a=z_\lambda$ the image consists of $z_\mu$ with $\mu\le\lambda$. At a pure lattice prime, inverting $e$ first gives $L$ and then inverts the complement of that prime. At a ring prime, the denominators are exactly the supported ring elements outside the prime, so the old-denominator argument of Proposition~\ref{prop:localization} gives $G_L(R_{\mathfrak p})$. Finally every prime chain is a lattice-prime chain followed by a ring-prime chain. When both strata are nonempty, maximal lengths therefore add with the one intervening strict inclusion. If $R=0$, there is no ring-prime stratum.
\end{proof}

\begin{proposition}\label{prop:latticebranches}
Every bounded lattice map $\epsilon:L\to\B$ gives a surjective semiring map $G_L(R)\to R_0$, retaining the supported ring and sending $z_\lambda$ to the new zero or the old zero according as $\epsilon(\lambda)=0$ or $1$. Its spectrum map selects the lattice prime $\ker\epsilon$ and all ordinary ring primes.

If $L$ is finite and $P$ is its poset of nonzero join-irreducibles, then $L\simeq\Id(P)$ and $\Spec_{\mathrm{lat}}L\simeq P$ with open sets the order ideals. The lattice maps to $\B$ are exactly the coordinate tests $I\mapsto 1_{p\in I}$.
\end{proposition}
\begin{proof}
The formula for the first map respects addition by join and multiplication by meet, and respects mixed operations because every proper-support element has ring coordinate zero. It is onto, since it contains both zeros and every ring element. Its fibers identify precisely the lattice elements with the same $\epsilon$-value and leave distinct ring amplitudes distinct. Pulling back the primes of $R_0$ proves the spectrum assertion.

For the finite statement, a join-irreducible in a distributive lattice is join-prime: if $p\le a\vee b$, then $p=(p\wedge a)\vee(p\wedge b)$, so $p\le a$ or $p\le b$. Every element is the join of the join-irreducibles below it. Thus $x\mapsto\{p:p\le x\}$ identifies $L$ with $\Id(P)$, as in Lemma~\ref{lem:lattice}. Each coordinate test preserves join and meet. Conversely the $1$-fiber of a bounded lattice map is a prime filter; in a finite lattice it has a least element $p$, and join-primality of that filter makes $p$ join-irreducible. The map is exactly the coordinate test for $p$. Its kernel is a prime ideal, and every prime ideal has such a characteristic map. The basic open indexed by $x$ corresponds to its order ideal of join-irreducibles, proving the topology.
\end{proof}

A finite chain can also be realized by the ideals $(p^a)$ of $\Z/p^n\Z$ under sum and intersection. This is a lattice realization, not an identification with the ring's ideal multiplication: the latter sends exponents $a,b$ to $\min(a+b,n)$, whereas intersection uses $\max(a,b)$. The same distinction applies to the ideals $(\varepsilon^a)$ of $k[\varepsilon]/(\varepsilon^n)$.

\begin{proposition}[The boundary of the associative framework]
\label{prop:nonassociative-boundary}
Suppose instead that $A$ is a unital, possibly nonassociative algebra
whose addition is an abelian group, and form the same set $G_L(A)$ with
coordinatewise operations. In a bracketed monomial, the amplitude is
the corresponding monomial in $A$, and the support is the meet of its
variable supports, independently of bracketing. Thus a multiplicative
identity with the same variables on both sides has exactly its ordinary
amplitude obstruction. The reflection of this object into associative
unital rings is $A/I_{\mathrm{ass}}$, where $I_{\mathrm{ass}}$ is the ideal
generated by all associators $(ab)c-a(bc)$; it is not generally $A$.
\end{proposition}
\begin{proof}
Induction on a bracketed product proves the two coordinate formulas;
associativity and commutativity of meet remove the bracketing and order
from the support coordinate. The two amplitudes therefore agree exactly
when the two products agree, once the variable lists match. A map into
an associative ring kills every additive idempotent and hence factors
through the amplitude algebra $A$. Its kernel contains every associator
and the ideal they generate. Conversely $A/I_{\mathrm{ass}}$ is
associative, because each associator vanishes there, and its amplitude
map gives the required factorization. Uniqueness follows from the old
copy of $A$. If $A$ is nonassociative, $G_L(A)$ is outside the ordinary
semiring category used in the preceding classification theorems.
\end{proof}

\section{Projectives over finite distributive support extensions}
\label{sec:finite-lattice-projectives}
We prove Theorem~\ref{thm:finite-lattice-projectives}. Throughout this
section $C$ is a finite poset, $L=\Id(C)$, and
$T=G_L(R)$ for a unital ring $R$, possibly noncommutative. The Boolean
projectivity input is the finite case of Horn and Kimura's theorem,
proved in Lemma~\ref{lem:lattice}; see also Macpherson
\cite[Theorem~3.9]{Macpherson}. The group-diagram argument below is
ordinary projective-module theory on a finite poset.

\begin{lemma}[module data]\label{lem:lattice-module-data}
A $T$-semimodule is equivalently an $L$-semimodule $N$ together with a
covariant diagram $G:N\to R\text{-}\Mod$. Morphisms consist of an
$L$-linear support map and compatible linear maps on groups. The fiber
over the least element of $N$ need not be zero.
\end{lemma}
\begin{proof}
The action of $(-1_R,C)$ supplies the additive inverses in the sense
of inverse monoids. Their group identities form $N$. Every
$(0_R,l)$ acts on identities as the corresponding lattice scalar;
$(0_R,C)$ acts as their identity. The old ring elements act within
each group, making it an $R$-module. Adding a larger group identity
gives the transition maps, as in Section~\ref{sec:diagrams}.

Conversely, transport addition to the join group, let an old ring
scalar act linearly in the same fiber, and let a proper lattice scalar
give the zero in its scalar-translated fiber. These formulas satisfy
the module identities. In particular, if two proper lattice elements
join to $C$, their sum in $T$ is $(0_R,C)$, the old zero; both sides of
the distributive identity are the group zero in the join fiber.
The zero scalar gives the zero in the least fiber. The constructions
on maps are inverse because a homomorphism preserves group identities
and additive inverses.
\end{proof}

\begin{lemma}[the support classification]\label{lem:lattice-support-projectives}
The finite projective $L$-semimodules are exactly the colored order-ideal
lattices in the first part of
Theorem~\ref{thm:finite-lattice-projectives}.
\end{lemma}
\begin{proof}
Write $N$ as a retract of $L^n$, with section $i$ and retraction $r$.
Since $L$ is a finite projective Boolean semimodule, so is $N$;
Lemma~\ref{lem:lattice} gives $N=\Id(P)$ for a finite poset $P$.
Let $d_l$ denote multiplication by $l\in L$. If $I\le J$ and
$d_lJ=J$, then $i(I)\le i(J)$ has every coordinate contained in $l$,
so $d_lI=I$. Thus the fixed elements are the ideals contained in a
fixed ideal $U_l\subseteq P$. The map $d_l$ sends $I$ to its greatest
$l$-fixed subideal, hence $d_lI=I\cap U_l$. The action laws give
\[
 U_{l\cap m}=U_l\cap U_m,\quad U_{l\cup m}=U_l\cup U_m,
 \quad U_\varnothing=\varnothing,\quad U_C=P.
\]
For each $p$, membership in $U_l$ is a bounded lattice character of
$\Id(C)$, hence evaluation at a unique $c(p)\in C$. The ideals $U_l$
make $c$ order preserving, and give the asserted scalar action.

Decompose $i(\mathord\downarrow p)$ into coordinate join-irreducibles
$(\mathord\downarrow a)e_j$ of $L^n$. Some term retracts to
$\mathord\downarrow p$. Its color is $a=c(p)$: the section of
$\mathord\downarrow p$ is fixed by $\mathord\downarrow c(p)$,
giving $a\le c(p)$, while its retraction is fixed by
$\mathord\downarrow a$, giving the reverse inequality. Fix that
coordinate $j$. For $a\le c(p)$ put
\[
 X=r((\mathord\downarrow a)e_j)
   =d_{\mathord\downarrow a}(\mathord\downarrow p).
\]
The $j$th coordinate of $i(X)$ is exactly $\mathord\downarrow a$.
If $X=Y\vee Z$, the join-primality of $a$ in $L$ puts $a$ in one
of the corresponding section coordinates, say that of $Y$.
Then $(\mathord\downarrow a)e_j\le i(Y)$ and $X\le Y$.
Therefore $X$ is a nonzero join-irreducible of $N$. Its color is exactly
$a$: it is fixed by $\mathord\downarrow a$, and every scalar fixing
it must fix its $j$th section coordinate. Thus $X=\mathord\downarrow p_a$
with $c(p_a)=a$, proving the required greatest-element condition.

Conversely, for such a colored poset define
\begin{equation}\label{eq:lattice-support-section}
 s(I)_p=\{a\le c(p):p_a\in I\},\qquad
 r(l e_p)=l\cdot\mathord\downarrow p.
\end{equation}
The greatest-element property gives $(p_b)_a=p_a$ for
$a\le b\le c(p)$, and $p_a\le q_a$ when $p\le q$.
These identities make $s(I)_p$ an ideal in $C$. Membership of $p_a$
proves that $s$ preserves zero and joins, and its exact color proves
$s(lI)=l s(I)$. Every term of $r(s(I))$ lies in $I$, while $p\in I$
implies $c(p)\in s(I)_p$ and hence $p\in r(s(I))$. Thus $rs=1_N$.
\end{proof}

\begin{lemma}[projective diagrams on a finite poset]
\label{lem:projective-poset-diagrams}
Let $X$ be a finite poset. Every finitely generated projective object
$G$ of $\operatorname{Fun}(X,R\text{-}\Mod)$ is isomorphic to
\[
 G_x=\bigoplus_{a\le x}Q_a
\]
with inclusion maps and finitely generated projective $R$-modules $Q_a$.
Their isomorphism classes are recovered by
\begin{equation}\label{eq:diagram-birth-quotient}
 Q_a\simeq G_a\big/\sum_{b<a}\im(G_b\longrightarrow G_a).
\end{equation}
\end{lemma}
\begin{proof}
Let $E_a(Q)$ be the diagram equal to $Q$ at $x\ge a$ and zero elsewhere,
with identity maps wherever both values are $Q$. The diagrams $E_a(R)$
are projective generators: maps from them are elements of the group
at $a$. Therefore $G$ is a retract of a finite sum of these diagrams.
The quotient functor on the right of
\eqref{eq:diagram-birth-quotient} sends $E_b(R)$ to $R$ when $a=b$
and to zero otherwise. It preserves direct sums and retracts, so its
value $Q_a$ on $G$ is finitely generated projective over $R$.
Choose a linear section $Q_a\to G_a$ of each quotient. Transporting
these sections gives a natural map $F=\bigoplus_aE_a(Q_a)\to G$.
It is pointwise surjective by induction on $X$, and splits because
$G$ is projective. Applying each quotient functor to the splitting
shows that its kernel has every such quotient zero. A second induction
on $X$ makes the kernel zero. This proves both assertions.
\end{proof}

\begin{proof}[Proof of Theorem~\ref{thm:finite-lattice-projectives}]
Let $M$ be a finitely generated projective $T$-semimodule, with support
$N$ and group diagram $G$. The support is projective over $L$.
Indeed, lift the support map $M\to N$ across a surjection of $L$-modules,
regarding both target and domain as zero-group $T$-modules, and restrict
the lift to group identities. It is finite because it is finitely
generated over the finite semiring $L$. Lemma~\ref{lem:lattice-support-projectives}
gives $N=\Id(P)$ with the stated coloring.

The diagram $G$ is projective in $\operatorname{Fun}(N,R\text{-}\Mod)$.
Any pointwise surjection $H\to G$ defines, by
Lemma~\ref{lem:lattice-module-data}, a surjective $T$-module map with
support identity. A section supplied by projectivity of $M$ also has
support identity, so is a diagram section. Moreover a finite generating
set of $M$ generates $G$ as a diagram: all group contributions are
linear combinations of transported generators, since proper lattice
scalars contribute only group zeros. Lemma~\ref{lem:projective-poset-diagrams}
therefore gives labels $Q_A$ for all $A\in N$.

We show that a nonzero label requires admissibility. Choose a free
retraction $M\longrightarrow T^n\longrightarrow M$, with support
maps $i$ and $r$, and put $A_j=r(Ce_j)$. The group fiber of $T^n$
at $i(A)$ has free coordinates exactly at indices with $i(A)_j=C$.
Each such coordinate satisfies $A_j\subseteq A$, and its retraction
comes from the group $G_{A_j}$. If every active coordinate had
$A_j\subsetneq A$, the identity composite on $G_A$ would induce zero
on the quotient \eqref{eq:diagram-birth-quotient}, forcing $Q_A=0$.
Consequently $Q_A\ne0$ gives a coordinate $j$ with
$i(A)_j=C$ and $A_j=A$.

Set $\psi(I)=i(I)_j$ and $\theta(l)=lA$. They satisfy
\[
 \psi\theta=1_L,\qquad \theta\psi(I)\subseteq I.
\]
For $a\in C$, the element $X=\theta(\mathord\downarrow a)$ has
$\psi(X)=\mathord\downarrow a$. If $X=Y\vee Z$, membership of $a$
in one of $\psi(Y),\psi(Z)$ forces $X\le Y$ or $X\le Z$.
Thus $X$ is nonzero join-irreducible. It is fixed by
$\mathord\downarrow a$, while any scalar fixing $X$ must fix
$\psi(X)$; its color is therefore exactly $a$.
This is \eqref{eq:admissible-ideal}.

Conversely, start from the data in the statement. For each admissible
$A$, let $p_a$ denote its color-cut maximum and define
\begin{equation}\label{eq:admissible-functional}
 \rho_A(I)=\{a\in C:p_a\in I\}.
\end{equation}
It is a zero-preserving $L$-linear join map. It has
$\rho_A(I)=C$ exactly when $A\subseteq I$, and
$\rho_A(I)A\subseteq I$. Start with the support retraction
\eqref{eq:lattice-support-section}, giving every one of those
coordinates zero amplitude. For each nonzero $Q_A$, choose split
linear maps $Q_A\to R^{m_A}\to Q_A$ and append $m_A$ coordinates.
Their support section is $\rho_A$, and each corresponding retraction
column sends full support to $A$. On groups use the chosen module
section and retraction for the $Q_A$ summand. Outside $I\supseteq A$
the corresponding coordinate has proper support and zero amplitude,
so these formulas define $T$-linear maps.

The support composite is the identity: the original coordinates
already recover $I$, and every appended contribution lies in
$\rho_A(I)A\subseteq I$. The group composite is the identity by
the chosen module splittings. Thus the constructed module is a
retract of a finite free one. Finally $N$ recovers $P$ and its coloring,
and \eqref{eq:diagram-birth-quotient} recovers every label. This proves
the isomorphism classification. Diagram isomorphisms may include
upper-triangular mixing of summands; only their isomorphism classes
are being classified.
\end{proof}

\begin{corollary}[Grothendieck groups and strong projectives]
\label{cor:lattice-projective-Kzero}
Assume $C\ne\varnothing$. For data in
Theorem~\ref{thm:finite-lattice-projectives}, join two vertices of $P$
when they are comparable, and join every pair of vertices belonging
to an admissible $A$ with $Q_A\ne0$. Let $\mathcal D_{C,R}$ be the set
of isomorphism classes of nonempty data connected in this graph.
Then
\[
 V(G_{\Id(C)}(R))\simeq\N^{(\mathcal D_{C,R})},\qquad
 K_0(G_{\Id(C)}(R))\simeq\Z^{(\mathcal D_{C,R})}.
\]
If $R\ne0$, every finitely generated strong projective is free, and
$SV(G_{\Id(C)}(R))\simeq\N$, $SK_0(G_{\Id(C)}(R))\simeq\Z$.
If $R=0$, this monoid and group are instead
$\N^{\pi_0(C)}$ and $\Z^{\pi_0(C)}$, respectively, where $\pi_0(C)$
is the set of comparability components of $C$. When $C=\varnothing$
all assertions reduce to the corresponding ordinary ring invariants.
\end{corollary}
\begin{proof}
A direct sum takes the disjoint union of the colored posets, retains
the labels on ideals wholly contained in a summand, and gives zero
labels to every mixed admissible ideal. Since $C$ is nonempty,
every admissible ideal is nonempty. A decomposition therefore amounts
exactly to a partition into comparability components which does not
separate any nonzero label. The displayed graph gives the unique
finest partition. Restricting data to any component retains the
color-cut condition, because its required greatest elements are
comparable to their original vertices; all its nonzero labels remain
admissible. This proves the free-monoid assertion and its group
completion.

The free rank-one semimodule has $P=C$, identity coloring, and the
single label $R$ on $A=C$. If $R\ne0$, this label connects all of $C$.
A direct summand of its $n$-fold direct sum is therefore a union of
whole copies of these data, hence free. If $R=0$, the labels vanish
and the indecomposable pieces of the free rank-one object are the
comparability components of $C$. They remain distinguishable even
if their uncolored posets are isomorphic, because colors are fixed.
An arbitrary finite multiplicity for each such component is a direct
summand of a sufficiently large free module, giving the stated
monoid and group. The empty-$C$ case is $T=R$.
\end{proof}


\section{Finite group monoid semirings}
\label{sec:finite-group-projectives}

We prove Theorem~\ref{thm:finite-group-projectives}. Put $S=R_0$ and
$T=S[\Gamma]$. A left $T$-semimodule is exactly an $S$-semimodule
with a $\Gamma$-action by $S$-linear automorphisms. Restriction gives
one direction; in the other, set
\[
 \left(\sum_{g\in\Gamma}a_g g\right)m
   =\sum_{g\in\Gamma}a_g(gm).
\]
The old zero coefficient remains an occupied coefficient, so its action
is a fiber zero rather than an omitted summand. The distributive and
unit laws follow from those for $S$ and the commuting group action.

The support case is related to the study of group representations over
idempotent semifields. Jun, Mincheva and Tolliver
\cite[arXiv version~3, \S1.1 and Definition~3.1]{JMT} classify actions on free underlying
semifield modules; their Corollary~4.9 and Lemma~4.11 give further
Boolean-module results. Our modules need not be free over $\B$ or
$R_0$. Macpherson's $\mathbb F_1$-domain hypothesis
\cite[Definition~4.9, p.~14]{Macpherson} excludes nontrivial
multiplicative torsion. Stabilizer group rings also occur in the
classical splitting theory for modules over categories whose
endomorphisms are isomorphisms, recalled by Fiore, L\"uck and Sauer
\cite[pp.~306--307, splitting functors and Theorem~1.5]{FLS}. The proof here
constructs the required equivariant splittings directly over an
arbitrary unital coefficient ring.

\begin{proof}[Proof of Theorem~\ref{thm:finite-group-projectives}]
As an $S$-semimodule, $T$ is free on the finite set $\Gamma$.
Consequently a finite projective $T$-semimodule $M$ is finite projective
over $S$. By Lemma~\ref{lem:noncomm}, its support is $\Id(P)$ for a
finite poset $P$, and its group transitions are injective. Identify all
its groups with submodules of the top group. The classification gives
\[
 M_I+M_J=M_{I\cup J},\qquad M_I\cap M_J=M_{I\cap J}.
\]
The group action induces order automorphisms of $P$. Write
\[
 I_p=\downarrow p,\qquad I_p^-=I_p\setminus\{p\},\qquad
 V_p=M_{I_p}/M_{I_p^-}.
\]
This quotient is an $R[\Gamma_p]$-module. Its underlying $R$-module
is finite projective; we must prove the stronger stabilizer assertion.

For one $p$ in each orbit choose a finite free $R[\Gamma_p]$-module
$W_p$ surjecting onto $M_{I_p}$. Such a cover exists because $M_{I_p}$
is finitely generated over $R$: send a group-ring basis to an
$R$-generating family and extend equivariantly. Transport these modules
and maps along their orbits. The diagram
\[
 N_I=\bigoplus_{q\in I}W_q
\]
with its transported action is a $T$-semimodule with the same support.
The chosen maps and the group transitions define a $T$-map $f:N\to M$
inducing the identity on support. It is surjective on every group:
$M_I$ is the sum of the principal groups $M_{I_q}$ for $q\in I$,
and each of those is covered. Projectivity of $M$ supplies a section
with identity support. Passing to the principal birth quotients gives
\[
 V_p\longrightarrow W_p\longrightarrow V_p
\]
with identity composite, linearly over $R[\Gamma_p]$. Thus $V_p$ is
finite projective over that group ring. This argument does not require
$N$ to be projective.

Choose an $R[\Gamma_p]$-linear section of $M_{I_p}\to V_p$ for each
orbit representative, and transport it along $\Gamma$. Stabilizer
linearity makes the transported section independent of the chosen
coset representative. The sum of these sections identifies
$\bigoplus_{p\in I}V_p$ with $M_I$. Indeed, for a maximal $p\in I$
and $J=I\setminus\{p\}$, the two displayed identities give
\[
 M_I=M_J+M_{I_p},\qquad
 M_J\cap M_{I_p}=M_{I_p^-}.
\]
The selected section therefore yields $M_I=M_J\oplus V_p$, and finite
induction proves the claim. The identifications commute with all
transitions and with $\Gamma$. This proves necessity and the asserted
form of the module.

For sufficiency, start with the stated data. For each orbit representative
$p_a$, let $H_a=\Gamma_{p_a}$ and choose a split pair
\[
 V_{p_a}\xrightarrow{i_a}R[H_a]^{m_a}
 \xrightarrow{r_a}V_{p_a},\qquad r_ai_a=1,
\]
with $m_a\ge1$. For a zero label take $m_a=1$ and zero maps; the
coordinate still records the support point. Let $n=\sum_a m_a$.
The underlying $S$-free coordinates of $T^n$ are $(g,a,j)$ for
$g\in\Gamma$ and $1\le j\le m_a$. Define support maps by
\[
 \sigma(I)=\{(g,a,j):gp_a\in I\},\qquad
 \rho(U)=\bigcup_{(g,a,j)\in U}\downarrow(gp_a).
\]
They preserve unions and $\Gamma$-actions, hence are
$\mathcal P(\Gamma)$-linear, and $\rho\sigma=1$.

Induce the split pairs from $H_a$ to $\Gamma$. The resulting top group
maps are a split pair between
\[
 \bigoplus_a R[\Gamma]\otimes_{R[H_a]}V_{p_a}
 \quad\text{and}\quad
 \bigoplus_a R[\Gamma]^{m_a}.
\]
The first module is the direct sum of all the transported labels.
The section sends the label at $gp_a$ into coordinates $(gh,a,j)$
with $h\in H_a$; all have attached point $ghp_a=gp_a$.
It therefore takes the group over $I$ into the free group over
$\sigma(I)$. The retraction sends coordinate $(g,a,j)$ into the label
at $gp_a$, and hence into the group over $\downarrow(gp_a)$.
Both maps respect transitions and are $R$-linear and
$\Gamma$-equivariant. Together with $\sigma,\rho$, they give
$T$-linear maps
\[
 M\longrightarrow T^n\longrightarrow M
\]
whose composite is the identity on groups and supports. Thus $M$ is
finite projective. If $P$ is empty, the zero-fold free module suffices.

Finally any $T$-isomorphism induces a $\Gamma$-equivariant isomorphism
of the support posets and of all the intrinsic quotients $V_p$ with
their stabilizer actions. Conversely such data identify the split
diagrams and thus the $T$-semimodules. This proves precisely the
claimed isomorphism-class classification; it does not assert that
all homomorphisms are diagonal on the labels.
\end{proof}

For example, take $R=\Z$, $\Gamma\ne1$ finite, and $P$ a single
fixed point. The trivial label $\Z$ fails the stabilizer condition.
Indeed an equivariant section of the augmentation
$\Z[\Gamma]\to\Z$ would send $1$ to an invariant element
$c\sum_{g\in\Gamma}g$, whose augmentation is $c|\Gamma|\ne1$.
The regular label $\Z[\Gamma]$ always satisfies the condition.
When $R=0$, all labels vanish and the explicit support retraction
above proves the asserted classification over $\B[\Gamma]$.

\begin{corollary}\label{cor:finite-group-kzero}
Let $\mathcal C_{R,\Gamma}$ be the isomorphism classes of the nonempty
data in Theorem~\ref{thm:finite-group-projectives} for which $\Gamma$
acts transitively on the connected components of the comparability graph.
Then
\[
 V(R_0[\Gamma])\simeq\N^{(\mathcal C_{R,\Gamma})},\qquad
 K_0(R_0[\Gamma])\simeq\Z^{(\mathcal C_{R,\Gamma})}.
\]
The labels retain their actual projective $R[\Gamma_p]$-module
isomorphism classes. Every finitely generated strongly projective
$R_0[\Gamma]$-semimodule is free, and
\[
 SV(R_0[\Gamma])\simeq\N,\qquad SK_0(R_0[\Gamma])\simeq\Z.
\]
These assertions include $R=0$.
\end{corollary}
\begin{proof}
Direct sum is disjoint union of the decorated $\Gamma$-posets.
The comparability graph has a unique set of connected components,
which $\Gamma$ permutes. The unions belonging to individual
$\Gamma$-orbits of components give a unique decomposition into
nonempty $\Gamma$-invariant pieces. Such a piece admits no further
disjoint invariant decomposition. The classification therefore gives
the free commutative monoid, and its group completion is the displayed
free abelian group. No stabilization of individual ring labels is made.

The decorated poset of $T^n$ is the disjoint union of $n$ regular
$\Gamma$-orbits, with the discrete order, trivial stabilizers, and
label $R$. If $M\oplus N\simeq T^n$, the poset of $M$ is an invariant
subunion of these regular orbits with their unchanged labels. It is
therefore the poset of $T^k$ for some $k$. Conversely free modules
are strongly projective. Distinct free ranks have distinct support
cardinalities $2^{|\Gamma|k}$, also when the ring is zero. This proves
the strong-projective formulas.
\end{proof}


\section{Adjunctions, finite products, and submodules}
\label{sec:furtherfoundations}
\begin{proposition}[The unit-fiber adjunction]\label{prop:unitfiberadjunction}
For a unital inverse semiring $T$, let $U_1(T)$ be its additive group fiber
containing $1$, with group zero $e_1=1+(-1)$. This is a unital ring, and
\[
 \operatorname{Hom}_{\mathrm{InvSemi}}(R_0,T)
 \cong\operatorname{Hom}_{\mathrm{Ring}}(R,U_1(T))
\]
naturally. Thus zero adjunction is fully faithful and left adjoint to
$U_1$. This is the composite of the zero-adjunction/zeroless adjunction
and the ring-of-scalars adjunction stated in \cite[Definition 3.6,
Section 3.1, Proposition 4.7 and Remark 4.14]{FGJM}; it differs from their
ordinary ring reflection.
\end{proposition}
\begin{proof}
The idempotent map is multiplicative and $e_1^2=e_1$, so the fiber is
closed under multiplication and contains $1$. Its group zero absorbs
multiplication within the fiber, since $e_1a=a+(-a)=e_1$. Thus it is a
ring. A semiring map preserves unique additive inverses, sends the old
zero of $R_0$ to $e_1$, and restricts to a ring map into this fiber. Conversely
such a ring map extends uniquely by sending the new zero to the global
zero of $T$. The constructions are inverse and natural. The adjunction
unit $R\to U_1(R_0)$ is an isomorphism, proving full faithfulness.
\end{proof}

\begin{corollary}\label{cor:no-sections}
There is no unital semiring map from any ring to $R_0$ or to $\B$.
In particular, the ring reflection $R_0\to R$ has no semiring section.
If $R\ne0$, there is no unital semiring map $\B\to R_0$ either.
\end{corollary}
\begin{proof}
The images of $1$ and $-1$ under a ring-to-semiring map would sum to the
new zero. This is impossible in either target, where a sum is the new
zero only when both summands are new zero, whereas $1$ is not new zero.
A map $\B\to R_0$ would give $1_R+1_R=1_R$, hence $1_R=0_R$.
\end{proof}

There is nevertheless a continuous map
$\Spec\Z_0\to\Spec\Z$ which preserves every ordinary prime and sends the
new generic point to $(0)$. For $n\ne0$, the inverse image of
$D_{\Z}(n)$ is $D_{\Z_0}(n)$; for $n=0$ it is empty. These basic opens
prove continuity. This topological collapse is not induced by a unital
map $\Z\to\Z_0$, since Corollary~\ref{cor:no-sections} excludes such a map.

\begin{proposition}\label{prop:finite-product-reflections}
For a finite set $I$, the universal ring reflection of $(R_0)^I$ is $R^I$,
with zero fiber the Boolean cube $\B^I$. Localization of a finite semiring
product at one tuple is the product of its coordinate localizations.
For a field $K$, this turns nonzero-amplitude coordinates into unchanged
$K_0$ factors, old-zero coordinates into $\B$, and new-zero coordinates
into terminal factors. Identifying an old-zero mask with a tuple of
nonzero amplitudes on the same selected coordinates Booleanizes exactly
those coordinates.
\end{proposition}
\begin{proof}
Every additive idempotent is killed by a ring-valued map. Adding the full
old-zero mask to a tuple replaces its absent coordinates by old zero,
without changing that map's value. Thus the map factors uniquely through
its tuple of amplitudes. The zero fiber is exactly the cube of old/new
zero choices, with union and intersection operations.

Coordinate fractions admit a common denominator exponent, the maximum of
finitely many exponents; equality witnesses likewise admit a common
maximum. This proves the finite-product localization assertion. Inverting
an old field unit has no effect, inverting the old zero identifies all old
scalars with $1$, and inverting the new zero gives the terminal semiring.
For the mask congruence, multiply the proposed relation by each coordinate
projector. At a selected coordinate it says $e=a$ with $a\ne0$, hence makes
$e$ invertible and collapses the old field to its Boolean value. The
coordinate decomposition proves universality and preserves every
unselected coordinate.
\end{proof}

\begin{proposition}\label{prop:diagram-consequences}
The additive-idempotent functor $M\mapsto eM$ is both left and right
adjoint to restriction of $\B$-semimodules along $R_0\to\B$.
The least group fiber $M_{\bot}$ is right adjoint to restriction of
$R$-modules along $R_0\to R$. The cancellative $R_0$-semimodules are
exactly these ordinary $R$-modules.

A subsemimodule is specified by a zero-preserving join-subsemilattice
$L'\subseteq L_M$ and submodules $N_l\subseteq M_l$ for $l\in L'$,
compatible with every transition between its labels. If $m_1,\ldots,m_n$
generate $M$, its support lattice is generated by $l_i=em_i$ and has at most
$2^n$ elements; its fiber over $l$ is generated by the transported $m_i$
with $l_i\le l$.
\end{proposition}
\begin{proof}
A map to an idempotent module satisfies $f(m)=f(em)$, proving the left
adjunction. A map from an idempotent module lands in the additive
idempotents, proving the right one. A map from an ordinary module sends
inverse pairs to pairs summing to the global zero, so lands in $M_{\bot}$;
conversely a linear map to that fiber is an $R_0$-map. In a cancellative
module the identity $m+em=m$ gives $em=0$, so its support has one point.
Its group fiber is the required ordinary module, and the converse is
immediate.

A subsemimodule contains its supports, is closed under $-1$, and is closed
under adding a larger support. It therefore yields exactly the indicated
submodules and transitions. Conversely those data are closed under the
sum formula at a join and under scalar action. Every generated element
is a sum of old scalar multiples of the $m_i$, after omitting new-zero
terms. Its support is the join of their $l_i$. Within a fixed fiber,
transport the terms to that join to express it as the stated linear
combination. The empty sum supplies the global zero; no assertion that
the entire least group vanishes is used.
\end{proof}

\section{The multiplicative skeleton and its limits}
\label{sec:multiplicativeskeleton}
\begin{proposition}\label{prop:contracted-presentation}
Let $M$ be the multiplicative monoid of $R_0$. Its contracted monoid semiring
\[
 \mathbb N[M]/([\tau]=0)
\]
modulo the relations $[a]+[b]=[a+b]$ for $a,b\in R$ is $R_0$.
If one instead starts with the ordinary uncontracted monoid semiring and
imposes only those relations and $[x]+[\tau]=[x]$, the assertion fails.
\end{proposition}
\begin{proof}
After contraction, every nonempty sum of old generators reduces to one
old-ring generator. Evaluation distinguishes all such generators and the
empty sum, giving the isomorphism. Without contraction, evaluate each
$[x]$ in the inner copy of $R_0$ in $(R_0)_0$. All the displayed relations
hold, but $[\tau]$ maps to the inner zero and the empty sum to the distinct
outer zero. Thus those relations do not force preservation of the source
additive zero by the proposed inverse map.
\end{proof}

\begin{proposition}\label{prop:monoid-spectrum-ring}
If $R$ is a UFD with infinitely many irreducible classes, the spectrum of
the multiplicative monoid of $R_0$ has infinite dimension. If $R$ is a
domain, its ordinary contracted integral monoid ring is
\[
 \mathbb Z[M]/([\tau])\cong\mathbb Z\times\mathbb Z[R\setminus\{0\}].
\]
This ordinary ring must not be confused with the inverse-semiring monoid
algebra $\mathbb Z_0[M]$ of Proposition~\ref{prop:monoid}.
\end{proposition}
\begin{proof}
For any set $E$ of irreducible representatives, the set consisting of the
two zeros and all old nonzero elements divisible by some member of $E$
is a proper monoid ideal with multiplicatively closed complement.
Different sets give the same inclusion relation as their associated
ideals. Increasing finite sets of irreducibles give arbitrarily long
prime chains.

In the contracted integral monoid ring, let $z=[0_R]$. It is a central
idempotent and $z[a]=z$ for every old nonzero $a$. Its factor is therefore
$\mathbb Z$. The complementary factor has basis $[a]-z$ indexed by old
nonzero $a$; the domain hypothesis makes their products exactly
$([a]-z)([b]-z)=[ab]-z$. This identifies the other factor with the stated
ordinary monoid ring.
\end{proof}

For a domain, an old nonzero element has exactly its ordinary ring
factorizations, so its units and irreducibles are classical. The old
zero $e$ is reducible, since $e=e^2$, and is divisibility-prime: $e\mid xy$
means $xy\in\{e,\tau\}$, forcing a factor into that set. Thus prime in the
divisibility sense does not imply irreducible. The semiring prime spectrum
proved earlier is distinct from the much larger monoid spectrum.

\begin{proposition}[Ideal finiteness and integer arithmetic]
\label{prop:integer-arithmetic}
For a commutative unital ring $R$, extension $J\mapsto J\cup\{\tau\}$
preserves products of ideals and their finite generating sets. The
semiring $R_0$ is Noetherian, respectively has every ideal principal,
exactly when $R$ has that property.

In $\Z_0$, the cancellative multipliers are exactly the old nonzero
integers; the units are $\pm1$, and associates differ by sign. Define a
prime element to be a nonunit different from $\tau$ satisfying the usual
divisibility implication. The irreducibles are $\pm p$ for rational
primes $p$, whereas the prime elements are those elements together with
$e$. The old zero is reducible and has no factorization into irreducibles.
The Pythagorean triples are precisely the ordinary integer triples,
\[
 (a,\tau,\pm a),\qquad(\tau,b,\pm b)\quad(a,b\in\Z),
 \qquad(\tau,\tau,\tau).
\]
\end{proposition}
\begin{proof}
Finite sums of products of old elements are computed in $R$, and an old
zero term is available from every pair of ring ideals. Thus
$(I\cup\{\tau\})(J\cup\{\tau\})=IJ\cup\{\tau\}$. The ideal generated by
old elements $r_i$ is $(r_i)_R\cup\{\tau\}$; the new zero generates
$\{\tau\}$. The ideal poset of $R_0$ is the ideal poset of $R$ with one
new least point, so ascending chains stabilize in one exactly when they
do in the other. These facts prove both finiteness assertions.

Multiplication by an old nonzero integer distinguishes old inputs from
$\tau$, and is cancellative on the old inputs by integer cancellation.
Multiplication by $\tau$ is constant, while $e\cdot1=e\cdot2$.
A product equal to $1$ uses old integer factors, giving exactly the
units $\pm1$. Mutual divisibility for old nonzero integers is the ordinary
integer relation, while $e$ and $\tau$ each have a singleton associate
class, as their sets of multiples are $\{e,\tau\}$ and $\{\tau\}$.
Thus both usual definitions of associates give precisely sign changes.
Every factorization of an old nonzero integer likewise
uses old nonzero integers, so its irreducibility and primality are the
ordinary ones. The multiples of $e$ are exactly $\{e,\tau\}$; a product
lies in this pair only if one factor does, proving that $e$ is prime.
It is reducible because $e=e^2$, and a product of old irreducibles is
old nonzero.

If both entries on the left of $x^2+y^2=z^2$ are old, so is $z$, and the
equation is the ordinary integer one. If exactly one is $\tau$, the
remaining equation is $a^2=z^2$ in $\Z$, hence $z=\pm a$. If both are
$\tau$, then $z=\tau$. Every listed triple satisfies the equation.
\end{proof}

Although $R_0$ has no zero divisors with respect to $\tau$, for every
nonzero $R$ the equality $e\cdot e=e\cdot1$ prevents cancellation by all
multipliers other than $\tau$. When $R=0$, the semiring is $\B$, whose
only nonzero multiplier is $1$ and is cancellative.

If a finite nonnegative multiplicative ideal norm extends a nonconstant
ordinary ideal norm $N$, it must vanish on both $\{\tau\}$ and
$\{e,\tau\}$. Each of these ideals is unchanged by multiplication by
$J\cup\{\tau\}$ for a nonzero ordinary ideal $J$. Choosing $N(J)\ne1$
gives $w=wN(J)$ and hence $w=0$. Merely assuming that ordinary quotients
are finite does not make their cardinalities multiplicative. For
$R=\Z[2i]$ and $\mathfrak m=(2,2i)$, the basis $1,2i$ gives
$[R:\mathfrak m]=2$, while $\mathfrak m^2=(4,4i)$ has index $8$.
This supplies a concrete counterexample to that omitted hypothesis.

\section{Residue fibers and internal zero adjunction}
\label{sec:residueinternal}
\begin{proposition}\label{prop:residue-fibers}
For a surjective ring map $A\to k$ with kernel $\mathfrak m$, the induced
map $A_0\to k_0$ has old-element fibers $\widetilde a+\mathfrak m$ and
new-zero fiber $\{\tau\}$. For a finite tuple $a\in k_0^I$ of support $J$,
allow every absent coordinate of $a$ to remain absent or become the old
zero. The inverse image of this support star is the disjoint union,
indexed by $E\subseteq I\setminus J$, of
\[
 \prod_{i\in J}(\widetilde a_i+\mathfrak m)
 \times\prod_{i\in E}\mathfrak m
 \times\prod_{i\notin J\cup E}\{\tau\}.
\]
If $k$ is a field, reduction modulo $\mathfrak m^2$ makes each fixed
stratum an affine torsor under $(\mathfrak m/\mathfrak m^2)^{J\cup E}$.
Moreover, for every ideal $I\subseteq A$,
\[
 \varprojlim_n(A/I^n)_0\cong(\varprojlim_n A/I^n)_0.
\]
\end{proposition}
\begin{proof}
Old elements preserve occupied support, and the old quotient fiber is
precisely the stated coset. Partitioning the choices of activated absent
coordinates gives the product formula, with disjointness detected by
support. Quotienting each coset by $\mathfrak m^2$ gives its translation
torsor under the $k$-module $\mathfrak m/\mathfrak m^2$. This module is the
cotangent space; no extra scheme of cubes is implied by the subset
indexing. In the inverse limit, compatibility forces all coordinates of
a sequence to be old or all to be new zero. The old sequences are the
ordinary ring limit, and operations agree coordinatewise.
\end{proof}

In projective coordinates the matrix swapping two positions is invertible
and exchanges $[e:1]$ with $[1:e]$, and $[\tau:1]$ with $[1:\tau]$.
It therefore identifies the corresponding pointed residue-star sets.
For a nonzero $h\in\mathfrak m$, writing $[1:h]=[h^{-1}:1]$ requires
extension to a field containing $h^{-1}$. A nonmonomial fractional-linear
matrix, including a translation by a nonzero scalar, has no global
invertible lift over $K_0$, since an invertible Boolean matrix acts as an automorphism of the Boolean cube and must permute its atoms.
Cancelling a common factor from a pair of polynomials gives the ordinary
local extension wherever the reduced pair is unimodular; the uncancelled
pair $(e,e)$ at a common zero is not a projective point and differs from
$(\tau,\tau)$.

\begin{proposition}[Internal zero adjunction]\label{prop:internalzero}
For a commutative semiring object $A$ in a topos, the coproduct $A\amalg1$
carries zero adjunction, naturally in $A$. It has amplitude map to $A$ and
support map to $1\amalg1$, and it commutes with every geometric
inverse-image functor. For a sheaf $\mathcal A$ on a space,
\[
 (\mathcal A\amalg1)(U)
 \cong\coprod_{V\in\operatorname{Clop}(U)}\mathcal A(V).
\]
On nonempty connected $U$ this is $\mathcal A(U)_0$.
\end{proposition}
\begin{proof}
A topos is extensive, so products distribute over the indicated
coproducts. Define the operations on each summand: the old operations on
old inputs, the new point as additive identity, and as multiplicative
absorber. Every axiom follows on these summands. The two maps are
specified by sending the new point to old zero or recording its summand.
An inverse-image functor preserves finite limits and coproducts, so
preserves this construction. A section of the coproduct sheaf chooses a
locally constant summand, hence a clopen old-support subset $V$, and a
section on $V$; this description is reversible. A nonempty connected open
has just two such choices. On the empty open the sheaf has one section,
so that qualification is necessary.
\end{proof}

Every endomorphism of $A$ consequently extends while fixing the new zero.
This includes Frobenius actions on an arithmetic-site semiring sheaf.
Stalkwise, localization at its old zero gives $\B$ and at its new zero the
terminal semiring: the first localization identifies every old element
with $1$, while the second identifies $0$ and $1$. These are structural
statements about the given sheaf. They do not identify its stalks with a
fixed scalar set or supply an unstated morphism of arithmetic topoi.


\section{What remains open}
The diagram description settles the module-data question for every inverse semiring. Proposition~\ref{prop:matrix} gives exact finite certificates for projectivity and isomorphism. Theorem~\ref{thm:canonical-ideal} gives a presentation-independent reduction: finite projectivity of the support, ordinary finite projectivity over its action category, and generation by the canonical ideal $\mathcal J_L$. Corollary~\ref{cor:canonical-ideal-matrices} identifies the relevant ordinary projectives by idempotents with entries in that ideal. Proposition~\ref{prop:canonical-ideal-core} permits replacing it by its largest idempotent subideal, by a standard ordinary ideal construction applied to the action category. This replaces the existential support-section test of Theorem~\ref{thm:ringoid-projective}; Proposition~\ref{prop:ringoid-obstruction} still shows why ordinary amplitude projectivity alone is insufficient.

Theorem~\ref{thm:action-colimits} identifies that action category by directed coefficient colimits and its canonical ideal by directed coefficient images. For finite $Q$, it gives exact residual coefficient blocks, with arbitrary $L$ and arbitrary coefficient groups. Theorem~\ref{thm:finite-frame-ideal} further computes the canonical ideal from any finite projective support frame using only scalar constraint sets, for arbitrary $Q$. Thus the search over admissible functions is eliminated once that frame is given. Theorem~\ref{thm:finite-coefficient-models} further models every specified finite family of projective amplitudes and all its morphisms by exact corners in one ordinary coefficient matrix ring. No action-colimit quotient remains in those corners. Theorem~\ref{thm:uniform-coefficient-category} makes this description uniform: one category of coefficient groups and one explicit ideal are fixed by $(T,L)$ before any amplitude is chosen. Its ideal-generated finite projectives give the entire fixed-support amplitude category, including all morphisms. This removes family-dependent coefficient choices from the formulation; it retains the classification of ordinary ideal-generated projectives and of the supports themselves.

Theorem~\ref{thm:upper-unit-quotient} first removes the largest coefficient ideal vanishing above the support unit. Theorem~\ref{thm:upper-idempotent-quotient} strengthens this reduction: only the upper multiplicative-idempotent ring fibers need detect coefficient ideals. Example~\ref{ex:upper-idempotent-strict} proves that the stronger quotient can remove additional coefficients from a previously reduced pair. Both quotients preserve finite-projective isomorphism classes. Theorem~\ref{thm:upper-levitzki-quotient} removes the larger compatible ideal that is locally nilpotent in every upper idempotent ring. Its quotient commutes with finite matrices. Proposition~\ref{prop:upper-transition-radical} evaluates every coefficient group of this ideal by ordinary radicals and upward transitions in explicit two-by-two fiber rings; for finite $Q$ these are finitely many tests. Example~\ref{ex:levitzki-boundary} specifies why the residual condition concerns compatible ideals rather than the radicals of the individual rings. Theorem~\ref{thm:upper-jacobson-quotient} permits the larger compatible Jacobson quotient for isomorphism detection without extra hypotheses. Theorem~\ref{thm:exchange-coefficient-lifting} lifts projectives through every compatible coefficient ideal when all upper rings are exchange, and gives projective-class invariance for the Jacobson quotient under that hypothesis. Example~\ref{ex:jacobson-lifting-boundary} rules out unconditional object lifting. Classifying arbitrary projective $Q$-supports and ordinary projectives generated by the computed ideal for residual diagrams outside the classes below remains open.

Proposition~\ref{prop:scalar-ideal-idempotent} proves idempotence of the underlying scalar order ideal, which must be distinguished from the coefficient ideal. Propositions~\ref{prop:canonical-nonidempotent} and~\ref{prop:canonical-arbitrary-ideal} show that the canonical ideal need not be idempotent, that its powers need not stabilize, and that arbitrary ordinary ideal-generated projective modules genuinely occur in this reduction.

Theorem~\ref{thm:twolevelprojectives} resolves the problem for every two-idempotent inverse semiring, including noninjective bimodule transitions, and Theorem~\ref{thm:finite-lattice-projectives} resolves it for every finite distributive support extension $G_L(R)$, including the iterates in Theorem~\ref{thm:iteratedprojective}. Theorem~\ref{thm:finite-group-projectives} also settles $R_0[\Gamma]$ for every finite group $\Gamma$, with a non-meet support multiplication and several nonzero coefficient fibers. Theorem~\ref{thm:entire-support-projectives} evaluates constant ring fibers over every entire idempotent $Q$, by Booleanizing the projective support. Theorem~\ref{thm:finite-inverse-projectives} classifies every finite support semiring with semiperfect upper coefficient rings by support tables and multiplicities of permitted ordinary projective covers. This includes every finite inverse semiring and allows infinite coefficient groups, as Example~\ref{ex:finite-support-semiperfect} demonstrates. Theorem~\ref{thm:upper-semiperfect-amplitudes} removes finiteness of $Q$ for the amplitude classification under the same upper-semiperfect hypothesis: local idempotents in the canonical ideal and their inverse-arrow relations supply the types. Corollary~\ref{cor:upper-local-diagonal} gives the local specialization. The remaining support problem is unchanged, and no algorithm for identifying those types is asserted.

Theorem~\ref{thm:integral-support-reduction} removes all transition kernels when the support unit is greatest, replacing the coefficient diagram by ideals in its unit ring without changing finite-projective classes. Theorem~\ref{thm:local-unit-projectives} evaluates the remaining amplitudes for a local unit ring by finite multiplicities on allowed support classes. Arbitrary ideal diagrams and the underlying projective support problem remain open. The outstanding problem is an equally explicit description for arbitrary pairs $(Q,F)$ without a condition on their upper coefficient rings: their transitions, ordinary projective amplitudes and idempotent actions may all be more complicated than the classes evaluated here.

Proposition~\ref{prop:completion-projectivity-boundary} also rules out recovering ordinary projectivity merely from ideal-completion projectivity, even for a finitely presented cyclic support.

For the special idempotent coefficients $Q=\B[A;A^+]$, Macpherson \cite[Definition~4.9 and Corollary~4.13, pp.~14--16]{Macpherson} gives an intrinsic classification when the monoid pair $(A;A^+)$ is an $\mathbb F_1$-domain in his sense, using its free module of primitive elements, a nondegenerate order, and a finite-quiver presentation. This supplies a substantial known support case; it does not classify arbitrary coefficient-group diagrams $F$.

The sheaf result is complete for ordinary semimodule sheaves in the specified semiring-scheme convention. Comparing it with broader notions of modules over arbitrary blueprints is a different question: one must specify the relations in the blue module category and determine whether they preserve the fiberwise abelian-group structure. No such comparison is asserted.

\subsection{Finite presentations and a decision obstruction}
Finite operation tables and finite presentations are different inputs.
The following reduction makes the ordinary ring difficulty explicit.
\begin{proposition}\label{prop:projective-decision-reduction}
A finite presentation of any unital associative ring $R$ can be
transformed effectively into a finite presentation of the inverse
semiring $T=R\times\B$, with a named idempotent $c=(1_R,0_\B)$.
Its additive-idempotent support is exactly $\B$. The cyclic projective
$Tc$ has singleton additive support and is zero exactly when $R=0$.
Thus uniform zero testing, and hence isomorphism testing, for specified
finite projectives over finitely presented inverse semirings is at
least as hard as triviality of finitely presented unital rings, even
with fixed Boolean support and a one-by-one idempotent matrix.
\end{proposition}
\begin{proof}
Write a finite ring presentation with generators $x_i$ and polynomial
relations $f_j(x;1)=0$. In the finitary inverse-semiring signature,
put $e=1+(-1)$ and use generators $c,z_i$ with relations
\[
 c+(-c)=0,\quad z_i+(-z_i)=0,\quad c+e=1,
 \quad cz_i=z_i=z_ic,\quad f_j(z;c)=0.
\]
The last relations replace every occurrence of the ring unit by $c$.
Distributivity and uniqueness of additive inverses give
$ce=ec=0$ and $ez_i=z_ie=0$. Multiplying $c+e=1$ by $c$ gives
$c^2=c$. Also $e+e=e$, $-e=e$ and $e^2=e$.
The bottom additive group therefore contains the ring generated by
$c,z_i$, with unit $c$, and receives the presented ring $R$.
Every term has the form $a+\epsilon e$, where $a$ is in that ring
and $\epsilon\in\{0,1\}$: these forms contain the generators and
$1=c+e$, and are closed under all operations because mixed products
with $e$ vanish.

The relations hold in $R\times\B$ under
$c=(1_R,0)$ and $z_i=(x_i,0)$. Conversely, the ring map into the
bottom group extends to
$(r,\epsilon)\mapsto r+\epsilon e$ in the presented inverse semiring.
The two maps are inverse on generators and normal forms, proving
that the finite presentation is exactly $R\times\B$, also when
$R=0$. Its two additive idempotents are $(0_R,0)$ and $(0_R,1)$.

Right multiplication by $c$ is an idempotent endomorphism of the
free left module $T$. Its image is $Tc=R\times\{0\}$, a cyclic
projective with a single additive idempotent. It vanishes exactly
when $R$ does. The construction supplies the certified idempotent
matrix $[c]$; no algorithm is asked to discover its idempotence.
Comparison with the zero projective also proves the isomorphism-test
reduction.
\end{proof}

The unit convention in an ordinary undecidability input must be
checked. Bokut's original presentation convention uses free
algebras without a multiplicative identity
\cite[Section~2, p.~137]{Bokut}. The following explicit bridge is
useful independently of that source application.
\begin{lemma}\label{lem:unital-idempotent-bridge}
Let $k$ be a finite field, $A$ a finitely presented possibly nonunital
$k$-algebra, and $e\in A$ a named idempotent. Let $S=k\oplus A$ be
its unitization and put
\[
 B=(S*_k k\langle u,v\rangle)/(uv-1,\ vu-e).
\]
The free product imposes no commutation of $u,v$ with $S$.
This is an effective finite unital presentation. Moreover, $B=0$
if and only if $e=0$ in $A$. If $e\ne0$, the natural map $S\to B$
is injective.
\end{lemma}
\begin{proof}
The same generators and nonconstant relations that present $A$,
now interpreted in the free unital algebra, present $S$.
Adjoining $u,v$ and the two displayed relations is effective.
If $e=0$, then $1=(uv)^2=u(vu)v=0$ in $B$.

Suppose $e\ne0$. Let $V$ be a countable direct sum of copies of the
left regular $S$-module, regarded as a $k$-vector space. Left
multiplication gives a faithful unital representation
$\lambda:S\to\operatorname{End}_k(V)$.
The algebra $S$ is spanned by countably many words. Both $V$ and
$eV$ have countably infinite dimension, since $S$ and $eS$ are
nonzero. Choose a vector-space isomorphism $w:V\to eV$.
Represent $v$ by $w$ followed by the inclusion into $V$, and
represent $u$ by $w^{-1}\lambda(e)$. Then $uv=1_V$ and
$vu=\lambda(e)$. These operators give a representation of $B$
whose restriction to $S$ is faithful. Thus $B\ne0$ and $S\to B$
is injective. The choice of $w$ proves existence of a model;
it is not needed to compute the finite presentation.
\end{proof}

\begin{corollary}\label{cor:projective-undecidable}
There is no uniform algorithm deciding whether a specified finitely
generated projective over a finitely presented inverse semiring is
zero, even when the additive-idempotent support is $\B$ and the
projective is given by a one-by-one idempotent matrix. The corresponding
isomorphism problem is also undecidable.
\end{corollary}
\begin{proof}
Use Bokut's effective construction
\cite[Section~3, p.~139; Section~4, pp.~142--143]{Bokut} over
$k=\mathbb F_2$. His proof starts with a fixed finitely presented
semigroup $\Pi$ with undecidable word problem and the nonunital free
product $k[\Pi]*ke$, where $e^2=e$ and $ke$ is the chosen
obstruction algebra. Both factors embed by the free-product
description in his Section~2. A difference $a$ of two semigroup
words is zero exactly when the words agree; otherwise it is a
nonzero vector, hence primitive over $k$. His lemma effectively
constructs a finitely presented algebra $A(a)$ that is zero when
$a=0$ and contains the original free product when $a\ne0$.
The named idempotent $e$ therefore vanishes in the first case and
survives in the second. Lemma~\ref{lem:unital-idempotent-bridge}
converts this into a finite unital $\mathbb F_2$-algebra presentation
that is zero exactly when the words agree. This is also a finite
unital ring presentation, with the additional relation $2\cdot1=0$.
Proposition~\ref{prop:projective-decision-reduction} completes the
effective reduction to the claimed projective zero and isomorphism
tests.
\end{proof}
This obstruction concerns finite presentations, not supplied finite
operation tables, and makes no assertion about commutative rings.
It does not make mathematical classification by ordinary-ring data
impossible; it rules out the stated uniform decision procedures.


\section*{Note on method and provenance}
The project's independent zero-adjunction construction began in a conversation with Gemini~2.5 Pro in the Gemini app and was formalized by Gemini~2.5 Deep Think in Deep Think mode in November 2025. Versions were deposited on Zenodo under the name HI-AI from 6 November 2025 and later withdrawn. A GPT-5-series Pro model in ChatGPT Pro mode at reasoning effort ultra extracted part of this material in March 2026; GPT-5-series Pro models in that mode and effort developed the corpus from April to July 2026. Codex runs at effort ultra produced the August--September 2026 workbench. None of that work knew \cite{FGJM}, which preceded it. The November 2026 foundation was checked and written by Claude Opus~5.5 at effort max and checked and edited by Claude Fable~5.2 at effort max; related readings also used Claude Opus~5.5 at effort high. Those earlier internal records are project notes; the public foundation version is cited above.

This text and its proofs were produced by GPT-6 Astra Pro in ChatGPT Pro mode. The effort setting of this execution was not exposed. Separate workers of the same model re-derived the principal statements from their statements before comparing proofs; the dated records distinguish these checks from the earlier checks and from finite computations. No human has checked the mathematics. The accompanying \emph{Finite inverse-semiring checks} programs, documented through the workbench, verify the stated finite matrix and congruence cases. Our text is dedicated to the public domain under CC0~1.0. The workbench is \expandafter\url\expandafter{\workbenchurl}; the permanent record address is \expandafter\url\expandafter{\recordurl}.

\begin{thebibliography}{99}
\bibitem{BorgerJun} Borger, J., \& Jun, J. (2025). Facets of module theory over semirings. \emph{Advances in Mathematics, 479}, 110431. \url{https://doi.org/10.1016/j.aim.2025.110431} Preprint: \url{https://arxiv.org/abs/2405.18645}
\bibitem{BandeltPetrich} Bandelt, H.-J., \& Petrich, M. (1982). Subdirect products of rings and distributive lattices. \emph{Proceedings of the Edinburgh Mathematical Society, 25}, 155--171. \url{https://doi.org/10.1017/S0013091500016643}
\bibitem{BM} Bulman-Fleming, S., \& McDowell, K. (1980). Flatness and projectivity in commutative regular semigroups. \emph{Semigroup Forum, 19}, 247--259. \url{https://doi.org/10.1007/BF02572519}
\bibitem{CFMM} Caviglia, E., Faul, P. F., Manuell, G., \& Mesiti, L. (2026). \emph{Clifford semigroups and the monoidal Grothendieck construction} (Version 2). arXiv:2607.12944. \url{https://arxiv.org/abs/2607.12944}
\bibitem{Clifford} Clifford, A. H. (1941). Semigroups admitting relative inverses. \emph{Annals of Mathematics, 42}, 1037--1049.
\bibitem{CC} Connes, A., \& Consani, C. (2019). Homological algebra in characteristic one. \emph{Higher Structures, 3}, 155--247. \url{https://arxiv.org/abs/1703.02325}
\bibitem{DLN} Di Nola, A., Lenzi, G., \& Nam, T. G. (2020). Ultramatricial algebras over commutative chain semirings and application to MV-algebras. \emph{Forum Mathematicum, 32}(2), 287--305. \url{https://doi.org/10.1515/forum-2019-0056}
\bibitem{DJV} Dokuchaev, M., Johari, F., \& Vilca-Rodr\'iguez, J. L. (2026). Inverse semialgebras and partial actions of Lie algebras. \emph{Revista Matem\'atica Iberoamericana, 42}(5), 1637--1680. \url{https://doi.org/10.4171/RMI/1628}
\bibitem{FLS} Fiore, T. M., L\"uck, W., \& Sauer, R. (2011). Euler characteristics of categories and homotopy colimits. \emph{Documenta Mathematica, 16}, 301--354. \url{https://doi.org/10.4171/DM/333}
\bibitem{Emmanouil} Emmanouil, I. (2025). Residually nilpotent groups of homological dimension~1. \emph{Bulletin of the London Mathematical Society, 57}(10), 3223--3232. \url{https://doi.org/10.1112/blms.70140}
\bibitem{FGJM} Faul, P. F., Goswami, A., Joubert, G., \& Manuell, G. (2024). \emph{The case for inverse semirings}. arXiv:2411.12318. \url{https://arxiv.org/abs/2411.12318}
\bibitem{Golan} Golan, J. S. (1999). \emph{Semirings and their applications}. Kluwer.
\bibitem{HK} Horn, A., \& Kimura, N. (1971). The category of semilattices. \emph{Algebra Universalis, 1}, 26--38.
\bibitem{JMT} Jun, J., Mincheva, K., \& Tolliver, J. (2026). Representation theory over semifields. \emph{Journal of Algebra, 706}, 243--265. \url{https://doi.org/10.1016/j.jalgebra.2026.05.022} Preprint: \url{https://arxiv.org/abs/2411.19883}
\bibitem{KNZ} Katsov, Y., Nam, T. G., \& Zumbr\"agel, J. (2018). On congruence-semisimple semirings and the $K_0$-group characterization of ultramatricial algebras over semifields. \emph{Journal of Algebra, 508}, 157--195. \url{https://doi.org/10.1016/j.jalgebra.2018.04.024} Preprint: \url{https://arxiv.org/abs/1711.05163}
\bibitem{Lorscheid} Lorscheid, O. (2012). The geometry of blueprints. Part I: Algebraic background and scheme theory. \emph{Advances in Mathematics, 229}, 1804--1846. \url{https://arxiv.org/abs/1103.1745}
\bibitem{Lorscheid2026} Lorscheid, O. (2026). \emph{Two comparison theorems for semiring schemes}. arXiv:2603.28616. \url{https://arxiv.org/abs/2603.28616}
\bibitem{Macpherson} Macpherson, A. W. (2018). \emph{Projective modules over polyhedral semirings} (July 14 preprint). arXiv:1507.07213. \url{https://arxiv.org/abs/1507.07213}
\bibitem{Mitchell} Mitchell, B. (1972). Rings with several objects. \emph{Advances in Mathematics, 8}(1), 1--161. \url{https://doi.org/10.1016/0001-8708(72)90002-3}
\bibitem{Murfet} Murfet, D. (2006). \emph{Rings with several objects} (October 5 notes). \url{https://therisingsea.org/notes/RingsWithSeveralObjects.pdf}
\bibitem{NicholsonExchange} Nicholson, W. K. (1977). Lifting idempotents and exchange rings. \emph{Transactions of the American Mathematical Society, 229}, 269--278. \url{https://doi.org/10.1090/S0002-9947-1977-0439876-2}
\bibitem{NielsenExchange} Nielsen, P. P. (2006). \emph{The exchange property for modules and rings} [Doctoral dissertation, University of California, Berkeley]. \url{https://mathdept.byu.edu/~pace/PaceThesis_web.pdf}
\bibitem{DasKarExchange} Das, D., \& Kar, S. (2022). \emph{Strongly clean semiring} (Version 1) [Preprint]. Research Square. \url{https://doi.org/10.21203/rs.3.rs-1751073/v1}
\bibitem{Pasku} Pasku, E. (2011). Clifford semigroups as functors and their cohomology. \emph{Semigroup Forum, 83}, 75--88.
\bibitem{SBM} Sen, M. K., Bhuniya, A. K., \& Maity, S. K. (2021). Inverse semimodules. \emph{Asian-European Journal of Mathematics, 14}(2), 2150022. \url{https://doi.org/10.1142/S1793557121500224}
\bibitem{Stacks} The Stacks Project Authors. (2026). \emph{The Stacks Project}. Lemmas~10.32.7 (Tag~05BU), 10.77.6 (Tag~0D47), and 15.6.9 (Tag~0D2G). \url{https://stacks.math.columbia.edu}
\bibitem{OlsonJenkins} Olson, D. M., \& Jenkins, T. L. (1973). A note on the Levitzki radical of a semiring. \emph{Fundamenta Mathematicae, 79}(2), 91--93. \url{https://doi.org/10.4064/fm-79-2-91-93}
\bibitem{Rich} Rich, M. (1982). Some properties of the Levitzki radical in alternative rings. \emph{Journal of the Australian Mathematical Society (Series A), 32}(1), 52--60. \url{https://doi.org/10.1017/S1446788700024393}
\bibitem{Rowen} Rowen, L. H. (2016). \emph{Algebras with a negation map}. arXiv:1602.00353. \url{https://arxiv.org/abs/1602.00353}
\bibitem{HerberaPrihoda} Herbera, D., \& P\v r\'ihoda, P. (2014). Reconstructing projective modules from its trace ideal. \emph{Journal of Algebra, 416}, 25--57. \url{https://doi.org/10.1016/j.jalgebra.2014.06.010}. Preprint: \url{https://arxiv.org/abs/1309.5021}
\bibitem{CortesTorrecillas} Cort\'es-Izurdiaga, M., \& Torrecillas, B. (2013). Projective and flat objects over rooted rings with several objects. \emph{Bulletin Math\'ematique de la Soci\'et\'e des Sciences Math\'ematiques de Roumanie, 56(104)}(3), 315--330.
\bibitem{FoundationPublic} \emph{Spectrum, semimodules and projectives of the inverse semiring $R_0$ of Faul, Goswami, Joubert and Manuell}. (2026). Hexagon 2611.00993v1. \url{https://hexagonmath.org/2611.00993v1}
\bibitem{KatsovNamRadicals} Katsov, Y., \& Nam, T. G. (2014). On radicals of semirings and related problems. \emph{Communications in Algebra, 42}(12), 5065--5099. Preprint: \url{https://arxiv.org/abs/1206.1580}
\bibitem{Krause} Krause, H. (2015). Krull--Schmidt categories and projective covers. \emph{Expositiones Mathematicae, 33}, 535--549. \url{https://doi.org/10.1016/j.exmath.2015.10.001}
\bibitem{CGQ} Cohen, G., Gaubert, S., \& Quadrat, J.-P. (2004). Duality and separation theorems in idempotent semimodules. \emph{Linear Algebra and its Applications, 379}, 395--422. \url{https://doi.org/10.1016/j.laa.2003.08.010} Preprint: \url{https://arxiv.org/abs/math/0212294}
\bibitem{DeorePatil} Deore, R. P., \& Patil, K. B. (2005). On the dual basis of projective semimodules and its applications. \emph{Sarajevo Journal of Mathematics, 1(14)}, 161--169. \url{https://doi.org/10.5644/SJM.01.2.02}
\bibitem{NW} Niefield, S. B., \& Wood, R. J. (2017). Coexponentiability and projectivity: Rigs, rings, and quantales. \emph{Theory and Applications of Categories, 32}(36), 1222--1228. \url{https://tac.mta.ca/tac/volumes/32/36/32-36.pdf}
\bibitem{Niefield} Niefield, S. (2016). Projectivity, continuity, and adjointness: Quantales, $Q$-posets, and $Q$-modules. \emph{Theory and Applications of Categories, 31}(30), 839--851. \url{https://tac.mta.ca/tac/volumes/31/30/31-30.pdf}
\bibitem{Bokut} Bokut', L. A. (1970). Unsolvability of certain algorithmic problems in a class of associative rings. \emph{Algebra i Logika, 9}(2), 137--144. \url{https://www.mathnet.ru/eng/al1237}
\end{thebibliography}
\end{document}
